% Options for packages loaded elsewhere
\PassOptionsToPackage{unicode}{hyperref}
\PassOptionsToPackage{hyphens}{url}
\documentclass[
]{article}
\usepackage{xcolor}
\usepackage{amsmath,amssymb}
\setcounter{secnumdepth}{-\maxdimen} % remove section numbering
\usepackage{iftex}
\ifPDFTeX
  \usepackage[T1]{fontenc}
  \usepackage[utf8]{inputenc}
  \usepackage{textcomp} % provide euro and other symbols
\else % if luatex or xetex
  \usepackage{unicode-math} % this also loads fontspec
  \defaultfontfeatures{Scale=MatchLowercase}
  \defaultfontfeatures[\rmfamily]{Ligatures=TeX,Scale=1}
\fi
\usepackage{lmodern}
\ifPDFTeX\else
  % xetex/luatex font selection
\fi
% Use upquote if available, for straight quotes in verbatim environments
\IfFileExists{upquote.sty}{\usepackage{upquote}}{}
\IfFileExists{microtype.sty}{% use microtype if available
  \usepackage[]{microtype}
  \UseMicrotypeSet[protrusion]{basicmath} % disable protrusion for tt fonts
}{}
\makeatletter
\@ifundefined{KOMAClassName}{% if non-KOMA class
  \IfFileExists{parskip.sty}{%
    \usepackage{parskip}
  }{% else
    \setlength{\parindent}{0pt}
    \setlength{\parskip}{6pt plus 2pt minus 1pt}}
}{% if KOMA class
  \KOMAoptions{parskip=half}}
\makeatother
\usepackage{longtable,booktabs,array}
\newcounter{none} % for unnumbered tables
\usepackage{calc} % for calculating minipage widths
% Correct order of tables after \paragraph or \subparagraph
\usepackage{etoolbox}
\makeatletter
\patchcmd\longtable{\par}{\if@noskipsec\mbox{}\fi\par}{}{}
\makeatother
% Allow footnotes in longtable head/foot
\IfFileExists{footnotehyper.sty}{\usepackage{footnotehyper}}{\usepackage{footnote}}
\makesavenoteenv{longtable}
\usepackage{graphicx}
\makeatletter
\newsavebox\pandoc@box
\newcommand*\pandocbounded[1]{% scales image to fit in text height/width
  \sbox\pandoc@box{#1}%
  \Gscale@div\@tempa{\textheight}{\dimexpr\ht\pandoc@box+\dp\pandoc@box\relax}%
  \Gscale@div\@tempb{\linewidth}{\wd\pandoc@box}%
  \ifdim\@tempb\p@<\@tempa\p@\let\@tempa\@tempb\fi% select the smaller of both
  \ifdim\@tempa\p@<\p@\scalebox{\@tempa}{\usebox\pandoc@box}%
  \else\usebox{\pandoc@box}%
  \fi%
}
% Set default figure placement to htbp
\def\fps@figure{htbp}
\makeatother
\setlength{\emergencystretch}{3em} % prevent overfull lines
\providecommand{\tightlist}{%
  \setlength{\itemsep}{0pt}\setlength{\parskip}{0pt}}
\usepackage{bookmark}
\IfFileExists{xurl.sty}{\usepackage{xurl}}{} % add URL line breaks if available
\urlstyle{same}
\hypersetup{
  hidelinks,
  pdfcreator={LaTeX via pandoc}}

\author{}
\date{}

\begin{document}

\begin{enumerate}
\def\labelenumi{\arabic{enumi}.}
\tightlist
\item
  \textbf{Algunos convenios}:
\end{enumerate}

\begin{enumerate}
\def\labelenumi{\arabic{enumi}.}
\item
  \begin{enumerate}
  \def\labelenumii{\arabic{enumii}.}
  \item
    \[0 \notin {\mathbb{N}}\]Ante la definición de los naturales
    nosotros adoptamos este convenio.
  \item
    \[{\mathbb{N}}_{0}{\qquad ≝ \qquad}{{\mathbb{N}} \cup \left\{ 0 \right\}}\]Conjunto
    de los naturales extendidos que ya tiene el cero.
  \item
    Definición recurrente de los
    conjuntos\[\lbrack 0,1\rbrack_{\mathbb{Q}}\]y\[\lbrack 0,1\rbrack_{\mathbb{Q}}^{n}\],
    dónde \[{n \in {\mathbb{N}}}\qquad{n > 1}\]:

    \begin{enumerate}
    \def\labelenumiii{\arabic{enumiii}.}
    \tightlist
    \item
      \begin{enumerate}
      \def\labelenumiv{\arabic{enumiv}.}
      \tightlist
      \item
        \[\lbrack 0,1\rbrack_{\mathbb{Q}}{\qquad ≝ \qquad}{\lbrack 0,1\rbrack \cap {\mathbb{Q}}}\]
      \item
        \[\phantom{{\forall{n \in {{\mathbb{N}} \smallsetminus \left\{ 0,1 \right\}}}}\mspace{144mu}}\lbrack 0,1\rbrack_{\mathbb{Q}}^{1}{\qquad ≝ \qquad}\lbrack 0,1\rbrack_{\mathbb{Q}}\]
      \item
        \[\phantom{{\forall{n \in {{\mathbb{N}} \smallsetminus \left\{ 0,1 \right\}}}}\mspace{144mu}}{{{\lbrack 0,1\rbrack_{\mathbb{Q}}}^{n}{: = {{\lbrack 0,1\rbrack_{\mathbb{Q}}}^{n - 1} \times {\lbrack 0,1\rbrack_{\mathbb{Q}}}^{1}}}}{\mspace{72mu}\forall{n \in {{\mathbb{N}} \smallsetminus {\{ 0,1\}}}}}}\]
      \end{enumerate}
    \end{enumerate}
  \item
    Por lo general los elementos de un conjunto se representarán por
    letras minúsculas (alfa\-betos griego y latino) con o sin
    suscriptores
    (ejemplo:\[a_{3}\]o\includegraphics[width=0.425cm,height=0.467cm]{./ObjectReplacements/Object 1513}o\includegraphics[width=0.982cm,height=0.374cm]{./ObjectReplacements/Object 1514}o\includegraphics[width=0.422cm,height=0.344cm]{./ObjectReplacements/Object 1515})
    y dí\-gitos decimales (\[\{{0,1,\ldots,9}\}\]), mientras que los
    conjuntos se representarán por letras mayúsculas (alfabetos griego y
    latino), igualmente con o sin suscriptores. Este convenio será
    válido a excepción que se exprese de forma explícita otro nombre
    para elementos y/o conjuntos.
  \item
    En ocasiones aseguraremos que existe un conjunto asociado a un
    elemento: en general serán letras mayúsculas como corresponde a un
    conjunto, pero con un subscriptor que estará escrito exactamente
    como el elemento. La única excepción que se dará es un elemento
    cubierto con la tilde circunfleja, para expresar el conjunto de
    elementos complementa\-rios con uno dado.
  \item
    A veces aparecerá una operación binaria,
    digamos\includegraphics[width=0.859cm,height=0.467cm]{./ObjectReplacements/Object 1573},
    dónde los puntos se sustitu\-yen con los argumentos. Cuando
    expresemos\[a \ast B\], o sea, el elemento \[a\]operado con un
    conjunto\[B\], se trata una operación binaria que no es la original.
    Se interpreta\-rá
    como\[{a \ast B} = {\{{{a \ast b}{\qquad \mid \qquad}{b \in B}}\}}\],
    esto es, un conjunto. Igualmente, puede ocurrir que esta misma
    operación aparezca entre dos
    conjuntos:\[{A \ast B} = {\{{{a \ast b}{\qquad \mid \qquad}{{a \in A},{b \in B}}}\}}\].
  \item
    El universo en el que nos moveremos
    será\includegraphics[width=0.51cm,height=0.563cm]{./ObjectReplacements/Object 1579},
    por lo que para el cuantificador univer\-sal (que evitaremos en lo
    posible) no pondremos ningún símbolo, de forma que, cuando aparezca
    una variable de elementos o de conjuntos sin haber sido
    cuantificada, supondremos que se aplica el cuantificador universal a
    la variable situándola dentro
    de\includegraphics[width=0.51cm,height=0.563cm]{./ObjectReplacements/Object 1580}.
    Si el cuantificador universal debiera aplicarse a una variable sobre
    otro conjunto subconjunto
    de\includegraphics[width=0.51cm,height=0.563cm]{./ObjectReplacements/Object 1581},
    se intentará omitir el cuantificador universal, pero se precederá la
    sentencia con una indicación sobre la pertenencia del elemento al
    conjunto. Para cualquier elemento en el que no aparezca su
    pertenencia es porque pertenece al
    universal\includegraphics[width=0.51cm,height=0.563cm]{./ObjectReplacements/Object 1582}.
    Al igual, cuando una variable que asume un valor que es un conjunto,
    éste será parte
    de\includegraphics[width=1.829cm,height=0.624cm]{./ObjectReplacements/Object 1583},
    que no aparecerá de forma explícita. Para las variables conjunto se
    repite todo lo anterior con la diferencia de que se
    cambia\includegraphics[width=0.51cm,height=0.563cm]{./ObjectReplacements/Object 1585}por
    \includegraphics[width=1.829cm,height=0.624cm]{./ObjectReplacements/Object 1321}.
  \end{enumerate}
\item
  Álgebra de Boole.

  \begin{enumerate}
  \def\labelenumii{\arabic{enumii}.}
  \item
    Axiomas para un álgebra de
    Boole\includegraphics[width=7.008cm,height=0.667cm]{./ObjectReplacements/Object 1}.

    La operación uno-aria que se utilizará frecuentemente (la
    complementación), en los axiomas será introducida como definición y
    no supuesta inicialmente más que como la existencia de un conjunto
    de elementos no vacío asociados a otro que cumplen unos
    requisi\-tos. Sin embargo esta operación es tan importante en las
    álgebras de Boole que la forma normal de escribirse la estructura es
    \[\left\langle {,{\left\{ {0_{},1_{}} \right\} \subseteq},{0_{} \neq 1_{}},{\mathit{op}\left\{ {\mathit{bin}:\left\{ {+ , \cdot} \right\}} \right\},\left\{ {\mathit{una}:\left\{ \overline{\phantom{A}} \right\}} \right\}}} \right\rangle\].

    Además, la distinción entre los dos neutros no siempre se requiere.
    Esto tiene poco impacto, ya que el único álgebra de Boole dónde los
    dos neutros son iguales coincide con el álgebra de Boole con un solo
    elemento, que es el álgebra de Boole trivial.

    Los postulados de Huntington (artículos en 1904,1932 -este último
    desarrolla un con\-junto de 3 axiomas, uno de ellos llamado
    específicamente Axioma de Huntington, y no es el caso aquí expuesto,
    en esta primera formalización-) definen qué sea un álgebra de Boole
    \includegraphics[width=3.173cm,height=0.572cm]{./ObjectReplacements/Object 1094}:

    \begin{enumerate}
    \def\labelenumiii{\arabic{enumiii}.}
    \item
      \textbf{{[}Axioma H0{]}} De
      ser\includegraphics[width=0.51cm,height=0.563cm]{./ObjectReplacements/Object 1323}un
      conjunto, las operaciones binarias funciones y de la existencia
      de\includegraphics[width=0.422cm,height=0.467cm]{./ObjectReplacements/Object 613}y\includegraphics[width=0.411cm,height=0.467cm]{./ObjectReplacements/Object 1072}en\[\]:

      \begin{enumerate}
      \def\labelenumiv{\arabic{enumiv}.}
      \tightlist
      \item
        \textbf{{[}H0.0.0{]}}
        \includegraphics[width=0.51cm,height=0.563cm]{./ObjectReplacements/Object 1316}es
        un conjunto (por ejemplo en la axiomatización de
        Zermelo-Fraenkel-Skolem o más brevemente, \textbf{ZFS}, o
        digamos la más extensa Neumann-Gödel-Bernays, abreviadamente,
        \textbf{NGB}).
      \item
        \textbf{{[}H0.0.1{]}} Para la suma
        \includegraphics[width=0.492cm,height=0.467cm]{./ObjectReplacements/Object 1073}:\[0 \in\]
      \item
        \textbf{{[}H0.0.2{]}} Para el producto
        \includegraphics[width=0.445cm,height=0.467cm]{./ObjectReplacements/Object 1092}:\includegraphics[width=1.053cm,height=0.563cm]{./ObjectReplacements/Object 1074}
      \item
        \textbf{{[}H0.1{]}} Para la suma
        \includegraphics[width=0.492cm,height=0.467cm]{./ObjectReplacements/Object 499}:\includegraphics[width=5.819cm,height=0.577cm]{./ObjectReplacements/Object 2}
      \item
        \textbf{{[}H0.2{]}} Para la multiplicación
        \includegraphics[width=0.445cm,height=0.467cm]{./ObjectReplacements/Object 18}:\includegraphics[width=5.666cm,height=0.577cm]{./ObjectReplacements/Object 3}
      \end{enumerate}
    \item
      \textbf{{[}Axioma H1{]}} Existencia del elemento identidad o
      neutro:

      \begin{enumerate}
      \def\labelenumiv{\arabic{enumiv}.}
      \tightlist
      \item
        \textbf{{[}H1.1{]}} Para la suma
        \includegraphics[width=0.492cm,height=0.467cm]{./ObjectReplacements/Object 210}:
        \includegraphics[width=3.006cm,height=0.563cm]{./ObjectReplacements/Object 500}
      \item
        \textbf{{[}H1.2{]}} Para la multiplicación
        \includegraphics[width=0.445cm,height=0.467cm]{./ObjectReplacements/Object 197}:
        \includegraphics[width=2.842cm,height=0.563cm]{./ObjectReplacements/Object 198}
      \end{enumerate}
    \item
      \textbf{{[}Axioma H2{]}} Conmutabilidad:

      \begin{enumerate}
      \def\labelenumiv{\arabic{enumiv}.}
      \tightlist
      \item
        \textbf{{[}H2.1{]}} Para la suma
        \includegraphics[width=0.492cm,height=0.467cm]{./ObjectReplacements/Object 501}:
        \includegraphics[width=5.022cm,height=0.577cm]{./ObjectReplacements/Object 4}
      \item
        \textbf{{[}H2.2{]}} Para la multiplicación
        \includegraphics[width=0.445cm,height=0.467cm]{./ObjectReplacements/Object 19}:
        \includegraphics[width=4.715cm,height=0.577cm]{./ObjectReplacements/Object 5}
      \end{enumerate}
    \item
      \textbf{{[}Axioma H3{]}} Propiedad distributiva:

      \begin{enumerate}
      \def\labelenumiv{\arabic{enumiv}.}
      \item
        \textbf{{[}H3.1{]}} De la
        suma\includegraphics[width=0.492cm,height=0.467cm]{./ObjectReplacements/Object 502}sobre
        el
        producto\includegraphics[width=0.445cm,height=0.467cm]{./ObjectReplacements/Object 523}(o,
        inversamente, \textbf{sacar sumando común}):

        \includegraphics[width=8.181cm,height=0.577cm]{./ObjectReplacements/Object 6}
      \item
        \textbf{{[}H3.2{]}} Del
        producto\includegraphics[width=0.445cm,height=0.467cm]{./ObjectReplacements/Object 524}sobre
        la suma
        \includegraphics[width=0.492cm,height=0.467cm]{./ObjectReplacements/Object 503}(o,
        inversamente, \textbf{sacar factor común}):

        \includegraphics[width=8.027cm,height=0.577cm]{./ObjectReplacements/Object 7}
      \end{enumerate}
    \item
      \textbf{{[}Axioma H4{]}} Existencia de complementario:

      \begin{enumerate}
      \def\labelenumiv{\arabic{enumiv}.}
      \item
        \textbf{{[}H4{]}}\includegraphics[width=4.253cm,height=0.563cm]{./ObjectReplacements/Object 8}
        tal que
        \includegraphics[width=1.464cm,height=0.467cm]{./ObjectReplacements/Object 1324}se
        cumplen las dos sub-fórmulas siguientes:

        \begin{enumerate}
        \def\labelenumv{\arabic{enumv}.}
        \tightlist
        \item
          \textbf{{[}H4.1{]}}\[{x + y} = 1\]
        \item
          \textbf{{[}H4.2{]}}\[x{\mspace{9mu} \cdot \mspace{9mu}}{y = 0}\]
        \end{enumerate}
      \end{enumerate}
    \end{enumerate}
  \end{enumerate}

  Como ejemplos que cumplen los anteriores postulados o axiomas (a los
  que llamaremos propiamente modelos de la teoría de Álgebras de Boole
  axiomatizada según los anteriores axiomas) vamos a desarrollar unos
  cuántos de ellos.

  \textbf{{[}Ejemplo 1{]}} \textbf{El álgebra de las partes de un
  conjunto}\includegraphics[width=0.527cm,height=0.469cm]{./ObjectReplacements/Object 1317}.

  El ejemplo primero y más sencillo de ver es el álgebra de las partes
  de un conjunto. Dado que para un conjunto
  cualquiera\includegraphics[width=0.529cm,height=0.467cm]{./ObjectReplacements/Object 83},\includegraphics[width=3.371cm,height=0.58cm]{./ObjectReplacements/Object 1046}es
  también un conjunto en cualquiera de los dos sistemas axiomáticos de
  teoría de conjuntos mencionados.

  Además\includegraphics[width=6.479cm,height=0.58cm]{./ObjectReplacements/Object 1051},
  dónde se verifica
  que\includegraphics[width=4.544cm,height=0.64cm]{./ObjectReplacements/Object 1068}.

  Definiremos\includegraphics[width=1.907cm,height=0.58cm]{./ObjectReplacements/Object 1050},
  \includegraphics[width=1.145cm,height=0.467cm]{./ObjectReplacements/Object 1325}y
  \includegraphics[width=1.094cm,height=0.467cm]{./ObjectReplacements/Object 1326}.
  Como producto lógico pondremos la intersección de
  conjuntos\includegraphics[width=4.644cm,height=0.58cm]{./ObjectReplacements/Object 1052}y
  como suma lógica pondremos la unión de
  conjuntos\includegraphics[width=4.798cm,height=0.58cm]{./ObjectReplacements/Object 1053}.

  Las conmutabilidades de la suma y el producto se intercambian
  directamente por la conmutabilidad de la unión y la intersección.

  El elemento neutro de la suma es trivialmente
  el\includegraphics[width=0.422cm,height=0.467cm]{./ObjectReplacements/Object 1327}ya
  que
  \includegraphics[width=0.57cm,height=0.385cm]{./ObjectReplacements/Object 1328}es
  neutro para la unión. Igualmente el neutro del producto es
  el\includegraphics[width=0.411cm,height=0.467cm]{./ObjectReplacements/Object 1329}por
  ser
  \includegraphics[width=0.549cm,height=0.467cm]{./ObjectReplacements/Object 1330}el
  neutro de la intersección en
  \includegraphics[width=1.261cm,height=0.58cm]{./ObjectReplacements/Object 1331}.

  Las distributividades recaen directamente en las de la unión sobre la
  intersección y de la intersección sobre la unión.

  Sea
  \includegraphics[width=6.387cm,height=0.612cm]{./ObjectReplacements/Object 1054}y
  a partir de ahí sabemos
  que\includegraphics[width=2.069cm,height=0.612cm]{./ObjectReplacements/Object 82}puesto
  que\includegraphics[width=4.688cm,height=0.531cm]{./ObjectReplacements/Object 1055}y
  en el caso
  que\includegraphics[width=1.21cm,height=0.467cm]{./ObjectReplacements/Object 1056}tenemos
  que
  \includegraphics[width=3.983cm,height=0.531cm]{./ObjectReplacements/Object 1057}de
  forma
  que\includegraphics[width=2.736cm,height=0.612cm]{./ObjectReplacements/Object 1058}por
  la definición. Ahora sólo se trata de darse cuenta
  que\[{Y_{X} \cdot X} = {Y_{X} \cap X} = {{({U \smallsetminus X})} \cap X} = \varnothing \equiv 0\]y
  que\includegraphics[width=6.334cm,height=0.54cm]{./ObjectReplacements/Object 1065},
  y ya tenemos
  que\includegraphics[width=5.978cm,height=0.612cm]{./ObjectReplacements/Object 1066}habiendo
  tomado\includegraphics[width=1.976cm,height=0.531cm]{./ObjectReplacements/Object 1067}.

  \textbf{{[}Ejemplo 2{]}} \textbf{El álgebra de los números que son
  producto de
  los\includegraphics[width=0.422cm,height=0.467cm]{./ObjectReplacements/Object 1318}primeros
  números primos.}

  Un ejemplo interesante fácil de construir es el álgebra de Boole de
  los números que son producto de los primeros núme\-ros primos
  (cantidad finita de ellos) y sus divisores. Consideramos el
  conjunto\includegraphics[width=3.097cm,height=0.531cm]{./ObjectReplacements/Object 184},
  con\-sideraremos
  el\includegraphics[width=0.411cm,height=0.467cm]{./ObjectReplacements/Object 1626}booleano
  cómo
  \includegraphics[width=1.852cm,height=0.896cm]{./ObjectReplacements/Object 185}y
  el\includegraphics[width=0.423cm,height=0.467cm]{./ObjectReplacements/Object 1627}cómo\includegraphics[width=1.371cm,height=0.587cm]{./ObjectReplacements/Object 200}.
  Consideramos
  a\includegraphics[width=3.822cm,height=0.656cm]{./ObjectReplacements/Object 1069},
  y las operaciones serán el mínimo común múltiplo
  (\includegraphics[width=3.621cm,height=0.506cm]{./ObjectReplacements/Object 1319})
  como suma booleana y el máximo común divisor a
  (\includegraphics[width=3.353cm,height=0.506cm]{./ObjectReplacements/Object 1320})
  como producto booleano. El complemento de un elemento resulta
  ser\includegraphics[width=6.731cm,height=0.915cm]{./ObjectReplacements/Object 1070}.
  Éste es un modelo fácil de desarrollar para poner ejemplos.

  \textbf{{[}Ejemplo 3{]}} Partimos de un álgebra de Boole
  \includegraphics[width=0.51cm,height=0.563cm]{./ObjectReplacements/Object 1312}cualquiera
  y un elemento \[x\]
  no\includegraphics[width=0.609cm,height=0.587cm]{./ObjectReplacements/Object 1076}ni\[1_{}\]cualquiera
  tal
  que\includegraphics[width=2.641cm,height=0.669cm]{./ObjectReplacements/Object 1077}(luego
  este álgebra de Boole no puede ser el trivial ni uno de dos
  elementos). Definimos ahora un
  conjunto\includegraphics[width=3.997cm,height=0.61cm]{./ObjectReplacements/Object 1078}y\includegraphics[width=4.126cm,height=0.61cm]{./ObjectReplacements/Object 1079}.
  Las álgebras de Boole nuevas a considerar son
  \includegraphics[width=10.156cm,height=0.695cm]{./ObjectReplacements/Object 1082}y\includegraphics[width=10.146cm,height=0.695cm]{./ObjectReplacements/Object 1083}.
  Consideramos el mismo producto booleano que en el conjunto inicial e
  idéntica\-mente con la suma booleana (solo que restringidos al
  subconjunto elegido). Sólo varía el complemento, de la siguiente
  forma\includegraphics[width=3.872cm,height=0.612cm]{./ObjectReplacements/Object 1080}y\includegraphics[width=3.715cm,height=0.61cm]{./ObjectReplacements/Object 1081}.
  Es fácil comprobar la validez de esta definición de un álgebra de
  Boole, de manera más concreta, que las operaciones son internas y el
  complemento declarado es también interno y se comporta como
  complemento del nuevo álgebra, esto
  es,\includegraphics[width=5.189cm,height=0.61cm]{./ObjectReplacements/Object 1084}
  que\includegraphics[width=5.2cm,height=0.61cm]{./ObjectReplacements/Object 1085}.

  \textbf{{[}Ejemplo 4{]}} El álgebra de las proposiciones. Éste es sin
  duda, el primer desarrollo que se hizo del álgebra de Boole, hecha por
  el propio George Boole en mitad del siglo XIX (edición 1851). Lo
  desarrolló como una ``Una investigación en las leyes del
  pensamiento''. Amigo suyo que lo ayudó a penetrar los ambientes
  académi\-cos es Augustus De Morgan. Desde Aristóteles (que funda por
  primera vez la ló\-gica como ciencia analítica) en el siglo IV a. C.
  no había habido ningún adelanto sustancial en la lógica. Kant (medio
  siglo antes de Boole) había considerado que la lógica era un cuerpo de
  doctrina cerrado y completo (esto es, no había nada más que decir que
  lo que ya había desarrollado y escrito Aristóteles en sus ``Tra\-tados
  de Lógica'' u ``Órganon'', hacía ya 2.100 años). Aunque la verdadera
  revolu\-ción se da algunos años más tarde con Frege, el lógico más
  importante desde Aristóteles. Si doy estos datos sobre la historia de
  la lógica que todos asociaréis más a la filosofía, que parece queda
  muy lejos del propósito de unos apuntes de matemáticas discretas que
  cubran de la forma más amplia posible los Fundamentos de Electrónica
  en su parte Digital, es porque en realidad no queda tan lejos. La idea
  de hacer un lenguaje dónde el razonamiento siguiera unas pautas claras
  de forma que siempre quedara todo tan cierto como en las matemáticas
  era ya antigua. Aristóteles ya advertía de una cierta indefinición
  insuperable de los términos más importantes de la filosofía (en
  realidad de casi todos los conceptos de la vida ordinaria): ``existen
  conceptos o ideas que corresponden con la realidad que no se usan de
  forma equívoca -- esto es, su uso no es equívoco, este mismo concepto,
  palabra o idea que hablamos no se refiere a realidades distintas y
  diferenciadas, de forma que nos llevan a confusión -- pero tampoco de
  forma unívoca -- como las definiciones desde axiomas en un lenguaje
  formal matemático, así tenemos que existen conceptos análogos'' (es
  una glosa de palabras de Aristóteles). En la Modernidad, dado que el
  concepto de analogía lleva aparejado un tratamiento difícil que no
  lleva fácilmente a certeza, se intentan buscar criterios de certeza
  absoluta y unas definiciones que aparentemente son unívocas y se
  tratan como tales. Es significativo el nombre (y la estructura
  interna) de una importante obra de Spino\-za: ``Ética demostrada según
  el orden geométrico''. Será Leibniz quién escriba ya cumplidamente
  sobre la necesidad de establecer un léxico completamente unívoco (una
  tarea mastodonte, o mejor, imposible) y un ``cálculo'' del
  pensamiento, de forma que ``una cuestión como la existencia de Dios
  pueda ser resuelta mediante la resolución de unas ecuaciones de
  pensamiento'' (de nuevo es una glosa). Se empezaba a buscar con
  ansiedad una mecanización del pensamiento. Este motivo fue muy
  fructífero, dando un primer paso hacia él, George Boole que hace un
  ál\-gebra de las proposiciones. Esta álgebra no era más amplia que la
  de Aristóteles (de hecho, era solo un subconjunto de la lógica
  aristotélica), pero permitía el cálculo al modo algebraico. De aquí a
  la llegada de Frege, ya, Charles Babbage diseña y realiza (sin éxito
  debido al trabajo de mecanizado ex\-cesivamente minucioso que requería
  el diseño) una computadora universal me\-cánica (mediante engranajes)
  prácticamente similar al modelo de Von Neumann. La condesa de Lovelace
  (Ada) es la matemática que hace los primeros progra\-mas en lenguaje
  ensamblador de la máquina de Babbage. La primera programa\-dora de la
  historia. Después viene Frege, al que debemos el invento de la
  Lógica-Matemática, o lógica extensional en aquel nacer. Ampliaba
  inmensamente la lógica aristotélica pues la estructura de la
  proposición no seguía ya el esquema sujeto-predicado que se mostraba
  insuficiente por lo que adopta el esquema función-argumentos, además
  de introducir los cuantificadores (``existe al menos un''
  \includegraphics[width=0.476cm,height=0.333cm]{./ObjectReplacements/Object 1266}y
  ``para todo''
  \includegraphics[width=0.591cm,height=0.358cm]{./ObjectReplacements/Object 1293}).
  Después de Frege siguen los desarrollos con gente como Ber\-trand
  Rusell, David Hilbert y otros, hasta los increíbles resultados de
  Gödel (y posteriormente Turing) que ponen punto final a muchas de las
  pretensiones de mecanización del pensamiento, pero que son ya la base
  de la computación moderna, siendo los trabajos definitivos los de Alan
  Turing. Como veis el camino recorrido es largo y complicado, siendo el
  momento crucial para el arranque de la ingeniería digital los trabajos
  de George Boole. No he mencionado intentos de mecanización más
  antiguos como los del mallorquín Raimundo Lulio, ni la calculadora
  binaria de Leibniz, ni mecanismos remotos como el mecanismo de
  Anticitera, así como el papel fundamental que jugó la crisis de las
  matemáticas de finales del siglo XIX y principios del XX ni el papel
  de las máquinas de cifrado y descifrado de mensajes en la Gran Guerra
  y la II Guerra Mundial (en las que intervinieron muchas de las mentes
  antes mencionadas).

  De manera un tanto informal podemos ver una proposición (una frase que
  afirma o niega una propiedad de un objeto, una relación entre objetos
  o la existencia del mismo, una frase que ha de ser o
  verdadero,\includegraphics[width=3.002cm,height=0.587cm]{./ObjectReplacements/Object 1086},
  o
  falso,\includegraphics[width=3.129cm,height=0.587cm]{./ObjectReplacements/Object 1087}).
  Este conjunto de proposiciones pueden ser operadas mediante la
  conjunción
  `y'(\includegraphics[width=0.554cm,height=0.467cm]{./ObjectReplacements/Object 1295}),
  \includegraphics[width=1.164cm,height=0.467cm]{./ObjectReplacements/Object 1297}es
  verdadero
  si\includegraphics[width=0.526cm,height=0.467cm]{./ObjectReplacements/Object 1298}es
  verdadero
  y\includegraphics[width=0.487cm,height=0.467cm]{./ObjectReplacements/Object 1299}es
  verdadero a la vez y falso en cualquier otro caso. La disyunción sería
  la
  `o'(\includegraphics[width=0.554cm,height=0.467cm]{./ObjectReplacements/Object 1300}),
  siendo\includegraphics[width=1.164cm,height=0.467cm]{./ObjectReplacements/Object 1301}verdadero
  con
  que\includegraphics[width=0.526cm,height=0.467cm]{./ObjectReplacements/Object 1302}sea
  verdadero o lo
  sea\includegraphics[width=0.487cm,height=0.467cm]{./ObjectReplacements/Object 1303},
  siendo falso sólo
  cuando\includegraphics[width=0.526cm,height=0.467cm]{./ObjectReplacements/Object 1304}es
  falso
  y\includegraphics[width=0.487cm,height=0.467cm]{./ObjectReplacements/Object 1305}es
  falso a la vez. La notación más habitual es
  \includegraphics[width=2.506cm,height=0.489cm]{./ObjectReplacements/Object 1088}y\includegraphics[width=3.29cm,height=0.467cm]{./ObjectReplacements/Object 1089}.
  Para el `no'
  (negación,\includegraphics[width=0.482cm,height=0.467cm]{./ObjectReplacements/Object 1306})
  tenemos
  \includegraphics[width=3.628cm,height=0.467cm]{./ObjectReplacements/Object 1090}.
  El conjunto de Boole es el conjunto de las proposiciones de la que
  partamos.

  De manera más formal se considera un elemento del álgebra de Boole de
  la lógica a las clases de equivalencia de las proposiciones
  equivalentes lógicamente (en su valor de verdad o falsedad) entre sí.

  \textbf{{[}Ejemplo 5{]}} El álgebra de conmutación (Shannon). Esta es
  el álgebra de Boole más sencilla que
  hay:\includegraphics[width=2.612cm,height=0.61cm]{./ObjectReplacements/Object 1071}.
  Las operaciones las concretaremos en tablas:

  \includegraphics[width=5.78cm,height=1.623cm]{./ObjectReplacements/Object 1091}

  y podréis comprobar fácilmente que se cumplen todos los postulados de
  Huntington. Ésta será usada frecuentemente durante el curso. Esta
  álgebra está contenida en todo álgebra de Boole.

  Como diagrama de Hasse (dónde se hace patente un orden en el álgebra)
  tiene
\end{enumerate}

\includegraphics[width=0.273cm,height=3.434cm]{Pictures/200000170000011100000D6AF113C5679B46B21D.svm}

\begin{enumerate}
\def\labelenumi{\arabic{enumi}.}
\setcounter{enumi}{2}
\item
  \textbf{{[}Ejemplo 6{]}} El álgebra de Boole de 4 elementos. Este es
  el álgebra de Boole generada por un conjunto de 2 elementos. Es
  singular en el sentido que sólo tiene 3 niveles, el más
  bajo\[\{{0 ≝ {\{\}} \equiv \varnothing}\}\], el
  intermedio\includegraphics[width=1.842cm,height=0.492cm]{./ObjectReplacements/Object 1200},
  y el
  superior\includegraphics[width=2.111cm,height=0.492cm]{./ObjectReplacements/Object 1202}.
  \includegraphics[width=3.471cm,height=0.61cm]{./ObjectReplacements/Object 1095}.
  Las operaciones las concretaremos en tablas:

  \includegraphics[width=9.488cm,height=2.732cm]{./ObjectReplacements/Object 1096}y
  podréis comprobar fácilmente que se cumplen todos los postulados de
  Huntington si
  cam\-biáis\includegraphics[width=0.441cm,height=0.467cm]{./ObjectReplacements/Object 1203}por\includegraphics[width=0.811cm,height=0.48cm]{./ObjectReplacements/Object 1204},\includegraphics[width=0.425cm,height=0.467cm]{./ObjectReplacements/Object 1205}por\includegraphics[width=0.757cm,height=0.48cm]{./ObjectReplacements/Object 1206},\includegraphics[width=0.411cm,height=0.467cm]{./ObjectReplacements/Object 1744}por\includegraphics[width=1.842cm,height=0.492cm]{./ObjectReplacements/Object 1207}y\includegraphics[width=0.423cm,height=0.467cm]{./ObjectReplacements/Object 1745}por
  el conjunto
  vacío\includegraphics[width=0.57cm,height=0.385cm]{./ObjectReplacements/Object 1208}.
  Su diagrama de Hasse es:
\end{enumerate}

\includegraphics[width=2.822cm,height=3.328cm]{Pictures/2000003600000B0600000D019F08206D3B15AC87.svm}

\begin{enumerate}
\def\labelenumi{\arabic{enumi}.}
\setcounter{enumi}{3}
\item
  \textbf{{[}Ejemplo 7{]}} El álgebra de Boole de 8 elementos. Este es
  el álgebra de Boole generada por un conjunto de 3 elementos. Es
  singular en el sentido que sólo tiene 4 niveles, el más
  bajo\[\{ 0\}\], el de
  átomos\includegraphics[width=1.64cm,height=0.492cm]{./ObjectReplacements/Object 1209},
  el de
  hiperátomos\includegraphics[width=1.884cm,height=0.492cm]{./ObjectReplacements/Object 1210}y
  el superior\[\{ 1\}\]. Los niveles de átomos y de hiperátomos son
  especialmente importantes, siendo esta álge\-bra de Boole, la más
  pequeña que los diferencia. Sería:

  \includegraphics[width=5.38cm,height=0.61cm]{./ObjectReplacements/Object 1097}

  Las operaciones las concretaremos en tablas:

  \includegraphics[width=12.206cm,height=4.948cm]{./ObjectReplacements/Object 1629}

  \includegraphics[width=15.085cm,height=3.577cm]{./ObjectReplacements/Object 1630}

  y podréis comprobar fácilmente que se cumplen todos los postulados de
  Huntington si te\-néis en cuenta los cambios aconsejados en el cuadro
  entre llaves, dónde las flechas quieren decir ``substituir por''.

  El diagrama de Hasse correspondiente sería:
\end{enumerate}

\includegraphics[width=3.478cm,height=4.634cm]{Pictures/2000003200000D9600001219FA1521857C60DA85.svm}

\begin{enumerate}
\def\labelenumi{\arabic{enumi}.}
\setcounter{enumi}{4}
\item
  \textbf{{[}Ejemplo 8{]}} El álgebra de Boole de 16 elementos. Este es
  el álgebra de Boole generada por un conjunto de 4 elementos. Es ya un
  álgebra de Boole completamente regular. Tiene 5 niveles, el más bajo
  el\[\{ 0\}\], el de
  átomos\includegraphics[width=2.124cm,height=0.492cm]{./ObjectReplacements/Object 1101},
  el de hiperátomos\[\{{Α,Β,\Gamma,\Delta}\}\], el
  intermedio\includegraphics[width=2.833cm,height=0.492cm]{./ObjectReplacements/Object 1103}y
  finalmente el nivel superior con el\[\{ 1\}\]. Sería:

  \includegraphics[width=9.453cm,height=0.61cm]{./ObjectReplacements/Object 1099}

  Las operaciones las concretaremos en tablas:

  \includegraphics[width=12.06cm,height=9.444cm]{./ObjectReplacements/Object 1100}
\end{enumerate}

\begin{quote}
\includegraphics[width=11.18cm,height=1.102cm]{./ObjectReplacements/Object 1105}
\end{quote}

\begin{quote}
\includegraphics[width=11.728cm,height=9.604cm]{./ObjectReplacements/Object 1628}
\end{quote}

\begin{enumerate}
\def\labelenumi{\arabic{enumi}.}
\setcounter{enumi}{5}
\item
  y podréis comprobar fácilmente que se cumplen todos los postulados de
  Huntington, con solo tener en cuenta que todos los elementos se pueden
  poner en función de
  \includegraphics[width=1.958cm,height=0.483cm]{./ObjectReplacements/Object 1509}y
  sumas de ellos. Las sumas de dos de los anteriores elementos son
  \includegraphics[width=2.94cm,height=0.467cm]{./ObjectReplacements/Object 1510}y
  las sumas de tres de ellos
  son\includegraphics[width=2.168cm,height=0.467cm]{./ObjectReplacements/Object 1511}.

  Como diagrama de Hasse tiene:
\end{enumerate}

\includegraphics[width=7.128cm,height=5.355cm]{Pictures/200000D100001BD7000014EC91BF966ED9304A6C.svm}

\begin{enumerate}
\def\labelenumi{\arabic{enumi}.}
\setcounter{enumi}{6}
\item
  \textbf{{[}Ejemplo 9{]}} El álgebra de Boole de los conjuntos que se
  pueden expresar como unión disyunta finita de sub-intervalos genéricos
  de\includegraphics[width=1.639cm,height=0.517cm]{./ObjectReplacements/Object 1250}.
  Definimos por conveniencia
  \includegraphics[width=2.006cm,height=0.54cm]{./ObjectReplacements/Object 1631}.
  Para esto haremos abstracción de cualquier conjunto finito de puntos
  de
  \includegraphics[width=0.58cm,height=0.515cm]{./ObjectReplacements/Object 1098},
  esto es, consideraremos que dos conjuntos son equivalentes si su
  diferencia simétrica (la unión de las diferencias, los elementos que
  no son comunes de ambos conjuntos) es vacía o es un conjunto finito de
  puntos. Esta álgebra de Boole tiene un cardinal infinito numerable
  (como el cardinal de los números naturales). Lo más importante es que
  no puede desarrollarse de manera semejante a como desarrolla\-mos el
  álgebra de las partes de un conjunto. Lo formalizaremos del siguiente
  modo:

  \begin{enumerate}
  \def\labelenumii{\arabic{enumii}.}
  \setcounter{enumii}{2}
  \item
    \includegraphics[width=9.567cm,height=0.723cm]{./ObjectReplacements/Object 1259}

    Si
    escribimos\includegraphics[width=1.304cm,height=0.54cm]{./ObjectReplacements/Object 1260}entonces
    \includegraphics[width=2.127cm,height=0.467cm]{./ObjectReplacements/Object 1261}

    \begin{enumerate}
    \def\labelenumiii{\arabic{enumiii}.}
    \item
      Sea\includegraphics[width=6.415cm,height=0.753cm]{./ObjectReplacements/Object 1632}
    \item
      Sea\includegraphics[width=12.719cm,height=0.69cm]{./ObjectReplacements/Object 1634}
    \item
      Sea\includegraphics[width=6.976cm,height=0.766cm]{./ObjectReplacements/Object 1633}
    \item
      \includegraphics[width=7.354cm,height=0.67cm]{./ObjectReplacements/Object 1252}Esta
      relación es de equivalencia.

      \begin{enumerate}
      \def\labelenumiv{\arabic{enumiv}.}
      \item
        Reflexiva
        \includegraphics[width=4.24cm,height=0.556cm]{./ObjectReplacements/Object 1262}Luego\includegraphics[width=1.177cm,height=0.467cm]{./ObjectReplacements/Object 1263}
      \item
        Simétrica\includegraphics[width=10.195cm,height=0.519cm]{./ObjectReplacements/Object 1264}
      \item
        Transitiva\includegraphics[width=3.946cm,height=0.467cm]{./ObjectReplacements/Object 1265}

        \begin{enumerate}
        \def\labelenumv{\arabic{enumv}.}
        \tightlist
        \item
          \includegraphics[width=11.404cm,height=0.556cm]{./ObjectReplacements/Object 20}Y
          queda de\-mostrada la propiedad transitiva.
        \end{enumerate}
      \end{enumerate}
    \item
      A partir de aquí hablaremos de
      \includegraphics[width=0.847cm,height=0.482cm]{./ObjectReplacements/Object 1254}para
      hablar de la clase de equivalencia de
      \includegraphics[width=1.542cm,height=0.515cm]{./ObjectReplacements/Object 1268}bajo
      la relación de
      equivalencia\includegraphics[width=0.579cm,height=0.467cm]{./ObjectReplacements/Object 1636}.
    \item
      A partir de aquí hablaremos de nuestro conjunto de
      Boole\includegraphics[width=4.553cm,height=0.601cm]{./ObjectReplacements/Object 1635}.
    \item
      Nuestro conjunto de Boole será
      \includegraphics[width=2.431cm,height=0.594cm]{./ObjectReplacements/Object 1269}.
    \item
      El\includegraphics[width=1.483cm,height=0.482cm]{./ObjectReplacements/Object 1270}
    \item
      El\includegraphics[width=1.665cm,height=0.587cm]{./ObjectReplacements/Object 1271}
    \item
      Ahora veremos unas operaciones muy cercanas a la unión, la
      intersección y el com\-plemento, que realmente nos dan un álgebra
      de Boole sobre
      \includegraphics[width=0.945cm,height=0.575cm]{./ObjectReplacements/Object 1637}:

      \includegraphics[width=3.201cm,height=2.108cm]{./ObjectReplacements/Object 1272}
    \item
      Convenio de
      notación:\[{⟦{a,b}⟧} ≝ {⟦\left\lbrack {a,b} \right\rbrack_{\mathbb{Q}}⟧}\].
      Estos conjuntos serán nuestros sub-intervalos genéricos del
      intervalo-unidad genéri\-co.
    \item
      Sea una sucesión finita de un número
      par\includegraphics[width=0.794cm,height=0.467cm]{./ObjectReplacements/Object 1218}de
      elementos
      de\includegraphics[width=1.155cm,height=0.54cm]{./ObjectReplacements/Object 1217},
      estricta\-mente
      creciente\includegraphics[width=6.034cm,height=0.531cm]{./ObjectReplacements/Object 1219}dispuestos
      como

      \includegraphics[width=4.41cm,height=0.531cm]{./ObjectReplacements/Object 1220}definirán
      los elementos
      de\includegraphics[width=0.945cm,height=0.575cm]{./ObjectReplacements/Object 1221},
      aparte de
      \includegraphics[width=5.011cm,height=0.639cm]{./ObjectReplacements/Object 1222}
    \item
      Si
      escribimos\includegraphics[width=4.41cm,height=0.531cm]{./ObjectReplacements/Object 1277},
      significamos ya (suponemos que es un hecho
      que)\includegraphics[width=6.034cm,height=0.531cm]{./ObjectReplacements/Object 1278}
    \item
      Ahora ya definimos (notación):
      \includegraphics[width=12.515cm,height=1.141cm]{./ObjectReplacements/Object 1279}
    \item
      Ahora ya tenemos el conjunto de Boole que buscábamos:

      \includegraphics[width=12.248cm,height=0.586cm]{./ObjectReplacements/Object 1224}
    \end{enumerate}
  \end{enumerate}

  Que las uniones, complementos e intersecciones de intervalos genéricos
  finitos siguen siendo intervalos genéricos finitos es claro desde el
  principio. Sin embargo voy a exponer la cabalística, hacer las cuentas
  vamos, para que no quede lugar a dudas. Toda esta comprobación (o
  redefinición) de
  que\includegraphics[width=0.51cm,height=0.563cm]{./ObjectReplacements/Object 1273}es
  cerrado bajo las distintas operaciones es quizás demasiado laboriosa,
  enojosa, pero quiero dejar claro este ejemplo de un álgebra de Boole
  no atómica, que por otra parte es la única vez que veremos.

  La operación de complemento queda de la siguiente manera, y aunque aún
  no podemos comprobar su corrección, si queda claro que es un operación
  un-aria interna:

  \includegraphics[width=9.992cm,height=5.239cm]{./ObjectReplacements/Object 1226}

  De dónde
  obtenemos\includegraphics[width=3.311cm,height=0.563cm]{./ObjectReplacements/Object 1229}.

  Tenemos que
  \includegraphics[width=3.471cm,height=0.667cm]{./ObjectReplacements/Object 1230}y\includegraphics[width=2.942cm,height=0.667cm]{./ObjectReplacements/Object 1231}y\includegraphics[width=2.26cm,height=0.517cm]{./ObjectReplacements/Object 1274}Además
  observamos con claridad
  que\includegraphics[width=2.932cm,height=0.563cm]{./ObjectReplacements/Object 1276}tal
  que\includegraphics[width=2.783cm,height=0.564cm]{./ObjectReplacements/Object 1638}y\includegraphics[width=2.265cm,height=0.529cm]{./ObjectReplacements/Object 1639},
  además
  de\includegraphics[width=2.912cm,height=0.563cm]{./ObjectReplacements/Object 1275}.
  Así nos queda
  \includegraphics[width=1.221cm,height=0.517cm]{./ObjectReplacements/Object 1640}
  si se verifican los demás axiomas.

  La suma quedará de la siguiente forma:

  \includegraphics[width=9.246cm,height=0.552cm]{./ObjectReplacements/Object 1286}

  El producto seguirá un camino par:

  \includegraphics[width=9.093cm,height=0.552cm]{./ObjectReplacements/Object 1287}

  Sólo queda ver que efectivamente las operaciones son internas:

  Prueba:

  \begin{enumerate}
  \def\labelenumii{\arabic{enumii}.}
  \item
    \begin{enumerate}
    \def\labelenumiii{\arabic{enumiii}.}
    \item
      Ahora vamos a desarrollar la suma de forma recurrente:

      \includegraphics[width=10.832cm,height=0.591cm]{./ObjectReplacements/Object 1289}

      Comenzaremos por
      \includegraphics[width=1.101cm,height=0.467cm]{./ObjectReplacements/Object 1288}
      y algunos casos especiales:
    \end{enumerate}
  \end{enumerate}

  \includegraphics[width=12.935cm,height=8.5cm]{./ObjectReplacements/Object 1234}

  Distintos casos con
  \includegraphics[width=1.087cm,height=0.467cm]{./ObjectReplacements/Object 1232}

  \includegraphics[width=9.23cm,height=2.05cm]{./ObjectReplacements/Object 1247}\includegraphics[width=9.215cm,height=2.789cm]{./ObjectReplacements/Object 1216}
  \includegraphics[width=9.2cm,height=2.05cm]{./ObjectReplacements/Object 1283}\includegraphics[width=9.809cm,height=2.789cm]{./ObjectReplacements/Object 1225}\includegraphics[width=9.486cm,height=3.378cm]{./ObjectReplacements/Object 1235}\includegraphics[width=9.744cm,height=3.378cm]{./ObjectReplacements/Object 1236}\includegraphics[width=13.213cm,height=4.119cm]{./ObjectReplacements/Object 1282}\includegraphics[width=11.913cm,height=4.119cm]{./ObjectReplacements/Object 1238}\includegraphics[width=11.88cm,height=4.119cm]{./ObjectReplacements/Object 1237}\includegraphics[width=11.578cm,height=4.119cm]{./ObjectReplacements/Object 1239}
  \includegraphics[width=12.305cm,height=4.119cm]{./ObjectReplacements/Object 1285}\includegraphics[width=12.305cm,height=4.119cm]{./ObjectReplacements/Object 1281}\includegraphics[width=11.915cm,height=4.119cm]{./ObjectReplacements/Object 1240}

  Para el
  caso\includegraphics[width=1.087cm,height=0.467cm]{./ObjectReplacements/Object 1644}o\includegraphics[width=1.101cm,height=0.467cm]{./ObjectReplacements/Object 1645}y
  especiales queda demostrada el cerramiento de
  \includegraphics[width=0.51cm,height=0.563cm]{./ObjectReplacements/Object 1646}bajo
  esta suma reducida. El caso siguiente se construye con facilidad por
  recurrencia en cualquier número finito de pasos. Si hacemos sumas
  comprobadas ya un número de veces finita, queda claro que ya no hay
  más que demostrar.

  Para
  cualquier\includegraphics[width=1.027cm,height=0.467cm]{./ObjectReplacements/Object 1284}:

  \includegraphics[width=12.37cm,height=3.38cm]{./ObjectReplacements/Object 1241}

  Queda demostrado que toda suma da como resultado un conjunto finito de
  intervalos gené\-ricos.

  A su vez el producto lo vamos a definir de forma recursiva también,
  comenzando primero
  con\includegraphics[width=0.499cm,height=0.467cm]{./ObjectReplacements/Object 1228}siendo
  la clase de un solo intervalo genérico, o la clase del vacío, además
  de algu\-nos casos especiales.

  \includegraphics[width=12.913cm,height=8.468cm]{./ObjectReplacements/Object 1242}\includegraphics[width=11.264cm,height=3.378cm]{./ObjectReplacements/Object 1243}\includegraphics[width=11.728cm,height=3.378cm]{./ObjectReplacements/Object 1244}\includegraphics[width=11.262cm,height=4.119cm]{./ObjectReplacements/Object 1290}\includegraphics[width=10.927cm,height=4.119cm]{./ObjectReplacements/Object 1292}\includegraphics[width=9.67cm,height=3.378cm]{./ObjectReplacements/Object 1291}\includegraphics[width=9.987cm,height=4.119cm]{./ObjectReplacements/Object 1245}\includegraphics[width=9.971cm,height=4.119cm]{./ObjectReplacements/Object 1246}\includegraphics[width=10.582cm,height=4.119cm]{./ObjectReplacements/Object 1248}

  Para el
  caso\includegraphics[width=1.087cm,height=0.467cm]{./ObjectReplacements/Object 1641}o\includegraphics[width=1.101cm,height=0.467cm]{./ObjectReplacements/Object 1642}y
  especiales queda demostrada el cerramiento de
  \includegraphics[width=0.51cm,height=0.563cm]{./ObjectReplacements/Object 1643}bajo
  esta suma reducida. El caso siguiente se construye con facilidad por
  recurrencia en cualquier número finito de pasos. Si hacemos sumas
  comprobadas ya un número de veces finita, queda claro que ya no hay
  más que demostrar.

  Para el caso
  \includegraphics[width=1.027cm,height=0.467cm]{./ObjectReplacements/Object 1280}:

  \includegraphics[width=9.313cm,height=3.38cm]{./ObjectReplacements/Object 1233}

  Y queda demostrado que la forma del conjunto producto es una clase de
  unión finita de sub-intervalos genéricos. Luego pertenece a nuestro
  álgebra de Boole.

  Este sistema es muy parecido a un álgebra de conjuntos subálgebra de
  algún conjunto po\-tencia, por lo que es fácil determinar que se trata
  de un álgebra de Boole. Sin embargo su cardinal
  es\includegraphics[width=4.77cm,height=0.476cm]{./ObjectReplacements/Object 1249}.
  Veámoslo:

  \includegraphics[width=13.457cm,height=4.768cm]{./ObjectReplacements/Object 1647}\includegraphics[width=10.114cm,height=1.683cm]{./ObjectReplacements/Object 1648}

  En definitiva es un álgebra numerable (del mismo cardinal que los
  números naturales). Puesto que el cardinal de los números naturales
  \includegraphics[width=0.533cm,height=0.385cm]{./ObjectReplacements/Object 1227}
  no es conmensurable con el de la potencia de ningún conjunto (es del
  cardinal infinito más pequeño posible y ningún conjunto finito tiene
  como potencia uno infinito), no existe ningún conjunto para el cual
  esta álgebra de Boole sea semejante (isomorfa) a un álgebra de las
  potencias de un conjunto. Este ejemplo será de utilidad más adelante,
  además de darnos un curioso ejemplo de álgebra de Boole nada común.

  \textbf{{[}Ejemplo 10{]}} El álgebra de las funciones de un álgebra de
  Boole sobre otra. Supongamos
  \includegraphics[width=1.88cm,height=0.563cm]{./ObjectReplacements/Object 1612}dónde\includegraphics[width=0.515cm,height=0.467cm]{./ObjectReplacements/Object 1613}es
  una función. Llamaremos
  \includegraphics[width=9.442cm,height=0.577cm]{./ObjectReplacements/Object 1614}a
  nuestro conjunto de Boole.
  \includegraphics[width=0.811cm,height=0.601cm]{./ObjectReplacements/Object 1617}es
  la función que asigna el cero
  de\includegraphics[width=0.683cm,height=0.563cm]{./ObjectReplacements/Object 1616}a
  cualquier elemento
  de\includegraphics[width=0.51cm,height=0.563cm]{./ObjectReplacements/Object 1615}.
  Igualmente
  \includegraphics[width=0.803cm,height=0.601cm]{./ObjectReplacements/Object 1618}es
  la función que asigna el uno
  de\includegraphics[width=0.683cm,height=0.563cm]{./ObjectReplacements/Object 1623}a
  cualquier elemento
  de\includegraphics[width=0.51cm,height=0.563cm]{./ObjectReplacements/Object 1624}.
  Las operaciones internas a introducir son:

  La suma de
  funciones:\includegraphics[width=5.898cm,height=0.577cm]{./ObjectReplacements/Object 1625}

  La multiplicación de
  funciones:\includegraphics[width=5.592cm,height=0.577cm]{./ObjectReplacements/Object 1739}

  La función
  complemento:\includegraphics[width=7.542cm,height=0.577cm]{./ObjectReplacements/Object 1740}

  De esta álgebra podemos entresacar otros conjuntos de funciones
  interesantes, como. Por ejemplo:

  \includegraphics[width=14.418cm,height=4.537cm]{./ObjectReplacements/Object 1741}

  y aún otros subconjuntos más pequeños serían las inyecciones de los
  anteriores homomorfismos. Si
  \includegraphics[width=1.344cm,height=0.563cm]{./ObjectReplacements/Object 1743}entonces
  uno de los conjuntos de aplicaciones más interesantes son los
  endomorfimos o isomorfimos en sí mismo.

  Además el kernel de cualquier homomorfismo es un subálgebra
  de\includegraphics[width=0.51cm,height=0.563cm]{./ObjectReplacements/Object 1742}.
  Los homomorfismos de las álgebras booleanas tienen propiedades
  interesantes que no veremos aquí.

  \begin{enumerate}
  \def\labelenumii{\arabic{enumii}.}
  \setcounter{enumii}{1}
  \item
    Desarrollo de las propiedades generales del álgebra de Boole a
    partir de los axiomas de Huntington.

    \begin{enumerate}
    \def\labelenumiii{\arabic{enumiii}.}
    \item
      Teoremas de Idempotencia:

      \begin{enumerate}
      \def\labelenumiv{\arabic{enumiv}.}
      \tightlist
      \item
        Para la suma:
        \includegraphics[width=3.644cm,height=0.563cm]{./ObjectReplacements/Object 1782}
      \end{enumerate}
    \end{enumerate}
  \end{enumerate}
\end{enumerate}

\begin{quote}
Prueba:
\end{quote}

\begin{enumerate}
\def\labelenumi{\arabic{enumi}.}
\setcounter{enumi}{7}
\tightlist
\item
  \begin{enumerate}
  \def\labelenumii{\arabic{enumii}.}
  \tightlist
  \item
    \begin{enumerate}
    \def\labelenumiii{\arabic{enumiii}.}
    \tightlist
    \item
      \begin{enumerate}
      \def\labelenumiv{\arabic{enumiv}.}
      \tightlist
      \item
        \begin{enumerate}
        \def\labelenumv{\arabic{enumv}.}
        \tightlist
        \item
          \includegraphics[width=3.794cm,height=0.467cm]{./ObjectReplacements/Object 1784}
        \end{enumerate}
      \end{enumerate}
    \end{enumerate}
  \end{enumerate}
\end{enumerate}

\begin{quote}
Elemento neutro de la suma.
\end{quote}

\begin{enumerate}
\def\labelenumi{\arabic{enumi}.}
\setcounter{enumi}{8}
\tightlist
\item
  \begin{enumerate}
  \def\labelenumii{\arabic{enumii}.}
  \tightlist
  \item
    \begin{enumerate}
    \def\labelenumiii{\arabic{enumiii}.}
    \tightlist
    \item
      \begin{enumerate}
      \def\labelenumiv{\arabic{enumiv}.}
      \tightlist
      \item
        \begin{enumerate}
        \def\labelenumv{\arabic{enumv}.}
        \tightlist
        \item
          \includegraphics[width=4.524cm,height=0.506cm]{./ObjectReplacements/Object 1785}
        \end{enumerate}
      \end{enumerate}
    \end{enumerate}
  \end{enumerate}
\end{enumerate}

\begin{quote}
Propiedad del complementario que produce 0.
\end{quote}

\begin{enumerate}
\def\labelenumi{\arabic{enumi}.}
\setcounter{enumi}{9}
\tightlist
\item
  \begin{enumerate}
  \def\labelenumii{\arabic{enumii}.}
  \tightlist
  \item
    \begin{enumerate}
    \def\labelenumiii{\arabic{enumiii}.}
    \tightlist
    \item
      \begin{enumerate}
      \def\labelenumiv{\arabic{enumiv}.}
      \tightlist
      \item
        \begin{enumerate}
        \def\labelenumv{\arabic{enumv}.}
        \tightlist
        \item
          \includegraphics[width=5.35cm,height=0.506cm]{./ObjectReplacements/Object 1786}
        \end{enumerate}
      \end{enumerate}
    \end{enumerate}
  \end{enumerate}
\end{enumerate}

\begin{quote}
Distributiva izquierda de la suma respecto al producto.
\end{quote}

\begin{enumerate}
\def\labelenumi{\arabic{enumi}.}
\setcounter{enumi}{10}
\tightlist
\item
  \begin{enumerate}
  \def\labelenumii{\arabic{enumii}.}
  \tightlist
  \item
    \begin{enumerate}
    \def\labelenumiii{\arabic{enumiii}.}
    \tightlist
    \item
      \begin{enumerate}
      \def\labelenumiv{\arabic{enumiv}.}
      \tightlist
      \item
        \begin{enumerate}
        \def\labelenumv{\arabic{enumv}.}
        \tightlist
        \item
          \includegraphics[width=4.503cm,height=0.506cm]{./ObjectReplacements/Object 1787}
        \end{enumerate}
      \end{enumerate}
    \end{enumerate}
  \end{enumerate}
\end{enumerate}

\begin{quote}
Propiedad del complementario que suma 1.
\end{quote}

\begin{enumerate}
\def\labelenumi{\arabic{enumi}.}
\setcounter{enumi}{11}
\item
  \begin{enumerate}
  \def\labelenumii{\arabic{enumii}.}
  \item
    \begin{enumerate}
    \def\labelenumiii{\arabic{enumiii}.}
    \item
      \begin{enumerate}
      \def\labelenumiv{\arabic{enumiv}.}
      \item
        \begin{enumerate}
        \def\labelenumv{\arabic{enumv}.}
        \item
          \includegraphics[width=3.512cm,height=0.467cm]{./ObjectReplacements/Object 1788}

          Elemento neutro del producto.
        \end{enumerate}
      \end{enumerate}
    \item
      Para el producto:
      \includegraphics[width=3.491cm,height=0.563cm]{./ObjectReplacements/Object 1790}
    \end{enumerate}
  \end{enumerate}
\end{enumerate}

\begin{quote}
Prueba:
\end{quote}

\begin{enumerate}
\def\labelenumi{\arabic{enumi}.}
\setcounter{enumi}{12}
\tightlist
\item
  \begin{enumerate}
  \def\labelenumii{\arabic{enumii}.}
  \tightlist
  \item
    \begin{enumerate}
    \def\labelenumiii{\arabic{enumiii}.}
    \tightlist
    \item
      \begin{enumerate}
      \def\labelenumiv{\arabic{enumiv}.}
      \tightlist
      \item
        \begin{enumerate}
        \def\labelenumv{\arabic{enumv}.}
        \tightlist
        \item
          \includegraphics[width=3.628cm,height=0.467cm]{./ObjectReplacements/Object 1791}
        \end{enumerate}
      \end{enumerate}
    \end{enumerate}
  \end{enumerate}
\end{enumerate}

\begin{quote}
Elemento neutro de la suma.
\end{quote}

\begin{enumerate}
\def\labelenumi{\arabic{enumi}.}
\setcounter{enumi}{13}
\tightlist
\item
  \begin{enumerate}
  \def\labelenumii{\arabic{enumii}.}
  \tightlist
  \item
    \begin{enumerate}
    \def\labelenumiii{\arabic{enumiii}.}
    \tightlist
    \item
      \begin{enumerate}
      \def\labelenumiv{\arabic{enumiv}.}
      \tightlist
      \item
        \begin{enumerate}
        \def\labelenumv{\arabic{enumv}.}
        \tightlist
        \item
          \includegraphics[width=4.524cm,height=0.506cm]{./ObjectReplacements/Object 1792}
        \end{enumerate}
      \end{enumerate}
    \end{enumerate}
  \end{enumerate}
\end{enumerate}

\begin{quote}
Propiedad del complementario que produce 0.
\end{quote}

\begin{enumerate}
\def\labelenumi{\arabic{enumi}.}
\setcounter{enumi}{14}
\tightlist
\item
  \begin{enumerate}
  \def\labelenumii{\arabic{enumii}.}
  \tightlist
  \item
    \begin{enumerate}
    \def\labelenumiii{\arabic{enumiii}.}
    \tightlist
    \item
      \begin{enumerate}
      \def\labelenumiv{\arabic{enumiv}.}
      \tightlist
      \item
        \begin{enumerate}
        \def\labelenumv{\arabic{enumv}.}
        \tightlist
        \item
          \includegraphics[width=5.195cm,height=0.506cm]{./ObjectReplacements/Object 1793}
        \end{enumerate}
      \end{enumerate}
    \end{enumerate}
  \end{enumerate}
\end{enumerate}

\begin{quote}
Distributiva izquierda de la suma respecto al producto.
\end{quote}

\begin{enumerate}
\def\labelenumi{\arabic{enumi}.}
\setcounter{enumi}{15}
\tightlist
\item
  \begin{enumerate}
  \def\labelenumii{\arabic{enumii}.}
  \tightlist
  \item
    \begin{enumerate}
    \def\labelenumiii{\arabic{enumiii}.}
    \tightlist
    \item
      \begin{enumerate}
      \def\labelenumiv{\arabic{enumiv}.}
      \tightlist
      \item
        \begin{enumerate}
        \def\labelenumv{\arabic{enumv}.}
        \tightlist
        \item
          \includegraphics[width=4.517cm,height=0.506cm]{./ObjectReplacements/Object 1794}
        \end{enumerate}
      \end{enumerate}
    \end{enumerate}
  \end{enumerate}
\end{enumerate}

\begin{quote}
Propiedad del complementario que suma 1.
\end{quote}

\begin{enumerate}
\def\labelenumi{\arabic{enumi}.}
\setcounter{enumi}{16}
\tightlist
\item
  \begin{enumerate}
  \def\labelenumii{\arabic{enumii}.}
  \tightlist
  \item
    \begin{enumerate}
    \def\labelenumiii{\arabic{enumiii}.}
    \tightlist
    \item
      \begin{enumerate}
      \def\labelenumiv{\arabic{enumiv}.}
      \tightlist
      \item
        \begin{enumerate}
        \def\labelenumv{\arabic{enumv}.}
        \tightlist
        \item
          \includegraphics[width=3.358cm,height=0.467cm]{./ObjectReplacements/Object 1789}
        \end{enumerate}
      \end{enumerate}
    \end{enumerate}
  \end{enumerate}
\end{enumerate}

\begin{quote}
Elemento neutro del producto.
\end{quote}

\begin{enumerate}
\def\labelenumi{\arabic{enumi}.}
\setcounter{enumi}{17}
\item
  \begin{enumerate}
  \def\labelenumii{\arabic{enumii}.}
  \item
    \begin{enumerate}
    \def\labelenumiii{\arabic{enumiii}.}
    \item
      \begin{enumerate}
      \def\labelenumiv{\arabic{enumiv}.}
      \item
        Teoremas de Absorción:

        \begin{enumerate}
        \def\labelenumv{\arabic{enumv}.}
        \tightlist
        \item
          Para
          el\includegraphics[width=0.411cm,height=0.467cm]{./ObjectReplacements/Object 1807}:\includegraphics[width=3.581cm,height=0.563cm]{./ObjectReplacements/Object 1796}
        \end{enumerate}
      \end{enumerate}
    \end{enumerate}
  \end{enumerate}
\end{enumerate}

\begin{quote}
Prueba:
\end{quote}

\begin{enumerate}
\def\labelenumi{\arabic{enumi}.}
\setcounter{enumi}{18}
\tightlist
\item
  \begin{enumerate}
  \def\labelenumii{\arabic{enumii}.}
  \tightlist
  \item
    \begin{enumerate}
    \def\labelenumiii{\arabic{enumiii}.}
    \tightlist
    \item
      \begin{enumerate}
      \def\labelenumiv{\arabic{enumiv}.}
      \tightlist
      \item
        \begin{enumerate}
        \def\labelenumv{\arabic{enumv}.}
        \tightlist
        \item
          \begin{enumerate}
          \def\labelenumvi{\arabic{enumvi}.}
          \tightlist
          \item
            \includegraphics[width=3.939cm,height=0.467cm]{./ObjectReplacements/Object 1797}
          \end{enumerate}
        \end{enumerate}
      \end{enumerate}
    \end{enumerate}
  \end{enumerate}
\end{enumerate}

\begin{quote}
Propiedad de
sumar\includegraphics[width=0.411cm,height=0.467cm]{./ObjectReplacements/Object 546}de
un elemento al sumar con uno de sus complementarios.
\end{quote}

\begin{enumerate}
\def\labelenumi{\arabic{enumi}.}
\setcounter{enumi}{19}
\tightlist
\item
  \begin{enumerate}
  \def\labelenumii{\arabic{enumii}.}
  \tightlist
  \item
    \begin{enumerate}
    \def\labelenumiii{\arabic{enumiii}.}
    \tightlist
    \item
      \begin{enumerate}
      \def\labelenumiv{\arabic{enumiv}.}
      \tightlist
      \item
        \begin{enumerate}
        \def\labelenumv{\arabic{enumv}.}
        \tightlist
        \item
          \begin{enumerate}
          \def\labelenumvi{\arabic{enumvi}.}
          \tightlist
          \item
            \includegraphics[width=4.493cm,height=0.506cm]{./ObjectReplacements/Object 1798}
          \end{enumerate}
        \end{enumerate}
      \end{enumerate}
    \end{enumerate}
  \end{enumerate}
\end{enumerate}

\begin{quote}
Elemento neutro de la multiplicación.
\end{quote}

\begin{enumerate}
\def\labelenumi{\arabic{enumi}.}
\setcounter{enumi}{20}
\tightlist
\item
  \begin{enumerate}
  \def\labelenumii{\arabic{enumii}.}
  \tightlist
  \item
    \begin{enumerate}
    \def\labelenumiii{\arabic{enumiii}.}
    \tightlist
    \item
      \begin{enumerate}
      \def\labelenumiv{\arabic{enumiv}.}
      \tightlist
      \item
        \begin{enumerate}
        \def\labelenumv{\arabic{enumv}.}
        \tightlist
        \item
          \begin{enumerate}
          \def\labelenumvi{\arabic{enumvi}.}
          \tightlist
          \item
            \includegraphics[width=5.316cm,height=0.506cm]{./ObjectReplacements/Object 1799}
          \end{enumerate}
        \end{enumerate}
      \end{enumerate}
    \end{enumerate}
  \end{enumerate}
\end{enumerate}

\begin{quote}
Distributividad izquierda de la suma respecto al producto.
\end{quote}

\begin{enumerate}
\def\labelenumi{\arabic{enumi}.}
\setcounter{enumi}{21}
\tightlist
\item
  \begin{enumerate}
  \def\labelenumii{\arabic{enumii}.}
  \tightlist
  \item
    \begin{enumerate}
    \def\labelenumiii{\arabic{enumiii}.}
    \tightlist
    \item
      \begin{enumerate}
      \def\labelenumiv{\arabic{enumiv}.}
      \tightlist
      \item
        \begin{enumerate}
        \def\labelenumv{\arabic{enumv}.}
        \tightlist
        \item
          \begin{enumerate}
          \def\labelenumvi{\arabic{enumvi}.}
          \tightlist
          \item
            \includegraphics[width=4.471cm,height=0.506cm]{./ObjectReplacements/Object 1800}
          \end{enumerate}
        \end{enumerate}
      \end{enumerate}
    \end{enumerate}
  \end{enumerate}
\end{enumerate}

\begin{quote}
Propiedad de sumar
\includegraphics[width=0.411cm,height=0.467cm]{./ObjectReplacements/Object 547}de
un elemento con cualquiera de sus complementarios.
\end{quote}

\begin{enumerate}
\def\labelenumi{\arabic{enumi}.}
\setcounter{enumi}{22}
\tightlist
\item
  \begin{enumerate}
  \def\labelenumii{\arabic{enumii}.}
  \tightlist
  \item
    \begin{enumerate}
    \def\labelenumiii{\arabic{enumiii}.}
    \tightlist
    \item
      \begin{enumerate}
      \def\labelenumiv{\arabic{enumiv}.}
      \tightlist
      \item
        \begin{enumerate}
        \def\labelenumv{\arabic{enumv}.}
        \tightlist
        \item
          \begin{enumerate}
          \def\labelenumvi{\arabic{enumvi}.}
          \tightlist
          \item
            \includegraphics[width=3.48cm,height=0.467cm]{./ObjectReplacements/Object 1801}
          \end{enumerate}
        \end{enumerate}
      \end{enumerate}
    \end{enumerate}
  \end{enumerate}
\end{enumerate}

\begin{quote}
Elemento neutro del producto.
\end{quote}

\begin{enumerate}
\def\labelenumi{\arabic{enumi}.}
\setcounter{enumi}{23}
\tightlist
\item
  \begin{enumerate}
  \def\labelenumii{\arabic{enumii}.}
  \tightlist
  \item
    \begin{enumerate}
    \def\labelenumiii{\arabic{enumiii}.}
    \tightlist
    \item
      \begin{enumerate}
      \def\labelenumiv{\arabic{enumiv}.}
      \tightlist
      \item
        \begin{enumerate}
        \def\labelenumv{\arabic{enumv}.}
        \tightlist
        \item
          Para
          el\includegraphics[width=0.423cm,height=0.467cm]{./ObjectReplacements/Object 1802}:\includegraphics[width=3.454cm,height=0.563cm]{./ObjectReplacements/Object 1803}
        \end{enumerate}
      \end{enumerate}
    \end{enumerate}
  \end{enumerate}
\end{enumerate}

\begin{quote}
Prueba:
\end{quote}

\begin{enumerate}
\def\labelenumi{\arabic{enumi}.}
\setcounter{enumi}{24}
\tightlist
\item
  \begin{enumerate}
  \def\labelenumii{\arabic{enumii}.}
  \tightlist
  \item
    \begin{enumerate}
    \def\labelenumiii{\arabic{enumiii}.}
    \tightlist
    \item
      \begin{enumerate}
      \def\labelenumiv{\arabic{enumiv}.}
      \tightlist
      \item
        \begin{enumerate}
        \def\labelenumv{\arabic{enumv}.}
        \tightlist
        \item
          \begin{enumerate}
          \def\labelenumvi{\arabic{enumvi}.}
          \tightlist
          \item
            \includegraphics[width=3.798cm,height=0.467cm]{./ObjectReplacements/Object 1804}
          \end{enumerate}
        \end{enumerate}
      \end{enumerate}
    \end{enumerate}
  \end{enumerate}
\end{enumerate}

\begin{quote}
Propiedad de
resultar\includegraphics[width=0.423cm,height=0.467cm]{./ObjectReplacements/Object 548}de
un elemento al multiplicar con uno de sus complementarios.
\end{quote}

\begin{enumerate}
\def\labelenumi{\arabic{enumi}.}
\setcounter{enumi}{25}
\tightlist
\item
  \begin{enumerate}
  \def\labelenumii{\arabic{enumii}.}
  \tightlist
  \item
    \begin{enumerate}
    \def\labelenumiii{\arabic{enumiii}.}
    \tightlist
    \item
      \begin{enumerate}
      \def\labelenumiv{\arabic{enumiv}.}
      \tightlist
      \item
        \begin{enumerate}
        \def\labelenumv{\arabic{enumv}.}
        \tightlist
        \item
          \begin{enumerate}
          \def\labelenumvi{\arabic{enumvi}.}
          \tightlist
          \item
            \includegraphics[width=4.505cm,height=0.506cm]{./ObjectReplacements/Object 1805}
          \end{enumerate}
        \end{enumerate}
      \end{enumerate}
    \end{enumerate}
  \end{enumerate}
\end{enumerate}

\begin{quote}
Elemento neutro de la suma.
\end{quote}

\begin{enumerate}
\def\labelenumi{\arabic{enumi}.}
\setcounter{enumi}{26}
\tightlist
\item
  \begin{enumerate}
  \def\labelenumii{\arabic{enumii}.}
  \tightlist
  \item
    \begin{enumerate}
    \def\labelenumiii{\arabic{enumiii}.}
    \tightlist
    \item
      \begin{enumerate}
      \def\labelenumiv{\arabic{enumiv}.}
      \tightlist
      \item
        \begin{enumerate}
        \def\labelenumv{\arabic{enumv}.}
        \tightlist
        \item
          \begin{enumerate}
          \def\labelenumvi{\arabic{enumvi}.}
          \tightlist
          \item
            \includegraphics[width=5.175cm,height=0.506cm]{./ObjectReplacements/Object 1806}
          \end{enumerate}
        \end{enumerate}
      \end{enumerate}
    \end{enumerate}
  \end{enumerate}
\end{enumerate}

\begin{quote}
Distributividad izquierda del producto respecto a la suma.
\end{quote}

\begin{enumerate}
\def\labelenumi{\arabic{enumi}.}
\setcounter{enumi}{27}
\tightlist
\item
  \begin{enumerate}
  \def\labelenumii{\arabic{enumii}.}
  \tightlist
  \item
    \begin{enumerate}
    \def\labelenumiii{\arabic{enumiii}.}
    \tightlist
    \item
      \begin{enumerate}
      \def\labelenumiv{\arabic{enumiv}.}
      \tightlist
      \item
        \begin{enumerate}
        \def\labelenumv{\arabic{enumv}.}
        \tightlist
        \item
          \begin{enumerate}
          \def\labelenumvi{\arabic{enumvi}.}
          \tightlist
          \item
            \includegraphics[width=4.498cm,height=0.506cm]{./ObjectReplacements/Object 1808}
          \end{enumerate}
        \end{enumerate}
      \end{enumerate}
    \end{enumerate}
  \end{enumerate}
\end{enumerate}

\begin{quote}
Propiedad de
resultar\includegraphics[width=0.423cm,height=0.467cm]{./ObjectReplacements/Object 549}de
un elemento al multiplicar con uno de sus complementarios.
\end{quote}

\begin{enumerate}
\def\labelenumi{\arabic{enumi}.}
\setcounter{enumi}{28}
\tightlist
\item
  \begin{enumerate}
  \def\labelenumii{\arabic{enumii}.}
  \tightlist
  \item
    \begin{enumerate}
    \def\labelenumiii{\arabic{enumiii}.}
    \tightlist
    \item
      \begin{enumerate}
      \def\labelenumiv{\arabic{enumiv}.}
      \tightlist
      \item
        \begin{enumerate}
        \def\labelenumv{\arabic{enumv}.}
        \tightlist
        \item
          \begin{enumerate}
          \def\labelenumvi{\arabic{enumvi}.}
          \tightlist
          \item
            \[\mspace{468mu} = {x \cdot 0}\]
          \end{enumerate}
        \end{enumerate}
      \end{enumerate}
    \end{enumerate}
  \end{enumerate}
\end{enumerate}

\begin{quote}
Elemento neutro de la suma.
\end{quote}

\begin{enumerate}
\def\labelenumi{\arabic{enumi}.}
\setcounter{enumi}{29}
\tightlist
\item
  \begin{enumerate}
  \def\labelenumii{\arabic{enumii}.}
  \tightlist
  \item
    \begin{enumerate}
    \def\labelenumiii{\arabic{enumiii}.}
    \tightlist
    \item
      Condición necesaria y suficiente particular para que haya un solo
      elemento en el conjunto de
      Boole:\includegraphics[width=3.487cm,height=0.73cm]{./ObjectReplacements/Object 1809}.
    \end{enumerate}
  \end{enumerate}
\end{enumerate}

\begin{quote}
Prueba:
\end{quote}

\begin{enumerate}
\def\labelenumi{\arabic{enumi}.}
\setcounter{enumi}{30}
\tightlist
\item
  \begin{enumerate}
  \def\labelenumii{\arabic{enumii}.}
  \tightlist
  \item
    \begin{enumerate}
    \def\labelenumiii{\arabic{enumiii}.}
    \tightlist
    \item
      \begin{enumerate}
      \def\labelenumiv{\arabic{enumiv}.}
      \tightlist
      \item
        \includegraphics[width=1.709cm,height=0.587cm]{./ObjectReplacements/Object 1810}
      \end{enumerate}
    \end{enumerate}
  \end{enumerate}
\end{enumerate}

\begin{quote}
Propiedad de absorción de la multiplicación con
el\includegraphics[width=0.423cm,height=0.467cm]{./ObjectReplacements/Object 550}
\end{quote}

\begin{enumerate}
\def\labelenumi{\arabic{enumi}.}
\setcounter{enumi}{31}
\tightlist
\item
  \begin{enumerate}
  \def\labelenumii{\arabic{enumii}.}
  \tightlist
  \item
    \begin{enumerate}
    \def\labelenumiii{\arabic{enumiii}.}
    \tightlist
    \item
      \begin{enumerate}
      \def\labelenumiv{\arabic{enumiv}.}
      \tightlist
      \item
        \includegraphics[width=1.693cm,height=0.587cm]{./ObjectReplacements/Object 1811}
      \end{enumerate}
    \end{enumerate}
  \end{enumerate}
\end{enumerate}

\begin{quote}
Sustituimos
el\includegraphics[width=0.423cm,height=0.467cm]{./ObjectReplacements/Object 1746}por
el\includegraphics[width=0.411cm,height=0.467cm]{./ObjectReplacements/Object 1747}gracias
a la hipótesis.
\end{quote}

\begin{enumerate}
\def\labelenumi{\arabic{enumi}.}
\setcounter{enumi}{32}
\tightlist
\item
  \begin{enumerate}
  \def\labelenumii{\arabic{enumii}.}
  \tightlist
  \item
    \begin{enumerate}
    \def\labelenumiii{\arabic{enumiii}.}
    \tightlist
    \item
      \begin{enumerate}
      \def\labelenumiv{\arabic{enumiv}.}
      \tightlist
      \item
        \includegraphics[width=1.681cm,height=0.587cm]{./ObjectReplacements/Object 1812}
      \end{enumerate}
    \end{enumerate}
  \end{enumerate}
\end{enumerate}

\begin{quote}
Elemento neutro de la multiplicación,
el\includegraphics[width=0.411cm,height=0.467cm]{./ObjectReplacements/Object 1748}.
\end{quote}

\begin{enumerate}
\def\labelenumi{\arabic{enumi}.}
\setcounter{enumi}{33}
\tightlist
\item
  \begin{enumerate}
  \def\labelenumii{\arabic{enumii}.}
  \tightlist
  \item
    \begin{enumerate}
    \def\labelenumiii{\arabic{enumiii}.}
    \tightlist
    \item
      \begin{enumerate}
      \def\labelenumiv{\arabic{enumiv}.}
      \tightlist
      \item
        \includegraphics[width=0.796cm,height=0.467cm]{./ObjectReplacements/Object 1813}
      \item
        \includegraphics[width=1.349cm,height=0.573cm]{./ObjectReplacements/Object 1126}
      \end{enumerate}
    \end{enumerate}
  \end{enumerate}
\end{enumerate}

\begin{quote}
Desde el punto (1) y (3), mediante la transitividad de la igualdad
equivale.
\end{quote}

\begin{enumerate}
\def\labelenumi{\arabic{enumi}.}
\setcounter{enumi}{34}
\tightlist
\item
  \begin{enumerate}
  \def\labelenumii{\arabic{enumii}.}
  \tightlist
  \item
    \begin{enumerate}
    \def\labelenumiii{\arabic{enumiii}.}
    \tightlist
    \item
      \begin{enumerate}
      \def\labelenumiv{\arabic{enumiv}.}
      \tightlist
      \item
        \includegraphics[width=2.159cm,height=0.573cm]{./ObjectReplacements/Object 1814}
      \end{enumerate}
    \end{enumerate}
  \end{enumerate}
\end{enumerate}

\begin{quote}
Por la hipótesis.
\end{quote}

\begin{enumerate}
\def\labelenumi{\arabic{enumi}.}
\setcounter{enumi}{35}
\item
  \begin{enumerate}
  \def\labelenumii{\arabic{enumii}.}
  \item
    \begin{enumerate}
    \def\labelenumiii{\arabic{enumiii}.}
    \item
      Condición necesaria y suficiente general para que haya un solo
      elemento en el conjunto de
      Boole:\includegraphics[width=3.545cm,height=0.73cm]{./ObjectReplacements/Object 1816}.

      Prueba:

      \begin{enumerate}
      \def\labelenumiv{\arabic{enumiv}.}
      \item
        \includegraphics[width=4.274cm,height=0.467cm]{./ObjectReplacements/Object 1818}

        Axioma de
        sumar\includegraphics[width=0.411cm,height=0.467cm]{./ObjectReplacements/Object 1751}un
        elemento y su complementario.
      \item
        \includegraphics[width=4.281cm,height=0.467cm]{./ObjectReplacements/Object 1819}

        Particularización de la fórmula (2) usando la hipótesis.
      \item
        \includegraphics[width=4.533cm,height=0.467cm]{./ObjectReplacements/Object 1822}

        Teorema de Idempotencia de la suma.
      \item
        \includegraphics[width=4.509cm,height=0.467cm]{./ObjectReplacements/Object 1749}

        Fórmulas (1) y (3) aplicando la transitividad de la igualdad.
      \item
        \includegraphics[width=4.217cm,height=0.467cm]{./ObjectReplacements/Object 1820}

        Axioma de
        multiplicar\includegraphics[width=0.423cm,height=0.467cm]{./ObjectReplacements/Object 1750}un
        elemento y su complementario.
      \item
        \includegraphics[width=4.21cm,height=0.467cm]{./ObjectReplacements/Object 1821}

        Particularización de la fórmula (5) con usando la hipótesis.
      \item
        \includegraphics[width=4.533cm,height=0.467cm]{./ObjectReplacements/Object 1815}

        Teorema de Idempotencia del producto.
      \item
        \includegraphics[width=4.479cm,height=0.467cm]{./ObjectReplacements/Object 1823}

        Fórmulas (5) y (7) aplicando la transitividad de la igualdad.
      \item
        \includegraphics[width=4.531cm,height=0.563cm]{./ObjectReplacements/Object 1825}

        Fórmulas (4) y (8) y transitividad de la igualdad.
      \item
        \includegraphics[width=0.988cm,height=0.467cm]{./ObjectReplacements/Object 1826}

        Puesto que se da la hipótesis de (9) en (4) y (8).
      \item
        \includegraphics[width=2.159cm,height=0.573cm]{./ObjectReplacements/Object 1827}

        Por el teorema (4) (anterior).
      \end{enumerate}
    \item
      Unicidad del
      complementario:\includegraphics[width=4.106cm,height=0.563cm]{./ObjectReplacements/Object 9}

      Prueba:

      \begin{enumerate}
      \def\labelenumiv{\arabic{enumiv}.}
      \item
        \includegraphics[width=4.881cm,height=0.563cm]{./ObjectReplacements/Object 545}

        Axioma de
        sumar\includegraphics[width=0.411cm,height=0.467cm]{./ObjectReplacements/Object 1753}cualquier
        elemento y su complementario.
      \item
        \includegraphics[width=1.827cm,height=0.467cm]{./ObjectReplacements/Object 1752}

        Elemento neutro del producto.
      \item
        \includegraphics[width=2.644cm,height=0.506cm]{./ObjectReplacements/Object 1764}

        Particularización de la fórmula (1) (axioma).
      \item
        \includegraphics[width=3.334cm,height=0.506cm]{./ObjectReplacements/Object 1765}

        Distributiva izquierda del producto respecto de la suma.
      \item
        \includegraphics[width=2.626cm,height=0.506cm]{./ObjectReplacements/Object 1766}

        Propiedad de los complementarios por la que
        multiplican\includegraphics[width=0.423cm,height=0.467cm]{./ObjectReplacements/Object 1776}.
      \item
        \includegraphics[width=1.82cm,height=0.467cm]{./ObjectReplacements/Object 1767}

        Propiedad del elemento neutro de la suma.
      \item
        \includegraphics[width=1.813cm,height=0.467cm]{./ObjectReplacements/Object 122}

        Propiedad conmutativa del producto
      \item
        \includegraphics[width=2.626cm,height=0.506cm]{./ObjectReplacements/Object 1768}

        Propiedad del elemento neutro de la suma.
      \item
        \includegraphics[width=3.314cm,height=0.506cm]{./ObjectReplacements/Object 1769}

        Propiedad complementaria de
        \includegraphics[width=1.018cm,height=0.467cm]{./ObjectReplacements/Object 1777}que
        multiplican
        \includegraphics[width=0.423cm,height=0.467cm]{./ObjectReplacements/Object 1778}
      \item
        \includegraphics[width=2.644cm,height=0.506cm]{./ObjectReplacements/Object 1770}

        Distributiva izquierda del producto respecto de la suma.
      \item
        \includegraphics[width=1.769cm,height=0.467cm]{./ObjectReplacements/Object 1771}

        Propiedad complementaria de
        \includegraphics[width=1.018cm,height=0.467cm]{./ObjectReplacements/Object 1779}que
        suma
        \includegraphics[width=0.411cm,height=0.467cm]{./ObjectReplacements/Object 1780}
      \item
        \includegraphics[width=1.071cm,height=0.467cm]{./ObjectReplacements/Object 1772}

        Elemento neutro del producto
      \end{enumerate}
    \item
      Definimos el elemento complementario de otro, aquel elemento
      de\includegraphics[width=0.51cm,height=0.563cm]{./ObjectReplacements/Object 1373}que\includegraphics[width=4.385cm,height=0.573cm]{./ObjectReplacements/Object 1371},
      o dicho de otro
      modo,\[\forall{a \in}\qquad\exists!{\overline{a} \in}\mspace{72mu}{\left( {{a + \overline{a}} = 1} \right) \land \left( {{a \cdot \overline{a}} = 0} \right)}\],
      y este elemento es único.
    \item
      Complementación doble es
      identidad:\[\forall{x \in}{}{}{}{\overline{\overline{x}} = x}\].

      Prueba:

      \begin{enumerate}
      \def\labelenumiv{\arabic{enumiv}.}
      \item
        \includegraphics[width=8.216cm,height=0.573cm]{./ObjectReplacements/Object 123}

        Versión del axioma de complementarios con el teorema (6) y la
        definición (7).
      \item
        \includegraphics[width=8.264cm,height=0.573cm]{./ObjectReplacements/Object 551}

        Aplicamos particularización de (1)
        sustituyendo\includegraphics[width=0.443cm,height=0.467cm]{./ObjectReplacements/Object 1755}por\[\overline{x}\]y\includegraphics[width=0.443cm,height=0.467cm]{./ObjectReplacements/Object 1774}por\includegraphics[width=0.443cm,height=0.467cm]{./ObjectReplacements/Object 1773}.
      \item
        \includegraphics[width=8.264cm,height=0.573cm]{./ObjectReplacements/Object 552}

        Aplicamos los axiomas de conmutabilidad de la suma y el producto
        a las fórmulas de (2).
      \item
        \includegraphics[width=6.736cm,height=0.573cm]{./ObjectReplacements/Object 553}

        De las fórmulas (1) y (3) y del axioma de complementarios se
        deriva lo anterior.
      \item
        \includegraphics[width=5.595cm,height=0.467cm]{./ObjectReplacements/Object 627}

        Por el teorema (6) de unicidad del complementario.
      \end{enumerate}
    \item
      Unicidad de los neutros:

      \begin{enumerate}
      \def\labelenumiv{\arabic{enumiv}.}
      \item
        El neutro para la suma es único
        \includegraphics[width=6.773cm,height=0.681cm]{./ObjectReplacements/Object 14}

        Prueba:

        \begin{enumerate}
        \def\labelenumv{\arabic{enumv}.}
        \item
          \includegraphics[width=3.612cm,height=0.563cm]{./ObjectReplacements/Object 1781}

          Hipótesis.
        \item
          \includegraphics[width=3.56cm,height=0.467cm]{./ObjectReplacements/Object 1783}

          Particularización de (1) con el neutro de la suma.
        \item
          \includegraphics[width=3.56cm,height=0.467cm]{./ObjectReplacements/Object 1817}

          Axioma de conmutabilidad para la suma.
        \item
          \includegraphics[width=3.56cm,height=0.467cm]{./ObjectReplacements/Object 1828}

          Simetría de la igualdad.
        \item
          \includegraphics[width=4.135cm,height=0.467cm]{./ObjectReplacements/Object 1824}

          Axioma del elemento neutro de la suma.
        \item
          \includegraphics[width=4.152cm,height=0.467cm]{./ObjectReplacements/Object 1829}

          Transitividad de la igualdad en las fórmulas (4) y (5).
        \end{enumerate}
      \item
        El neutro para el producto es único
        \includegraphics[width=6.678cm,height=0.681cm]{./ObjectReplacements/Object 1775}

        Prueba:

        \begin{enumerate}
        \def\labelenumv{\arabic{enumv}.}
        \item
          \includegraphics[width=3.484cm,height=0.563cm]{./ObjectReplacements/Object 1830}

          Hipótesis.
        \item
          \includegraphics[width=3.403cm,height=0.467cm]{./ObjectReplacements/Object 1831}

          Particularización de (1) con el neutro del producto.
        \item
          \includegraphics[width=3.403cm,height=0.467cm]{./ObjectReplacements/Object 1832}

          Axioma de conmutatividad para el producto.
        \item
          \includegraphics[width=3.403cm,height=0.467cm]{./ObjectReplacements/Object 1833}

          Simetría de la igualdad.
        \item
          \includegraphics[width=4.015cm,height=0.467cm]{./ObjectReplacements/Object 1834}

          Axioma del elemento neutro del producto.
        \item
          \includegraphics[width=4.165cm,height=0.467cm]{./ObjectReplacements/Object 1835}

          Transitividad de la igualdad en las fórmulas (4) y (5).
        \end{enumerate}
      \end{enumerate}
    \item
      Para los elementos
      \includegraphics[width=0.982cm,height=0.467cm]{./ObjectReplacements/Object 1047}:
      \includegraphics[width=2.124cm,height=0.467cm]{./ObjectReplacements/Object 1048}

      Prueba:

      \begin{enumerate}
      \def\labelenumiv{\arabic{enumiv}.}
      \item
        \includegraphics[width=1.489cm,height=0.467cm]{./ObjectReplacements/Object 1757}

        Axioma de elemento neutro de la suma.
      \item
        \includegraphics[width=1.595cm,height=0.467cm]{./ObjectReplacements/Object 1049}

        Axioma de elemento neutro del producto.
      \end{enumerate}
    \item
      Leyes de cancelación:

      \begin{enumerate}
      \def\labelenumiv{\arabic{enumiv}.}
      \item
        Para la suma:

        \includegraphics[width=9.881cm,height=0.658cm]{./ObjectReplacements/Object 1003}

        Prueba:

        \begin{enumerate}
        \def\labelenumv{\arabic{enumv}.}
        \item
          \includegraphics[width=3.496cm,height=0.563cm]{./ObjectReplacements/Object 1060}

          Hipótesis (1).
        \item
          \includegraphics[width=3.498cm,height=0.563cm]{./ObjectReplacements/Object 1061}

          Hipótesis (2).
        \item
          \includegraphics[width=4.789cm,height=0.519cm]{./ObjectReplacements/Object 1062}

          Por ser el producto aplicación, el multiplicar las expresiones
          derechas de las igualdades entre sí y las izquierdas entre sí,
          no altera la igualdad.
        \item
          \includegraphics[width=3.219cm,height=0.519cm]{./ObjectReplacements/Object 1063}

          Por el axioma de distributividad de la suma respecto al
          producto por la izquierda.
        \item
          \includegraphics[width=2.074cm,height=0.467cm]{./ObjectReplacements/Object 1064}

          Por la propiedad de resultar
          \includegraphics[width=0.423cm,height=0.467cm]{./ObjectReplacements/Object 1836}la
          multiplicación de un elemento y su complementario en el axioma
          de complementarios.
        \item
          \includegraphics[width=1.05cm,height=0.467cm]{./ObjectReplacements/Object 1005}

          Por el axioma del elemento neutro de la suma.
        \end{enumerate}
      \item
        Para el
        producto:\includegraphics[width=9.266cm,height=0.658cm]{./ObjectReplacements/Object 1026}

        Prueba:

        \begin{enumerate}
        \def\labelenumv{\arabic{enumv}.}
        \item
          \includegraphics[width=3.187cm,height=0.563cm]{./ObjectReplacements/Object 1837}

          Hipótesis (1).
        \item
          \includegraphics[width=3.191cm,height=0.563cm]{./ObjectReplacements/Object 1838}

          Hipótesis (2).
        \item
          \includegraphics[width=4.48cm,height=0.519cm]{./ObjectReplacements/Object 1007}

          Por ser la suma aplicación, el sumar las expresiones derechas
          de las igualdades entre sí y las izquierdas entre sí, no
          altera la igualdad.
        \item
          \includegraphics[width=3.219cm,height=0.519cm]{./ObjectReplacements/Object 1008}

          Por el axioma de distributividad del producto respecto de la
          suma por la izquierda.
        \item
          \includegraphics[width=1.744cm,height=0.467cm]{./ObjectReplacements/Object 1839}

          Por la propiedad de resultar
          \includegraphics[width=0.411cm,height=0.467cm]{./ObjectReplacements/Object 1842}la
          suma de un elemento y su complementario en el axioma de
          complementarios.
        \item
          \includegraphics[width=1.05cm,height=0.467cm]{./ObjectReplacements/Object 1843}

          Por el axioma del elemento neutro de la multiplicación.
        \end{enumerate}
      \end{enumerate}
    \item
      Propiedades de simplificación (para retículos en general):

      \begin{enumerate}
      \def\labelenumiv{\arabic{enumiv}.}
      \item
        \includegraphics[width=4.173cm,height=0.577cm]{./ObjectReplacements/Object 12}

        Prueba:

        \begin{enumerate}
        \def\labelenumv{\arabic{enumv}.}
        \item
          \includegraphics[width=2cm,height=0.506cm]{./ObjectReplacements/Object 133}
        \item
          \includegraphics[width=4.48cm,height=0.506cm]{./ObjectReplacements/Object 41}

          Por ser
          el\includegraphics[width=0.411cm,height=0.467cm]{./ObjectReplacements/Object 86}el
          elemento neutro del producto (Axioma).
        \item
          \includegraphics[width=3.81cm,height=0.506cm]{./ObjectReplacements/Object 42}

          Por el axioma de distributividad del producto respecto de la
          suma por la izquierda.
        \item
          \includegraphics[width=2.967cm,height=0.467cm]{./ObjectReplacements/Object 84}

          Por ser
          el\includegraphics[width=0.411cm,height=0.467cm]{./ObjectReplacements/Object 87}elemento
          absorbente en la suma {[}Teorema (2.1){]}.
        \item
          \[\mspace{306mu}{= x}\]

          Por ser
          el\includegraphics[width=0.411cm,height=0.467cm]{./ObjectReplacements/Object 124}el
          elemento neutro del producto (Axioma).
        \end{enumerate}
      \item
        \includegraphics[width=4.173cm,height=0.577cm]{./ObjectReplacements/Object 13}

        \begin{enumerate}
        \def\labelenumv{\arabic{enumv}.}
        \item
          Prueba:
        \item
          \includegraphics[width=2cm,height=0.506cm]{./ObjectReplacements/Object 128}
        \item
          \includegraphics[width=4.648cm,height=0.506cm]{./ObjectReplacements/Object 130}

          Por ser
          el\includegraphics[width=0.423cm,height=0.467cm]{./ObjectReplacements/Object 131}el
          elemento neutro de la suma (Axioma).
        \item
          \includegraphics[width=3.822cm,height=0.506cm]{./ObjectReplacements/Object 132}

          Por el axioma de distributividad de la suma respecto del
          producto por la izquierda.
        \item
          \includegraphics[width=2.967cm,height=0.467cm]{./ObjectReplacements/Object 134}

          Por ser
          el\includegraphics[width=0.423cm,height=0.467cm]{./ObjectReplacements/Object 241}elemento
          absorbente en el producto {[}Teorema (2.2){]}.
        \item
          \includegraphics[width=2.267cm,height=0.467cm]{./ObjectReplacements/Object 242}

          Por ser
          el\includegraphics[width=0.423cm,height=0.467cm]{./ObjectReplacements/Object 127}el
          elemento neutro de la suma (Axioma).
        \end{enumerate}
      \end{enumerate}
    \item
      Una propiedad muy general (para retículos no sólo para álgebras de
      Boole):

      \begin{enumerate}
      \def\labelenumiv{\arabic{enumiv}.}
      \item
        \includegraphics[width=5.83cm,height=0.577cm]{./ObjectReplacements/Object 203}

        Prueba:

        \begin{enumerate}
        \def\labelenumv{\arabic{enumv}.}
        \item
          \begin{enumerate}
          \def\labelenumvi{\arabic{enumvi}.}
          \item
            \includegraphics[width=1.452cm,height=0.467cm]{./ObjectReplacements/Object 125}

            Hipótesis.
          \item
            \includegraphics[width=1.416cm,height=0.467cm]{./ObjectReplacements/Object 244}

            Verdad universal para la igualdad.
          \item
            \includegraphics[width=2.828cm,height=0.506cm]{./ObjectReplacements/Object 245}

            Como la suma es una aplicación, al sumar las expresiones
            izquierdas entre si, e idénticamente con las derechas, la
            igualdad se mantiene.
          \item
            \includegraphics[width=2.828cm,height=0.506cm]{./ObjectReplacements/Object 246}

            Propiedad conmutativa del producto (axioma).
          \item
            \includegraphics[width=1.589cm,height=0.467cm]{./ObjectReplacements/Object 247}

            Propiedad de simplificación del teorema anterior (11.2).
          \item
            \includegraphics[width=1.589cm,height=0.467cm]{./ObjectReplacements/Object 248}

            Simetría de la igualdad.
          \item
            \includegraphics[width=1.589cm,height=0.467cm]{./ObjectReplacements/Object 249}

            Propiedad conmutativa de la suma (axioma).
          \item
            \includegraphics[width=3.856cm,height=0.506cm]{./ObjectReplacements/Object 126}

            De (1) hemos derivado (7).
          \item
            \includegraphics[width=1.572cm,height=0.467cm]{./ObjectReplacements/Object 555}

            Hipótesis.
          \item
            \includegraphics[width=1.453cm,height=0.467cm]{./ObjectReplacements/Object 556}

            Verdad universal de la igualdad.
          \item
            \includegraphics[width=2.693cm,height=0.506cm]{./ObjectReplacements/Object 557}

            Como el producto es aplicación, el multiplicar los términos
            izquierdos de las igualdades (9) y (10) entre sí y los
            derechos de (9) y (10) idénticamente, se mantiene la
            igualdad.
          \item
            \includegraphics[width=1.455cm,height=0.467cm]{./ObjectReplacements/Object 558}

            Teorema de simplificación inmediatamente anterior (11.1).
          \item
            \includegraphics[width=1.455cm,height=0.467cm]{./ObjectReplacements/Object 559}

            Por la simetría de la igualdad.
          \item
            \includegraphics[width=3.856cm,height=0.506cm]{./ObjectReplacements/Object 243}

            De (9) hemos derivado (13).
          \item
            \includegraphics[width=8.675cm,height=1.09cm]{./ObjectReplacements/Object 628}

            (11) y (14) significan exactamente la definición del ``si y
            solo si''.
          \end{enumerate}
        \end{enumerate}
      \end{enumerate}
    \item
      Otra propiedad de simplificación (Shannon):

      \begin{enumerate}
      \def\labelenumiv{\arabic{enumiv}.}
      \item
        \includegraphics[width=5.022cm,height=0.577cm]{./ObjectReplacements/Object 15}

        Prueba:

        \begin{enumerate}
        \def\labelenumv{\arabic{enumv}.}
        \item
          \includegraphics[width=1.905cm,height=0.467cm]{./ObjectReplacements/Object 250}

          Por el axioma de elemento neutro de la suma.
        \item
          \includegraphics[width=2.686cm,height=0.506cm]{./ObjectReplacements/Object 554}

          Por el axioma de complementarios en su afirmación del producto
          de complementarios.
        \item
          \includegraphics[width=3.171cm,height=0.506cm]{./ObjectReplacements/Object 560}

          Por la distributividad de la suma respecto al producto por la
          izquierda.
        \end{enumerate}
      \item
        \includegraphics[width=4.867cm,height=0.577cm]{./ObjectReplacements/Object 16}

        Prueba:

        \begin{enumerate}
        \def\labelenumv{\arabic{enumv}.}
        \item
          \includegraphics[width=1.905cm,height=0.467cm]{./ObjectReplacements/Object 135}

          Por el axioma de elemento neutro de la suma.
        \item
          \includegraphics[width=2.686cm,height=0.506cm]{./ObjectReplacements/Object 136}

          Por el axioma de complementarios en su afirmación del producto
          de complementarios.
        \item
          \includegraphics[width=3.171cm,height=0.506cm]{./ObjectReplacements/Object 561}

          Por la distributividad de la suma respecto al producto por la
          izquierda.
        \end{enumerate}
      \end{enumerate}
    \item
      Otra propiedad de simplificación más:

      \begin{enumerate}
      \def\labelenumiv{\arabic{enumiv}.}
      \item
        Para la suma respecto del producto:
        \includegraphics[width=4.731cm,height=0.577cm]{./ObjectReplacements/Object 28}

        Prueba:

        \begin{enumerate}
        \def\labelenumv{\arabic{enumv}.}
        \item
          \includegraphics[width=4.618cm,height=0.506cm]{./ObjectReplacements/Object 137}

          Por el axioma de distributividad de la suma respecto al
          producto.
        \item
          \includegraphics[width=3.81cm,height=0.506cm]{./ObjectReplacements/Object 1591}

          Por la propiedad de sumar 1 los complementarios en el axioma
          de complementarios.
        \item
          \includegraphics[width=2.812cm,height=0.467cm]{./ObjectReplacements/Object 11}

          Por el axioma del elemento neutro de la suma.
        \end{enumerate}
      \item
        Para el producto respecto de la
        suma:\includegraphics[width=4.731cm,height=0.577cm]{./ObjectReplacements/Object 138}

        Prueba:

        \begin{enumerate}
        \def\labelenumv{\arabic{enumv}.}
        \item
          \includegraphics[width=4.463cm,height=0.506cm]{./ObjectReplacements/Object 1593}

          Por el axioma de distributividad del producto respecto de la
          suma.
        \item
          \includegraphics[width=3.822cm,height=0.506cm]{./ObjectReplacements/Object 1594}

          Por la propiedad de multiplicar
          \includegraphics[width=0.423cm,height=0.467cm]{./ObjectReplacements/Object 1595}
          los complementarios en el axioma de complementarios.
        \item
          \includegraphics[width=2.656cm,height=0.467cm]{./ObjectReplacements/Object 1592}

          Por el axioma del elemento neutro del producto.
        \end{enumerate}
      \end{enumerate}
    \item
      Algunos lemas técnicos pre-asociativos útiles:

      \begin{enumerate}
      \def\labelenumiv{\arabic{enumiv}.}
      \item
        Asociatividad cuando dos variables se repiten:

        \begin{enumerate}
        \def\labelenumv{\arabic{enumv}.}
        \item
          Para la
          suma:\includegraphics[width=6.565cm,height=0.577cm]{./ObjectReplacements/Object 118}

          Prueba:

          \begin{enumerate}
          \def\labelenumvi{\arabic{enumvi}.}
          \item
            \includegraphics[width=2.865cm,height=0.506cm]{./ObjectReplacements/Object 129}

            Ponemos una expresión izquierda de una igualdad.
          \item
            \includegraphics[width=5.526cm,height=0.506cm]{./ObjectReplacements/Object 1596}

            Por el axioma de conmutatividad del producto.
          \item
            \includegraphics[width=6.216cm,height=0.506cm]{./ObjectReplacements/Object 1853}

            Distributividad del producto respecto a la suma por la
            izquierda.
          \item
            \includegraphics[width=6.906cm,height=0.506cm]{./ObjectReplacements/Object 1758}

            Distributividad del producto respecto a la suma por la
            izquierda.
          \item
            \includegraphics[width=6.216cm,height=0.506cm]{./ObjectReplacements/Object 1759}

            Por el teorema de idempotencia del producto.
          \item
            \includegraphics[width=6.216cm,height=0.506cm]{./ObjectReplacements/Object 1854}

            Por axioma de conmutatividad de la suma.
          \item
            \includegraphics[width=6.216cm,height=0.506cm]{./ObjectReplacements/Object 1855}

            Por axioma de conmutatividad de la suma.
          \item
            \includegraphics[width=6.216cm,height=0.506cm]{./ObjectReplacements/Object 1858}

            Por axioma de conmutatividad de la suma.
          \item
            \includegraphics[width=6.216cm,height=0.506cm]{./ObjectReplacements/Object 1859}

            Por axioma de conmutatividad del producto.
          \item
            \includegraphics[width=4.701cm,height=0.506cm]{./ObjectReplacements/Object 1840}

            Por el teorema de simplificación de la suma respecto al
            producto.
          \item
            \includegraphics[width=4.701cm,height=0.506cm]{./ObjectReplacements/Object 1856}

            Por axioma de conmutatividad de la suma.
          \item
            \includegraphics[width=4.701cm,height=0.506cm]{./ObjectReplacements/Object 1857}

            Por axioma de conmutatividad del producto.
          \item
            \includegraphics[width=6.316cm,height=0.467cm]{./ObjectReplacements/Object 1841}

            Por el teorema de simplificación de la suma respecto al
            producto.
          \item
            \includegraphics[width=3.129cm,height=0.506cm]{./ObjectReplacements/Object 202}

            Fórmulas (1) y (13) por la transitividad de la igualdad.
          \item
            \includegraphics[width=2.829cm,height=0.506cm]{./ObjectReplacements/Object 564}

            Ponemos una expresión izquierda de una igualdad.
          \item
            \includegraphics[width=6.177cm,height=0.506cm]{./ObjectReplacements/Object 565}

            Por el axioma de distributividad del producto respecto de la
            suma.
          \item
            \includegraphics[width=6.177cm,height=0.506cm]{./ObjectReplacements/Object 1848}

            Por el axioma de conmutatividad de la suma.
          \item
            \includegraphics[width=4.664cm,height=0.506cm]{./ObjectReplacements/Object 1844}

            Teorema de simplificación del producto respecto de la suma.
          \item
            \includegraphics[width=6.339cm,height=0.467cm]{./ObjectReplacements/Object 1847}

            Teorema de simplificación del producto respecto de la suma.
          \item
            \includegraphics[width=3.06cm,height=0.506cm]{./ObjectReplacements/Object 10}

            Fórmulas (15) y (19) y transitividad de la igualdad.
          \item
            \includegraphics[width=1.385cm,height=0.467cm]{./ObjectReplacements/Object 567}
          \item
            \includegraphics[width=7.227cm,height=0.506cm]{./ObjectReplacements/Object 1860}

            Fórmulas (14) y (20) y simetría de la igualdad. De otro lado
            está que al ser la suma una aplicación, se puede sumar lado
            derecho de una igualdad con el lado derecho de otra y e
            igualmente sus lados izquierdos, y el resultado sigue siendo
            una igualdad.
          \item
            \includegraphics[width=7.227cm,height=0.506cm]{./ObjectReplacements/Object 1861}

            Axioma de conmutatividad del producto.
          \item
            \includegraphics[width=4.898cm,height=0.506cm]{./ObjectReplacements/Object 568}

            Por el axioma de distributividad por la izquierda del
            producto respecto de la suma.
          \item
            \includegraphics[width=6.521cm,height=0.563cm]{./ObjectReplacements/Object 1862}

            Por el teorema de idempotencia de la suma.
          \item
            \includegraphics[width=6.482cm,height=0.563cm]{./ObjectReplacements/Object 566}

            Por el axioma de conmutabilidad del producto.
          \item
            \includegraphics[width=6.482cm,height=0.563cm]{./ObjectReplacements/Object 1863}

            Por el axioma de conmutatividad de la suma.
          \item
            \includegraphics[width=4.898cm,height=0.506cm]{./ObjectReplacements/Object 569}

            Por el axioma de distributividad de la suma respecto al
            producto.
          \item
            \includegraphics[width=4.898cm,height=0.506cm]{./ObjectReplacements/Object 1846}

            Axioma de conmutatividad del producto.
          \item
            \includegraphics[width=4.898cm,height=0.506cm]{./ObjectReplacements/Object 1864}

            Axioma de conmutatividad de la suma.
          \item
            \includegraphics[width=3.39cm,height=0.506cm]{./ObjectReplacements/Object 1849}

            Teorema de simplificación del producto respecto de la suma.
          \item
            \includegraphics[width=7.922cm,height=0.506cm]{./ObjectReplacements/Object 570}

            Por el axioma de conmutatividad de la suma.
          \item
            \includegraphics[width=2.979cm,height=0.506cm]{./ObjectReplacements/Object 1845}

            De las fórmulas (21) y (32) y la transitividad de la
            igualdad.
          \item
            \includegraphics[width=1.863cm,height=0.519cm]{./ObjectReplacements/Object 1865}

            Teorema de idempotencia de la suma.
          \item
            \includegraphics[width=1.076cm,height=0.467cm]{./ObjectReplacements/Object 1866}

            Propiedad fundamental de la igualdad.
          \item
            \includegraphics[width=2.983cm,height=0.506cm]{./ObjectReplacements/Object 1867}

            Al ser la suma una aplicación y sumar los lados derechos de
            las igualdades y los lados izquierdos permanece la igualdad.
          \item
            \includegraphics[width=4.951cm,height=0.506cm]{./ObjectReplacements/Object 212}

            Concatenación de las fórmulas (36) y (33) por transitividad
            de la igualdad.
          \end{enumerate}
        \item
          Para el
          producto:\includegraphics[width=6.068cm,height=0.577cm]{./ObjectReplacements/Object 119}

          Prueba:

          \begin{enumerate}
          \def\labelenumvi{\arabic{enumvi}.}
          \item
            \includegraphics[width=2.709cm,height=0.506cm]{./ObjectReplacements/Object 88}

            Ponemos una expresión izquierda bien formada de una
            igualdad.
          \item
            \includegraphics[width=5.373cm,height=0.506cm]{./ObjectReplacements/Object 211}

            Por el axioma de conmutatividad de la suma.
          \item
            \includegraphics[width=6.216cm,height=0.506cm]{./ObjectReplacements/Object 213}

            Distributividad de la suma respecto al producto por la
            izquierda.
          \item
            \includegraphics[width=7.061cm,height=0.506cm]{./ObjectReplacements/Object 214}

            Distributividad de la suma respecto al producto por la
            izquierda.
          \item
            \includegraphics[width=6.216cm,height=0.506cm]{./ObjectReplacements/Object 215}

            Por el teorema de idempotencia de la suma.
          \item
            \includegraphics[width=6.216cm,height=0.506cm]{./ObjectReplacements/Object 571}

            Por axioma de conmutatividad del producto.
          \item
            \includegraphics[width=6.216cm,height=0.506cm]{./ObjectReplacements/Object 572}

            Por axioma de conmutatividad del producto.
          \item
            \includegraphics[width=6.216cm,height=0.506cm]{./ObjectReplacements/Object 573}

            Por axioma de conmutatividad del producto.
          \item
            \includegraphics[width=6.216cm,height=0.506cm]{./ObjectReplacements/Object 574}

            Por axioma de conmutatividad de la suma.
          \item
            \includegraphics[width=4.701cm,height=0.506cm]{./ObjectReplacements/Object 575}

            Por el teorema de simplificación del producto respecto a la
            suma.
          \item
            \includegraphics[width=4.701cm,height=0.506cm]{./ObjectReplacements/Object 576}

            Por axioma de conmutatividad del producto.
          \item
            \includegraphics[width=4.701cm,height=0.506cm]{./ObjectReplacements/Object 1850}

            Por axioma de conmutatividad de la suma.
          \item
            \includegraphics[width=6.316cm,height=0.467cm]{./ObjectReplacements/Object 1868}

            Por el teorema de simplificación del producto respecto a la
            suma.
          \item
            \includegraphics[width=2.976cm,height=0.506cm]{./ObjectReplacements/Object 1869}

            Fórmulas (1) y (13) por la transitividad de la igualdad.
          \item
            \includegraphics[width=2.676cm,height=0.506cm]{./ObjectReplacements/Object 1870}

            Ponemos una expresión izquierda de una igualdad.
          \item
            \includegraphics[width=6.177cm,height=0.506cm]{./ObjectReplacements/Object 1871}

            Por el axioma de distributividad de la suma respecto del
            producto.
          \item
            \includegraphics[width=6.177cm,height=0.506cm]{./ObjectReplacements/Object 1872}

            Por el axioma de conmutatividad del producto.
          \item
            \includegraphics[width=4.664cm,height=0.506cm]{./ObjectReplacements/Object 1873}

            Teorema de simplificación de la suma respecto del producto.
          \item
            \includegraphics[width=6.339cm,height=0.467cm]{./ObjectReplacements/Object 1874}

            Teorema de simplificación del producto respecto de la suma.
          \item
            \[{{({x \cdot {({x \cdot y})}})} + x} = x\]

            Fórmulas (15) y (19) y transitividad de la igualdad.
          \item
            \includegraphics[width=1.229cm,height=0.467cm]{./ObjectReplacements/Object 1876}

            Expresión bien formada en la parte izquierda de la igualdad.
          \item
            \includegraphics[width=6.765cm,height=0.506cm]{./ObjectReplacements/Object 1877}

            Fórmulas (14) y (20) y simetría de la igualdad. De otro lado
            está que al ser el producto una aplicación, se puede
            multiplicar lado derecho de una igualdad con el lado derecho
            de otra y e igualmente sus lados izquierdos, y el resultado
            sigue siendo una igualdad.
          \item
            \includegraphics[width=6.765cm,height=0.506cm]{./ObjectReplacements/Object 1878}

            Axioma de conmutatividad de la suma.
          \item
            \includegraphics[width=4.59cm,height=0.506cm]{./ObjectReplacements/Object 1879}

            Por el axioma de distributividad por la izquierda de la suma
            respecto del producto.
          \item
            \includegraphics[width=5.906cm,height=0.563cm]{./ObjectReplacements/Object 1880}

            Por el teorema de idempotencia del producto.
          \item
            \includegraphics[width=5.867cm,height=0.563cm]{./ObjectReplacements/Object 1881}

            Por el axioma de conmutabilidad de la suma.
          \item
            \includegraphics[width=5.867cm,height=0.563cm]{./ObjectReplacements/Object 1882}

            Por el axioma de conmutatividad del producto.
          \item
            \includegraphics[width=4.59cm,height=0.506cm]{./ObjectReplacements/Object 1883}

            Por el axioma de distributividad del producto respecto a la
            suma.
          \item
            \includegraphics[width=4.59cm,height=0.506cm]{./ObjectReplacements/Object 1884}

            Axioma de conmutatividad de la suma.
          \item
            \includegraphics[width=4.59cm,height=0.506cm]{./ObjectReplacements/Object 1885}

            Axioma de conmutatividad del producto.
          \item
            \includegraphics[width=3.082cm,height=0.506cm]{./ObjectReplacements/Object 1886}

            Teorema de simplificación de la suma respecto del producto.
          \item
            \includegraphics[width=7.615cm,height=0.506cm]{./ObjectReplacements/Object 1887}

            Por el axioma de conmutatividad del producto.
          \item
            \includegraphics[width=2.519cm,height=0.506cm]{./ObjectReplacements/Object 1888}

            De las fórmulas (21) y (32) y la transitividad de la
            igualdad.
          \item
            \includegraphics[width=1.707cm,height=0.519cm]{./ObjectReplacements/Object 1889}

            Teorema de idempotencia del producto.
          \item
            \includegraphics[width=1.076cm,height=0.467cm]{./ObjectReplacements/Object 1890}

            Propiedad fundamental de la igualdad.
          \item
            \includegraphics[width=2.521cm,height=0.506cm]{./ObjectReplacements/Object 1891}

            Al ser el producto una aplicación y multiplicar los lados
            derechos de las igualdades y los lados izquierdos permanece
            la igualdad.
          \item
            \includegraphics[width=4.182cm,height=0.506cm]{./ObjectReplacements/Object 1892}

            Concatenación de las fórmulas (36) y (33) por transitividad
            de la igualdad.
          \end{enumerate}
        \end{enumerate}
      \item
        Asociatividad cuando una de las tres variables es la
        complementaria de otra:

        \begin{enumerate}
        \def\labelenumv{\arabic{enumv}.}
        \item
          Para la
          suma:\includegraphics[width=4.322cm,height=0.506cm]{./ObjectReplacements/Object 120}

          Prueba:

          \begin{enumerate}
          \def\labelenumvi{\arabic{enumvi}.}
          \item
            \includegraphics[width=2.831cm,height=0.506cm]{./ObjectReplacements/Object 577}

            Expresión derecha de la igualdad.
          \item
            \includegraphics[width=5.592cm,height=0.506cm]{./ObjectReplacements/Object 1923}

            Axioma de conmutatividad del producto.
          \item
            \includegraphics[width=6.262cm,height=0.506cm]{./ObjectReplacements/Object 1893}

            Axioma de distributividad del producto del producto respecto
            de la suma.
          \item
            \includegraphics[width=5.572cm,height=0.506cm]{./ObjectReplacements/Object 1914}

            Propiedad de
            multiplicar\includegraphics[width=0.423cm,height=0.467cm]{./ObjectReplacements/Object 1925}los
            complementarios por el axioma de complementarios.
          \item
            \includegraphics[width=4.766cm,height=0.506cm]{./ObjectReplacements/Object 1915}

            Axioma de elemento neutro de la suma.
          \item
            \includegraphics[width=4.766cm,height=0.506cm]{./ObjectReplacements/Object 1924}

            Axioma de conmutatividad del producto.
          \item
            \includegraphics[width=4.766cm,height=0.506cm]{./ObjectReplacements/Object 1916}

            Axioma de conmutatividad de la suma.
          \item
            \includegraphics[width=6.634cm,height=0.467cm]{./ObjectReplacements/Object 1894}

            Teorema de simplificación del producto respecto de la suma.
          \item
            \includegraphics[width=2.831cm,height=0.506cm]{./ObjectReplacements/Object 1895}

            Transitividad de la igualdad desde la fórmula (1) a la (8).
          \item
            \includegraphics[width=2.833cm,height=0.506cm]{./ObjectReplacements/Object 578}

            Expresión derecha de la igualdad.
          \item
            \includegraphics[width=5.48cm,height=0.506cm]{./ObjectReplacements/Object 1917}

            Axioma de conmutatividad de la suma.
          \item
            \includegraphics[width=6.623cm,height=0.467cm]{./ObjectReplacements/Object 1918}

            Teorema de simplificación del producto respecto de la suma.
          \item
            \includegraphics[width=3.106cm,height=0.506cm]{./ObjectReplacements/Object 1898}

            Transitividad de la igualdad des la fórmula (10) hasta la
            (12).
          \item
            \includegraphics[width=0.764cm,height=0.467cm]{./ObjectReplacements/Object 580}

            Expresión bien formada como lado izquierdo de la igualdad.
          \item
            \includegraphics[width=3.046cm,height=0.467cm]{./ObjectReplacements/Object 1930}

            Propiedad de sumar
            \includegraphics[width=0.411cm,height=0.467cm]{./ObjectReplacements/Object 1926}los
            complementarios en el axioma de existencia de
            complementarios.
          \item
            \includegraphics[width=7.504cm,height=0.631cm]{./ObjectReplacements/Object 579}

            Sustitución de los lados derechos de las igualdades (9) y
            (13) por sus lados izquierdos.
          \item
            \includegraphics[width=5.369cm,height=0.506cm]{./ObjectReplacements/Object 581}

            Axioma de distributividad izquierda del producto respecto de
            la suma.
          \item
            \includegraphics[width=4.551cm,height=0.506cm]{./ObjectReplacements/Object 1899}

            Propiedad de
            sumar\includegraphics[width=0.411cm,height=0.467cm]{./ObjectReplacements/Object 1927}los
            complementarios del axioma de existencia de complementario.
          \item
            \includegraphics[width=7.939cm,height=0.506cm]{./ObjectReplacements/Object 1900}

            Axioma de elemento neutro del producto.
          \item
            \includegraphics[width=2.395cm,height=0.506cm]{./ObjectReplacements/Object 1896}

            Transitividad de la igualdad desde la fórmula (14) hasta la
            (19).
          \item
            \includegraphics[width=0.764cm,height=0.467cm]{./ObjectReplacements/Object 582}

            Expresión bien formada en lado izquierdo de la igualdad.
          \item
            \includegraphics[width=1.967cm,height=0.467cm]{./ObjectReplacements/Object 1929}

            Teorema del elemento absorbente de la suma.
          \item
            \includegraphics[width=7.927cm,height=0.506cm]{./ObjectReplacements/Object 583}

            Propiedad de
            sumar\includegraphics[width=0.411cm,height=0.467cm]{./ObjectReplacements/Object 1928}los
            complementarios en el axioma de existencias de
            complementarios.
          \item
            \includegraphics[width=2.365cm,height=0.506cm]{./ObjectReplacements/Object 584}

            Transitividad de la igualdad desde la fórmula (21) hasta la
            (23).
          \item
            \includegraphics[width=4.355cm,height=0.506cm]{./ObjectReplacements/Object 1902}

            Transitividad de la igualdad y simetría de la misma en las
            fórmulas (20) y (24).
          \end{enumerate}
        \item
          Para el
          producto:\includegraphics[width=3.717cm,height=0.506cm]{./ObjectReplacements/Object 121}

          Prueba:

          \begin{enumerate}
          \def\labelenumvi{\arabic{enumvi}.}
          \item
            \includegraphics[width=2.678cm,height=0.506cm]{./ObjectReplacements/Object 585}

            Expresión bien formada como lado izquierdo de la igualdad.
          \item
            \includegraphics[width=5.073cm,height=0.506cm]{./ObjectReplacements/Object 1931}

            Axioma de conmutatividad de la suma.
          \item
            \includegraphics[width=5.837cm,height=0.575cm]{./ObjectReplacements/Object 586}

            Axioma de distributividad izquierda de la suma respecto al
            producto.
          \item
            \includegraphics[width=5.837cm,height=0.575cm]{./ObjectReplacements/Object 1912}

            Axioma de conmutatividad del producto.
          \item
            \includegraphics[width=5.027cm,height=0.575cm]{./ObjectReplacements/Object 1913}

            Propiedad de
            sumar\includegraphics[width=0.411cm,height=0.467cm]{./ObjectReplacements/Object 1933}los
            complementarios del axioma de existencia de complementarios.
          \item
            \includegraphics[width=4.424cm,height=0.519cm]{./ObjectReplacements/Object 1932}

            Axioma de elemento neutro del producto.
          \item
            \includegraphics[width=4.424cm,height=0.519cm]{./ObjectReplacements/Object 1901}

            Axioma de conmutatividad de la suma.
          \item
            \includegraphics[width=4.424cm,height=0.519cm]{./ObjectReplacements/Object 1919}

            Axioma de conmutatividad del producto.
          \item
            \includegraphics[width=4.955cm,height=0.467cm]{./ObjectReplacements/Object 587}

            Teorema de simplificación de la suma respecto al producto.
          \item
            \includegraphics[width=2.951cm,height=0.506cm]{./ObjectReplacements/Object 588}

            Transitividad de la igualdad desde la fórmula (1) a la (9).
          \item
            \includegraphics[width=2.586cm,height=0.575cm]{./ObjectReplacements/Object 589}

            Expresión bien formada como lado izquierdo de la igualdad.
          \item
            \includegraphics[width=5.59cm,height=0.575cm]{./ObjectReplacements/Object 1920}

            Axioma de distributividad izquierda de la suma respecto al
            producto.
          \item
            \includegraphics[width=4.822cm,height=0.575cm]{./ObjectReplacements/Object 590}

            Teorema de idempotencia de la suma.
          \item
            \includegraphics[width=4.822cm,height=0.575cm]{./ObjectReplacements/Object 1921}

            Axioma de conmutatividad del producto.
          \item
            \includegraphics[width=5.877cm,height=0.467cm]{./ObjectReplacements/Object 591}

            Teorema de simplificación del producto respecto de la suma.
          \item
            \includegraphics[width=2.861cm,height=0.575cm]{./ObjectReplacements/Object 1903}

            Transitividad de la igualdad desde la fórmula (11) hasta la
            (15).
          \item
            \includegraphics[width=0.774cm,height=0.467cm]{./ObjectReplacements/Object 1904}

            Expresión bien formada como lado izquierdo de la igualdad.
          \item
            \includegraphics[width=2.829cm,height=0.467cm]{./ObjectReplacements/Object 1934}

            Propiedad de multiplicar 0 los complementarios del axioma de
            existencia de complementarios.
          \item
            \includegraphics[width=6.959cm,height=0.631cm]{./ObjectReplacements/Object 1905}

            Sustitución de los lados izquierdos de las fórmulas (10) y
            (16) por sus lados derechos.
          \item
            \includegraphics[width=4.978cm,height=0.506cm]{./ObjectReplacements/Object 1906}

            Axioma de distributividad izquierda de la suma respecto al
            producto.
          \item
            \includegraphics[width=4.327cm,height=0.506cm]{./ObjectReplacements/Object 1907}

            Propiedad de
            multiplicar\includegraphics[width=0.423cm,height=0.467cm]{./ObjectReplacements/Object 1936}los
            complementarios del axioma de existencia de complementarios.
          \item
            \includegraphics[width=6.175cm,height=0.506cm]{./ObjectReplacements/Object 1908}

            Axioma de elemento neutro de la suma.
          \item
            \includegraphics[width=2.101cm,height=0.506cm]{./ObjectReplacements/Object 1897}

            Transitividad de la igualdad desde la fórmula (17) hasta la
            (22).
          \item
            \includegraphics[width=0.774cm,height=0.467cm]{./ObjectReplacements/Object 1909}

            Expresión bien formada en lado izquierdo de la igualdad.
          \item
            \includegraphics[width=2.746cm,height=0.467cm]{./ObjectReplacements/Object 1935}

            Teorema de elemento absorbente de la multiplicación.
          \item
            \includegraphics[width=3.436cm,height=0.506cm]{./ObjectReplacements/Object 1910}

            Propiedad de
            multiplicar\includegraphics[width=0.423cm,height=0.467cm]{./ObjectReplacements/Object 1937}los
            complementarios del axioma de existencia de complementarios.
          \item
            \includegraphics[width=4.763cm,height=0.506cm]{./ObjectReplacements/Object 1922}

            Axioma de conmutatividad del producto.
          \item
            \includegraphics[width=2.069cm,height=0.506cm]{./ObjectReplacements/Object 1911}

            Transitividad de la igualdad desde la fórmula (24) a la
            (27).
          \item
            \includegraphics[width=3.752cm,height=0.506cm]{./ObjectReplacements/Object 139}

            Transitividad y simetría de la igualdad con las fórmulas
            (23) y (28).
          \end{enumerate}
        \end{enumerate}
      \end{enumerate}
    \item
      Leyes de Morgan:

      \begin{enumerate}
      \def\labelenumiv{\arabic{enumiv}.}
      \item
        De la suma en
        producto:\includegraphics[width=3.918cm,height=0.563cm]{./ObjectReplacements/Object 23}

        Prueba:

        Informal:

        \includegraphics[width=11.723cm,height=0.506cm]{./ObjectReplacements/Object 140}

        \includegraphics[width=11.998cm,height=0.506cm]{./ObjectReplacements/Object 141}

        Formal:

        \begin{enumerate}
        \def\labelenumv{\arabic{enumv}.}
        \item
          \includegraphics[width=2.713cm,height=0.506cm]{./ObjectReplacements/Object 1938}

          Expresión bien formada en lado izquierdo de la igualdad.
        \item
          \includegraphics[width=5.276cm,height=0.506cm]{./ObjectReplacements/Object 1955}

          Axioma de conmutatividad del producto.
        \item
          \includegraphics[width=6.632cm,height=0.506cm]{./ObjectReplacements/Object 1956}

          Axioma de distributividad izquierda del producto respecto de
          la suma.
        \item
          \includegraphics[width=6.632cm,height=0.506cm]{./ObjectReplacements/Object 1939}

          Axioma de conmutatividad del producto.
        \item
          \includegraphics[width=6.63cm,height=0.506cm]{./ObjectReplacements/Object 1940}

          Lema de asociatividad del producto cuando uno de los tres
          operandos es negado de otro.
        \item
          \includegraphics[width=5.96cm,height=0.506cm]{./ObjectReplacements/Object 1976}

          Propiedad de
          multiplicar\includegraphics[width=0.423cm,height=0.467cm]{./ObjectReplacements/Object 1942}los
          complementarios del axioma de existencia de complementarios.
        \item
          \includegraphics[width=5.265cm,height=0.506cm]{./ObjectReplacements/Object 1977}

          Teorema de elemento absorbente del producto.
        \item
          \includegraphics[width=4.461cm,height=0.506cm]{./ObjectReplacements/Object 1941}

          Axioma del elemento neutro de la suma.
        \item
          \includegraphics[width=4.463cm,height=0.506cm]{./ObjectReplacements/Object 1978}

          Teorema de la doble negación.
        \item
          \includegraphics[width=6.865cm,height=0.467cm]{./ObjectReplacements/Object 1943}

          Lema de asociatividad del producto cuando uno de los tres
          operandos es negado de otro.
        \item
          \includegraphics[width=2.976cm,height=0.506cm]{./ObjectReplacements/Object 1944}

          Transitividad de la igualdad desde la fórmula (1) hasta la
          (10).
        \item
          \includegraphics[width=2.866cm,height=0.506cm]{./ObjectReplacements/Object 1945}

          Expresión bien formada en lado izquierdo de la igualdad.
        \item
          \includegraphics[width=7.093cm,height=0.506cm]{./ObjectReplacements/Object 1946}

          Axioma de distributividad izquierda de la suma respecto al
          producto.
        \item
          \includegraphics[width=7.093cm,height=0.506cm]{./ObjectReplacements/Object 1947}

          Axioma de conmutatividad de la suma.
        \item
          \includegraphics[width=7.093cm,height=0.506cm]{./ObjectReplacements/Object 1979}

          Teorema de la doble negación.
        \item
          \includegraphics[width=7.093cm,height=0.506cm]{./ObjectReplacements/Object 1952}

          Axioma de asociatividad del producto cuando uno de los
          elementos es negado de otro.
        \item
          \includegraphics[width=5.352cm,height=0.506cm]{./ObjectReplacements/Object 1953}

          Propiedad de sumar 1 los complementarios del axioma de
          existencia de complementarios.
        \item
          \includegraphics[width=5.352cm,height=0.506cm]{./ObjectReplacements/Object 1948}

          Axioma de conmutatividad de la suma.
        \item
          \includegraphics[width=3.69cm,height=0.467cm]{./ObjectReplacements/Object 1949}

          Teorema del elemento absorbente de la suma.
        \item
          \includegraphics[width=6.854cm,height=0.467cm]{./ObjectReplacements/Object 1950}

          Axioma del elemento neutro del producto.
        \item
          \includegraphics[width=3.119cm,height=0.506cm]{./ObjectReplacements/Object 1951}

          Transitividad de la igualdad desde la fórmula (12) hasta la
          (21).
        \item
          \includegraphics[width=6.653cm,height=0.559cm]{./ObjectReplacements/Object 1954}

          Fórmulas (11) y (21).
        \item
          \includegraphics[width=2.598cm,height=0.63cm]{./ObjectReplacements/Object 1981}

          Implicado por el axioma de existencia de complementarios.
        \item
          \includegraphics[width=2.619cm,height=0.563cm]{./ObjectReplacements/Object 1980}

          Por el teorema de unicidad del complementario.
        \end{enumerate}
      \item
        Del producto en
        suma:\includegraphics[width=3.918cm,height=0.563cm]{./ObjectReplacements/Object 24}

        Prueba:

        Informal:

        \includegraphics[width=12.614cm,height=0.506cm]{./ObjectReplacements/Object 592}

        \includegraphics[width=12.185cm,height=0.506cm]{./ObjectReplacements/Object 142}

        Formal:

        \begin{enumerate}
        \def\labelenumv{\arabic{enumv}.}
        \item
          \includegraphics[width=2.713cm,height=0.506cm]{./ObjectReplacements/Object 1958}

          Expresión bien formada como parte izquierda de la igualdad.
        \item
          \includegraphics[width=6.63cm,height=0.506cm]{./ObjectReplacements/Object 1959}

          Axioma de distribución del producto respecto de la suma.
        \item
          \includegraphics[width=5.232cm,height=0.506cm]{./ObjectReplacements/Object 1982}

          Lema de asociatividad del producto cuando uno de los factores
          es complementario de otro.
        \item
          \includegraphics[width=4.427cm,height=0.506cm]{./ObjectReplacements/Object 1960}

          Axioma del elemento neutro de la suma.
        \item
          \includegraphics[width=4.427cm,height=0.506cm]{./ObjectReplacements/Object 1961}

          Axioma de conmutatividad del producto.
        \item
          \includegraphics[width=6.865cm,height=0.467cm]{./ObjectReplacements/Object 1965}

          Axioma de asociatividad del producto cuando uno de los
          operandos es complementario de otro.
        \item
          \includegraphics[width=2.976cm,height=0.506cm]{./ObjectReplacements/Object 1966}

          Transitividad de la igualdad desde la fórmula (1) a la (6).
        \item
          \includegraphics[width=2.866cm,height=0.506cm]{./ObjectReplacements/Object 1962}

          Expresión de lado izquierdo de la igualdad bien formada.
        \item
          \includegraphics[width=5.514cm,height=0.506cm]{./ObjectReplacements/Object 1968}

          Axioma de conmutatividad de la suma.
        \item
          \includegraphics[width=7.177cm,height=0.506cm]{./ObjectReplacements/Object 1963}

          Axioma de distribución izquierda de la suma respecto del
          producto.
        \item
          \includegraphics[width=7.177cm,height=0.506cm]{./ObjectReplacements/Object 1964}

          Axioma de conmutatividad de la suma.
        \item
          \includegraphics[width=5.493cm,height=0.506cm]{./ObjectReplacements/Object 1967}

          Lema asociativo de la suma cuando un operando es
          complementario de otro.
        \item
          \includegraphics[width=4.854cm,height=0.506cm]{./ObjectReplacements/Object 1969}

          Axioma de elemento neutro del producto.
        \item
          \[\mspace{486mu}{= {{({\overline{x} + \overline{y}})} + \overline{\overline{y}}} =}\]

          Teorema de la doble negación.
        \item
          \includegraphics[width=6.854cm,height=0.467cm]{./ObjectReplacements/Object 1975}

          Lema de asociatividad de la suma cuando un operando es
          complementario de otro.
        \item
          \includegraphics[width=2.965cm,height=0.506cm]{./ObjectReplacements/Object 1957}

          Transitividad de la igualdad desde la fórmula (8) hasta la
          (15).
        \item
          \includegraphics[width=6.653cm,height=0.559cm]{./ObjectReplacements/Object 1970}

          Fórmulas (7) y (16).
        \item
          \includegraphics[width=2.598cm,height=0.63cm]{./ObjectReplacements/Object 1971}

          Derivación por el axioma de existencia de complementarios.
        \item
          \[{{({\overline{x} \cdot \overline{y}})} = \overline{({x+y})}}{}\]

          Por el teorema de unicidad del complementario.
        \end{enumerate}
      \end{enumerate}
    \item
      Transformación de la suma y el producto mediante Morgan y doble
      complemento:

      \begin{enumerate}
      \def\labelenumiv{\arabic{enumiv}.}
      \item
        De la suma en
        producto:\includegraphics[width=3.918cm,height=0.563cm]{./ObjectReplacements/Object 27}

        Prueba:

        \begin{enumerate}
        \def\labelenumv{\arabic{enumv}.}
        \item
          \includegraphics[width=1.341cm,height=0.467cm]{./ObjectReplacements/Object 1974}

          Expresión bien formada en la parte izquierda de la igualdad.
        \item
          \includegraphics[width=2.738cm,height=0.573cm]{./ObjectReplacements/Object 1983}

          Por el teorema de la doble negación es identidad.
        \item
          \includegraphics[width=3.369cm,height=0.573cm]{./ObjectReplacements/Object 1984}

          Por el teorema de Morgan sobre la negación de una suma.
        \item
          \includegraphics[width=2.224cm,height=0.573cm]{./ObjectReplacements/Object 143}

          Por la transitividad de la igualdad desde la fórmula (1) a la
          (3).
        \end{enumerate}
      \item
        Del producto en
        suma:\includegraphics[width=3.918cm,height=0.563cm]{./ObjectReplacements/Object 26}

        Prueba:

        \begin{enumerate}
        \def\labelenumv{\arabic{enumv}.}
        \item
          \includegraphics[width=1.187cm,height=0.467cm]{./ObjectReplacements/Object 1985}

          Expresión bien formada en la parte izquierda de la igualdad.
        \item
          \includegraphics[width=2.582cm,height=0.573cm]{./ObjectReplacements/Object 1986}

          Por el teorema de doble negación es identidad.
        \item
          \includegraphics[width=3.522cm,height=0.573cm]{./ObjectReplacements/Object 1987}

          Por el teorema de Morgan de la negación del producto.
        \item
          \includegraphics[width=2.224cm,height=0.573cm]{./ObjectReplacements/Object 144}

          Transitividad de la igualdad desde la fórmula (1) a la (3).
        \end{enumerate}
      \end{enumerate}
    \item
      Asociatividad de las operaciones binarias:

      \begin{enumerate}
      \def\labelenumiv{\arabic{enumiv}.}
      \item
        De la suma:
        \includegraphics[width=6.061cm,height=0.577cm]{./ObjectReplacements/Object 22}

        Prueba:

        \begin{enumerate}
        \def\labelenumv{\arabic{enumv}.}
        \item
          \includegraphics[width=2.658cm,height=0.506cm]{./ObjectReplacements/Object 157}

          Definición dónde la parte derecha de la igualdad está bien
          formada.
        \item
          \includegraphics[width=2.658cm,height=0.506cm]{./ObjectReplacements/Object 158}

          Definición dónde la parte derecha de la igualdad está bien
          formada.
        \item
          \includegraphics[width=0.778cm,height=0.467cm]{./ObjectReplacements/Object 205}

          Expresión bien formada en la parte izquierda de la igualdad.
        \item
          \includegraphics[width=3.09cm,height=0.63cm]{./ObjectReplacements/Object 1991}

          Fórmula derivada de (2) por aplicación de una operación unaria
          unívoca sobre ambos lados de la igualdad.
        \item
          \includegraphics[width=2.949cm,height=0.637cm]{./ObjectReplacements/Object 1988}

          Teorema de Morgan de la negación de la suma.
        \item
          \includegraphics[width=4.45cm,height=0.575cm]{./ObjectReplacements/Object 1989}

          Teorema de Morgan de la negación de la suma.
        \item
          \includegraphics[width=2.281cm,height=0.575cm]{./ObjectReplacements/Object 1990}

          Transitividad de la igualdad desde la fórmula (3) a la (6).
        \item
          \includegraphics[width=1.293cm,height=0.467cm]{./ObjectReplacements/Object 206}

          Expresión bien formada en el lado izquierdo de la igualdad.
        \item
          \includegraphics[width=3.974cm,height=0.575cm]{./ObjectReplacements/Object 1992}

          Sustitución de un término de la fórmula (8) por el lado
          derecho de la igualdad (7).
        \item
          \includegraphics[width=4.246cm,height=0.575cm]{./ObjectReplacements/Object 1993}

          Axioma de distribución del producto respecto de la suma.
        \item
          \includegraphics[width=5.368cm,height=0.575cm]{./ObjectReplacements/Object 1994}

          Axioma de distribución del producto respecto de la suma.
        \item
          \includegraphics[width=10.409cm,height=0.506cm]{./ObjectReplacements/Object 209}

          Sustitución de un término de la fórmula (11) por el lado
          derecho de la igualdad definición (1).
        \item
          \includegraphics[width=7.962cm,height=0.506cm]{./ObjectReplacements/Object 504}

          Lema de asociatividad de la suma cuando un término es
          complementario de otro.
        \item
          \includegraphics[width=7.324cm,height=0.506cm]{./ObjectReplacements/Object 505}

          Axioma de elemento neutro del producto.
        \item
          \includegraphics[width=4.874cm,height=0.506cm]{./ObjectReplacements/Object 506}

          Axioma de distribución del producto respecto de la suma.
        \item
          \includegraphics[width=5.027cm,height=0.563cm]{./ObjectReplacements/Object 507}

          Teorema de Morgan de la negación de la suma.
        \item
          \includegraphics[width=5.027cm,height=0.563cm]{./ObjectReplacements/Object 208}

          Axioma de conmutatividad de la suma.
        \item
          \includegraphics[width=9.206cm,height=0.467cm]{./ObjectReplacements/Object 509}

          Lema de asociatividad de la suma cuando un término es
          complemento de otro.
        \item
          \includegraphics[width=1.519cm,height=0.467cm]{./ObjectReplacements/Object 2004}

          Transitividad de la igualdad desde la fórmula (8) a la (18).
        \item
          \includegraphics[width=1.139cm,height=0.467cm]{./ObjectReplacements/Object 510}

          Expresión bien formada en el lado izquierdo de la igualdad.
        \item
          \includegraphics[width=3.821cm,height=0.575cm]{./ObjectReplacements/Object 508}

          Sustitución de un término de la fórmula (21) por el lado
          derecho de la igualdad (7).
        \item
          \includegraphics[width=5.018cm,height=0.575cm]{./ObjectReplacements/Object 511}

          Sustitución de un término de la fórmula (21) por el lado
          derecho de la igualdad definición (1).
        \item
          \includegraphics[width=7.177cm,height=0.506cm]{./ObjectReplacements/Object 512}

          Axioma de distribución del producto respecto de la suma.
        \item
          \includegraphics[width=7.177cm,height=0.506cm]{./ObjectReplacements/Object 1995}

          Axioma de conmutatividad del producto.
        \item
          \[\mspace{144mu}{= {{({x \cdot {({\overline{x} \cdot {({\overline{y} \cdot \overline{z}})}})}})} + {({{({\overline{x} \cdot \overline{({y+z})}})} \cdot {({y + z})}})}} =}\]

          Teorema de Morgan de la negación de la suma.
        \item
          \includegraphics[width=7.624cm,height=0.617cm]{./ObjectReplacements/Object 1997}

          Teorema de doble negación es identidad.
        \item
          \includegraphics[width=4.355cm,height=0.506cm]{./ObjectReplacements/Object 1998}

          Lema de asociatividad del producto cuando un término es
          complemento de otro.
        \item
          \includegraphics[width=3.551cm,height=0.506cm]{./ObjectReplacements/Object 1999}

          Axioma de elemento neutro de la suma.
        \item
          \includegraphics[width=3.551cm,height=0.506cm]{./ObjectReplacements/Object 2000}

          Axioma de conmutatividad del producto.
        \item
          \includegraphics[width=3.551cm,height=0.506cm]{./ObjectReplacements/Object 2001}

          Axioma de conmutatividad del producto.
        \item
          \includegraphics[width=3.844cm,height=0.563cm]{./ObjectReplacements/Object 2002}

          Teorema de doble negación es identidad.
        \item
          \includegraphics[width=10.225cm,height=0.467cm]{./ObjectReplacements/Object 2003}

          Lema de asociatividad del producto cuando un término es
          complemento de otro.
        \item
          \includegraphics[width=1.378cm,height=0.467cm]{./ObjectReplacements/Object 2005}

          Transitividad del la igualdad desde la fórmula (20) a la (32).
        \item
          \includegraphics[width=0.982cm,height=0.467cm]{./ObjectReplacements/Object 519}

          De (33) y (19) por el axioma de existencia de elemento
          complementario.
        \item
          \includegraphics[width=1.002cm,height=0.467cm]{./ObjectReplacements/Object 513}

          Teorema de unicidad del complementario.
        \item
          \includegraphics[width=1.002cm,height=0.467cm]{./ObjectReplacements/Object 514}

          Simetría de la igualdad.
        \item
          \includegraphics[width=1.002cm,height=0.517cm]{./ObjectReplacements/Object 515}

          Por el teorema de unicidad del complementario la igualdad
          permanece.
        \item
          \includegraphics[width=1.018cm,height=0.467cm]{./ObjectReplacements/Object 2006}

          Teorema doble negación es identidad.
        \item
          \includegraphics[width=3.651cm,height=0.519cm]{./ObjectReplacements/Object 2007}

          Sustitución en (38) de
          \includegraphics[width=0.441cm,height=0.467cm]{./ObjectReplacements/Object 516}por
          el lado derecho de (1) y de
          \includegraphics[width=0.425cm,height=0.467cm]{./ObjectReplacements/Object 517}por
          el lado izquierdo de (2).
        \end{enumerate}
      \item
        De la operación producto:
        \includegraphics[width=5.445cm,height=0.577cm]{./ObjectReplacements/Object 520}

        Prueba:

        \begin{enumerate}
        \def\labelenumv{\arabic{enumv}.}
        \item
          \includegraphics[width=2.351cm,height=0.506cm]{./ObjectReplacements/Object 528}

          Definición dónde la parte derecha de la igualdad está bien
          formada.
        \item
          \includegraphics[width=2.351cm,height=0.506cm]{./ObjectReplacements/Object 529}

          Definición dónde la parte derecha de la igualdad está bien
          formada.
        \item
          \includegraphics[width=0.778cm,height=0.467cm]{./ObjectReplacements/Object 530}

          Expresión bien formada en la parte izquierda de la igualdad.
        \item
          \includegraphics[width=2.616cm,height=0.63cm]{./ObjectReplacements/Object 531}

          Fórmula derivada de (2) por aplicación de una operación unaria
          unívoca sobre ambos lados de la igualdad.
        \item
          \includegraphics[width=2.782cm,height=0.637cm]{./ObjectReplacements/Object 532}

          Teorema de Morgan de la negación del producto.
        \item
          \includegraphics[width=4.759cm,height=0.575cm]{./ObjectReplacements/Object 533}

          Teorema de Morgan de la negación del producto.
        \item
          \includegraphics[width=2.589cm,height=0.575cm]{./ObjectReplacements/Object 534}

          Transitividad de la igualdad desde la fórmula (3) a la (6).
        \item
          \includegraphics[width=1.139cm,height=0.467cm]{./ObjectReplacements/Object 535}

          Expresión bien formada en el lado izquierdo de la igualdad.
        \item
          \includegraphics[width=3.455cm,height=0.575cm]{./ObjectReplacements/Object 536}

          Sustitución de un término de la fórmula (8) por el lado
          derecho de la igualdad (7).
        \item
          \includegraphics[width=4.076cm,height=0.575cm]{./ObjectReplacements/Object 537}

          Axioma de distribución de la suma respecto del producto.
        \item
          \includegraphics[width=5.046cm,height=0.575cm]{./ObjectReplacements/Object 538}

          Axioma de distribución de la suma respecto del producto.
        \item
          \includegraphics[width=9.163cm,height=0.506cm]{./ObjectReplacements/Object 539}

          Sustitución de un término de la fórmula (11) por el lado
          derecho de la igualdad definición (1).
        \item
          \includegraphics[width=7.19cm,height=0.506cm]{./ObjectReplacements/Object 540}

          Lema de asociatividad del producto cuando un término es
          complementario de otro.
        \item
          \includegraphics[width=6.385cm,height=0.506cm]{./ObjectReplacements/Object 541}

          Axioma de elemento neutro de la suma.
        \item
          \includegraphics[width=4.397cm,height=0.506cm]{./ObjectReplacements/Object 542}

          Axioma de distribución de la suma respecto del producto.
        \item
          \includegraphics[width=4.242cm,height=0.563cm]{./ObjectReplacements/Object 543}

          Teorema de Morgan de la negación del producto.
        \item
          \includegraphics[width=4.242cm,height=0.563cm]{./ObjectReplacements/Object 2008}

          Axioma de conmutatividad del producto.
        \item
          \includegraphics[width=2.104cm,height=0.467cm]{./ObjectReplacements/Object 2009}

          Lema de asociatividad del producto cuando un término es
          complemento de otro.
        \item
          \includegraphics[width=1.378cm,height=0.467cm]{./ObjectReplacements/Object 2010}

          Transitividad de la igualdad desde la fórmula (8) a la (18).
        \item
          \includegraphics[width=1.293cm,height=0.467cm]{./ObjectReplacements/Object 2011}

          Expresión bien formada en el lado izquierdo de la igualdad.
        \item
          \includegraphics[width=3.149cm,height=0.575cm]{./ObjectReplacements/Object 2012}

          Sustitución de un término de la fórmula (21) por el lado
          derecho de la igualdad (7).
        \item
          \includegraphics[width=5.173cm,height=0.575cm]{./ObjectReplacements/Object 2013}

          Sustitución de un término de la fórmula (21) por el lado
          derecho de la igualdad definición (1).
        \item
          \includegraphics[width=7.795cm,height=0.506cm]{./ObjectReplacements/Object 2014}

          Axioma de distribución de la suma respecto del producto.
        \item
          \includegraphics[width=7.795cm,height=0.506cm]{./ObjectReplacements/Object 2015}

          Axioma de conmutatividad de la suma.
        \item
          \includegraphics[width=7.639cm,height=0.563cm]{./ObjectReplacements/Object 2016}

          Teorema de Morgan de la negación del producto.
        \item
          \includegraphics[width=7.932cm,height=0.617cm]{./ObjectReplacements/Object 2017}

          Teorema de doble negación es identidad.
        \item
          \includegraphics[width=4.651cm,height=0.506cm]{./ObjectReplacements/Object 2018}

          Lema de asociatividad de la suma cuando un término es
          complemento de otro.
        \item
          \includegraphics[width=4.013cm,height=0.506cm]{./ObjectReplacements/Object 2019}

          Axioma de elemento neutro del producto.
        \item
          \includegraphics[width=4.013cm,height=0.506cm]{./ObjectReplacements/Object 2020}

          Axioma de conmutatividad de la suma.
        \item
          \includegraphics[width=4.013cm,height=0.506cm]{./ObjectReplacements/Object 2021}

          Axioma de conmutatividad de la suma.
        \item
          \includegraphics[width=4.306cm,height=0.563cm]{./ObjectReplacements/Object 2022}

          Teorema de doble negación es identidad.
        \item
          \includegraphics[width=10.55cm,height=0.467cm]{./ObjectReplacements/Object 2023}

          Lema de asociatividad de la suma cuando un término es
          complemento de otro.
        \item
          \includegraphics[width=1.519cm,height=0.467cm]{./ObjectReplacements/Object 2024}

          Transitividad del la igualdad desde la fórmula (20) a la (32).
        \item
          \includegraphics[width=0.982cm,height=0.467cm]{./ObjectReplacements/Object 2025}

          De (33) y (19) por el axioma de existencia de elemento
          complementario.
        \item
          \includegraphics[width=1.002cm,height=0.467cm]{./ObjectReplacements/Object 2026}

          Teorema de unicidad del complementario.
        \item
          \[\overline{a} = \overline{b}\]

          Simetría de la igualdad.
        \item
          \includegraphics[width=1.002cm,height=0.517cm]{./ObjectReplacements/Object 2028}

          Por el teorema de unicidad del complementario la igualdad
          permanece.
        \item
          \includegraphics[width=1.018cm,height=0.467cm]{./ObjectReplacements/Object 2029}

          Teorema doble negación es identidad.
        \item
          \includegraphics[width=3.036cm,height=0.519cm]{./ObjectReplacements/Object 2030}

          Sustitución en (38) de
          \includegraphics[width=0.441cm,height=0.467cm]{./ObjectReplacements/Object 21}por
          el lado derecho de (1) y de
          \includegraphics[width=0.425cm,height=0.467cm]{./ObjectReplacements/Object 518}por
          el lado izquierdo de (2).
        \end{enumerate}
      \end{enumerate}
    \item
      Propiedad de simplificación de Quine:

      \begin{enumerate}
      \def\labelenumiv{\arabic{enumiv}.}
      \item
        \includegraphics[width=5.967cm,height=0.575cm]{./ObjectReplacements/Object 563}

        Prueba:

        \begin{enumerate}
        \def\labelenumv{\arabic{enumv}.}
        \tightlist
        \item
          \includegraphics[width=3.946cm,height=0.575cm]{./ObjectReplacements/Object 521}
        \item
          \includegraphics[width=3.519cm,height=0.506cm]{./ObjectReplacements/Object 522}
        \item
          \includegraphics[width=3.812cm,height=0.506cm]{./ObjectReplacements/Object 525}
        \item
          \includegraphics[width=5.293cm,height=0.506cm]{./ObjectReplacements/Object 526}
        \item
          \includegraphics[width=5.293cm,height=0.506cm]{./ObjectReplacements/Object 527}
        \item
          \includegraphics[width=6.93cm,height=0.506cm]{./ObjectReplacements/Object 1006}
        \item
          \includegraphics[width=6.93cm,height=0.506cm]{./ObjectReplacements/Object 1009}
        \item
          \includegraphics[width=5.299cm,height=0.506cm]{./ObjectReplacements/Object 1012}
        \item
          \includegraphics[width=4.658cm,height=0.506cm]{./ObjectReplacements/Object 2031}
        \item
          \includegraphics[width=4.658cm,height=0.506cm]{./ObjectReplacements/Object 2032}
        \item
          \includegraphics[width=3.022cm,height=0.506cm]{./ObjectReplacements/Object 2033}
        \item
          \includegraphics[width=3.022cm,height=0.506cm]{./ObjectReplacements/Object 2034}
        \item
          \includegraphics[width=3.717cm,height=0.506cm]{./ObjectReplacements/Object 2035}
        \item
          \includegraphics[width=3.717cm,height=0.506cm]{./ObjectReplacements/Object 2036}
        \item
          \includegraphics[width=2.683cm,height=0.506cm]{./ObjectReplacements/Object 2040}
        \item
          \includegraphics[width=6.073cm,height=0.575cm]{./ObjectReplacements/Object 2041}
        \end{enumerate}
      \item
        \includegraphics[width=6.017cm,height=0.519cm]{./ObjectReplacements/Object 1296}
      \end{enumerate}
    \end{enumerate}
  \end{enumerate}
\end{enumerate}

Prueba:

\begin{enumerate}
\def\labelenumi{\arabic{enumi}.}
\setcounter{enumi}{36}
\item
  \begin{enumerate}
  \def\labelenumii{\arabic{enumii}.}
  \item
    \begin{enumerate}
    \def\labelenumiii{\arabic{enumiii}.}
    \item
      \begin{enumerate}
      \def\labelenumiv{\arabic{enumiv}.}
      \tightlist
      \item
        \begin{enumerate}
        \def\labelenumv{\arabic{enumv}.}
        \tightlist
        \item
          \[{\left( {\left( {x + y} \right) \cdot \left( {x + \overline{z}} \right)} \right) \cdot \left( {y + z} \right)} = {}\]
        \item
          \includegraphics[width=3.812cm,height=0.506cm]{./ObjectReplacements/Object 1294}
        \item
          \includegraphics[width=3.812cm,height=0.506cm]{./ObjectReplacements/Object 1587}
        \item
          \includegraphics[width=5.293cm,height=0.506cm]{./ObjectReplacements/Object 1588}
        \item
          \includegraphics[width=5.293cm,height=0.506cm]{./ObjectReplacements/Object 1589}
        \item
          \includegraphics[width=6.622cm,height=0.506cm]{./ObjectReplacements/Object 1590}
        \item
          \includegraphics[width=6.623cm,height=0.506cm]{./ObjectReplacements/Object 2037}
        \item
          \includegraphics[width=5.309cm,height=0.506cm]{./ObjectReplacements/Object 2038}
        \item
          \includegraphics[width=4.505cm,height=0.506cm]{./ObjectReplacements/Object 2039}
        \item
          \includegraphics[width=4.505cm,height=0.506cm]{./ObjectReplacements/Object 2042}
        \item
          \[{} = {{({{({y + z})} \cdot x})} + y} = {}\]
        \item
          \includegraphics[width=3.177cm,height=0.506cm]{./ObjectReplacements/Object 2044}
        \item
          \includegraphics[width=4.023cm,height=0.506cm]{./ObjectReplacements/Object 2045}
        \item
          \includegraphics[width=4.023cm,height=0.506cm]{./ObjectReplacements/Object 2046}
        \item
          \includegraphics[width=2.836cm,height=0.506cm]{./ObjectReplacements/Object 2047}
        \item
          \includegraphics[width=6.073cm,height=0.575cm]{./ObjectReplacements/Object 562}
        \end{enumerate}
      \end{enumerate}
    \item
      Todos los axiomas de los que hemos partido son reducibles al
      siguiente grupo de axiomas:

      \begin{enumerate}
      \def\labelenumiv{\arabic{enumiv}.}
      \item
        Axiomas para un álgebra de Boole con
        solo\includegraphics[width=1.224cm,height=0.621cm]{./ObjectReplacements/Object 152}(Sheffer
        1933):

        \begin{enumerate}
        \def\labelenumv{\arabic{enumv}.}
        \tightlist
        \item
          \includegraphics[width=1.663cm,height=0.624cm]{./ObjectReplacements/Object 1649}
        \item
          \includegraphics[width=7.848cm,height=0.563cm]{./ObjectReplacements/Object 145}
        \item
          \includegraphics[width=6.292cm,height=0.577cm]{./ObjectReplacements/Object 147}
        \item
          \includegraphics[width=6.851cm,height=0.577cm]{./ObjectReplacements/Object 149}
        \item
          \includegraphics[width=11.654cm,height=0.577cm]{./ObjectReplacements/Object 148}
        \end{enumerate}
      \item
        Es muy fácil demostrar que los axiomas de Huntington de 1904,
        implican estos 5 axiomas de Sheffer. Los puntos 1 y 2 son
        inmediatos, y ya se han definido. También es fácil ver
        que\includegraphics[width=1.566cm,height=0.467cm]{./ObjectReplacements/Object 146}(ya
        se ha visto),
        \[\overline{x}@{\overline{x} = \overline{\overline{x}} = x}\]y
        con ello queda probado
        3.\includegraphics[width=5.611cm,height=0.563cm]{./ObjectReplacements/Object 151},
        y de otro lado queda
        que\includegraphics[width=1.549cm,height=0.467cm]{./ObjectReplacements/Object 1620},
        y el punto 4 queda demostrado. El punto 5 vamos a reducirlo a
        expresiones en sumas y productos y negaciones, aplicaremos la
        distributiva y agruparemos de nuevo la expresión en la forma
        deseada. La parte izquierda de la igualdad 5 es
        \includegraphics[width=4.277cm,height=0.506cm]{./ObjectReplacements/Object 1651}que
        deshaciendo en sumas y negaciones
        sale\includegraphics[width=1.506cm,height=0.503cm]{./ObjectReplacements/Object 1650}donde
        quitando la doble negación y aplicando
        Morgan\includegraphics[width=1.658cm,height=0.506cm]{./ObjectReplacements/Object 1652},
        ahora solo aplicamos
        distributiva\includegraphics[width=2.498cm,height=0.506cm]{./ObjectReplacements/Object 1653}
        y ahora aplicamos
        conmutabilidad\includegraphics[width=2.498cm,height=0.506cm]{./ObjectReplacements/Object 1621}y
        reagrupamos en funciones
        ``NOR''\includegraphics[width=2.043cm,height=0.467cm]{./ObjectReplacements/Object 1654}y
        ya tenemos la igualdad del axioma 5 de Sheffer en su parte
        derecha:\includegraphics[width=4.277cm,height=0.506cm]{./ObjectReplacements/Object 1655}.
      \end{enumerate}
    \item
      Los axiomas de Huntington de 1903 (H03) implican los de Sheffer de
      1933 (S33).

      Diremos\includegraphics[width=2.223cm,height=0.467cm]{./ObjectReplacements/Object 1656}.

      \begin{enumerate}
      \def\labelenumiv{\arabic{enumiv}.}
      \item
        Demostraremos
        que\includegraphics[width=2.223cm,height=0.467cm]{./ObjectReplacements/Object 1657}.

        \begin{enumerate}
        \def\labelenumv{\arabic{enumv}.}
        \item
          Definición\includegraphics[width=2.18cm,height=0.467cm]{./ObjectReplacements/Object 154}.
          Es una aplicación.
        \item
          Axioma 3 de Sheffer en una nueva
          escritura\includegraphics[width=1.386cm,height=0.467cm]{./ObjectReplacements/Object 1619}.
        \item
          Conmutabilidad.
          \includegraphics[width=2.134cm,height=0.467cm]{./ObjectReplacements/Object 150}.

          \begin{enumerate}
          \def\labelenumvi{\arabic{enumvi}.}
          \tightlist
          \item
            \includegraphics[width=7.308cm,height=0.506cm]{./ObjectReplacements/Object 153}
          \item
            \includegraphics[width=10.338cm,height=0.506cm]{./ObjectReplacements/Object 155}
          \item
            \includegraphics[width=6.812cm,height=0.506cm]{./ObjectReplacements/Object 156}
          \item
            \includegraphics[width=6.396cm,height=0.506cm]{./ObjectReplacements/Object 1611}
          \end{enumerate}
        \item
          Existencia de constante.
          \includegraphics[width=2.482cm,height=0.467cm]{./ObjectReplacements/Object 1622}

          \begin{enumerate}
          \def\labelenumvi{\arabic{enumvi}.}
          \tightlist
          \item
            \includegraphics[width=3.438cm,height=0.506cm]{./ObjectReplacements/Object 1658}
          \item
            \includegraphics[width=4.15cm,height=0.506cm]{./ObjectReplacements/Object 1659}
          \item
            \includegraphics[width=3.881cm,height=0.506cm]{./ObjectReplacements/Object 1660}
          \item
            \includegraphics[width=1.842cm,height=0.506cm]{./ObjectReplacements/Object 1661}
          \end{enumerate}
        \item
          Definición\includegraphics[width=1.723cm,height=0.467cm]{./ObjectReplacements/Object 1662}.
          Es una única constante.
        \item
          Definición\includegraphics[width=2.926cm,height=0.506cm]{./ObjectReplacements/Object 1663}.
          Es una única constante.
        \item
          Definición\includegraphics[width=2.623cm,height=0.506cm]{./ObjectReplacements/Object 1664}.
          Es una aplicación.
        \item
          Definición\includegraphics[width=2.351cm,height=0.467cm]{./ObjectReplacements/Object 1665}.
          Es una aplicación.
        \item
          Elemento neutro de la suma.
          \includegraphics[width=1.552cm,height=0.467cm]{./ObjectReplacements/Object 1666}.

          \begin{enumerate}
          \def\labelenumvi{\arabic{enumvi}.}
          \tightlist
          \item
            Prueba:
          \item
            \includegraphics[width=5.801cm,height=0.506cm]{./ObjectReplacements/Object 1667}
          \item
            \includegraphics[width=4.925cm,height=0.506cm]{./ObjectReplacements/Object 1668}
          \end{enumerate}
        \item
          Elemento neutro del producto.
          \includegraphics[width=1.385cm,height=0.467cm]{./ObjectReplacements/Object 1670}.

          \begin{enumerate}
          \def\labelenumvi{\arabic{enumvi}.}
          \tightlist
          \item
            \includegraphics[width=3.607cm,height=0.467cm]{./ObjectReplacements/Object 1671}
          \end{enumerate}
        \item
          Conmutabilidad de la
          suma:\includegraphics[width=2.141cm,height=0.467cm]{./ObjectReplacements/Object 1669}.
          Solo hay que aplicar repetidas veces la conmutabilidad de la
          ``NOR''.
        \item
          Conmutabilidad del
          producto:\includegraphics[width=1.834cm,height=0.467cm]{./ObjectReplacements/Object 1672}.
          Solo hay que aplicar repetidas veces la conmutabilidad de la
          ``NAND''.
        \item
          Distributividad de la suma respecto del producto.
          \includegraphics[width=4.277cm,height=0.506cm]{./ObjectReplacements/Object 1673}

          \begin{enumerate}
          \def\labelenumvi{\arabic{enumvi}.}
          \tightlist
          \item
            \[{{x + {({y \cdot z})}} = {({x@{({y'@z'})}})}}'{' = {({{({y@x})}@{({z@x})}})}}'{' = {{({x + y})} \cdot {({x + z})}}}\]
          \end{enumerate}
        \item
          Distributividad del producto respecto de la
          suma.\includegraphics[width=4.124cm,height=0.506cm]{./ObjectReplacements/Object 1676}

          \begin{enumerate}
          \def\labelenumvi{\arabic{enumvi}.}
          \tightlist
          \item
            \includegraphics[width=7.684cm,height=0.506cm]{./ObjectReplacements/Object 1675}
          \item
            \includegraphics[width=7.19cm,height=0.506cm]{./ObjectReplacements/Object 1677}
          \item
            \includegraphics[width=5.796cm,height=0.506cm]{./ObjectReplacements/Object 1678}
          \end{enumerate}
        \item
          Existe elemento complementario.

          \begin{enumerate}
          \def\labelenumvi{\arabic{enumvi}.}
          \item
            \includegraphics[width=3.263cm,height=0.467cm]{./ObjectReplacements/Object 1679}

            \begin{enumerate}
            \def\labelenumvii{\arabic{enumvii}.}
            \tightlist
            \item
              \includegraphics[width=2.514cm,height=0.467cm]{./ObjectReplacements/Object 1680}
            \item
              \includegraphics[width=4.226cm,height=0.506cm]{./ObjectReplacements/Object 1681}
            \end{enumerate}
          \item
            \includegraphics[width=3.124cm,height=0.467cm]{./ObjectReplacements/Object 1683}

            \begin{enumerate}
            \def\labelenumvii{\arabic{enumvii}.}
            \tightlist
            \item
              \includegraphics[width=2.514cm,height=0.467cm]{./ObjectReplacements/Object 1684}
            \item
              \includegraphics[width=4.168cm,height=0.467cm]{./ObjectReplacements/Object 1682}
            \end{enumerate}
          \end{enumerate}
        \item
          Definición:\includegraphics[width=3.194cm,height=0.467cm]{./ObjectReplacements/Object 1597}
        \end{enumerate}
      \end{enumerate}
    \item
      Todos los axiomas de los que hemos partido son reducibles al
      siguiente grupo de axiomas (nueva interpretación de la flecha de
      Sheffer):

      \begin{enumerate}
      \def\labelenumiv{\arabic{enumiv}.}
      \item
        Axiomas para un álgebra de Boole con
        solo\includegraphics[width=1.224cm,height=0.621cm]{./ObjectReplacements/Object 1686}(Sheffer
        1933):

        \begin{enumerate}
        \def\labelenumv{\arabic{enumv}.}
        \tightlist
        \item
          \includegraphics[width=1.663cm,height=0.624cm]{./ObjectReplacements/Object 1687}
        \item
          \includegraphics[width=7.848cm,height=0.563cm]{./ObjectReplacements/Object 1688}
        \item
          \includegraphics[width=4.778cm,height=0.577cm]{./ObjectReplacements/Object 1689}
        \item
          \includegraphics[width=5.844cm,height=0.577cm]{./ObjectReplacements/Object 1690}
        \item
          \[\forall x,y,{z \in}{({x@{({y@z})}})}@{{({x@{({y@z})}})} = {({{({y@y})}@x})}}@{({{({z@z})}@x})}\]
        \end{enumerate}
      \end{enumerate}
    \item
      Es muy fácil demostrar que los axiomas de Huntington de 1903,
      implican estos 5 axiomas de Sheffer. Los puntos 1 y 2 son
      inmediatos, y ya se han definido. También es fácil ver
      que\includegraphics[width=1.566cm,height=0.467cm]{./ObjectReplacements/Object 1692}(ya
      se ha visto),
      \includegraphics[width=2.163cm,height=0.467cm]{./ObjectReplacements/Object 1693}y
      con ello queda probado
      3.\includegraphics[width=5.457cm,height=0.563cm]{./ObjectReplacements/Object 1694},
      y de otro lado queda que\[x@{1 = \overline{x}}\], y el punto 4
      queda demostrado. El punto 5 lo vamos a reducirlo a expresiones en
      sumas y productos y negaciones, aplicaremos la distributiva y
      agruparemos de nuevo la expresión en la forma deseada. La parte
      izquierda de la igualdad 5 es
      \[{({x@{({y@z})}})}@{({x@{({y@z})}})}\]que deshaciendo en sumas y
      negaciones
      sale\[\overline{\overline{x\cdot\overline{y\cdot z}}}\]donde
      quitando la doble negación y aplicando Morgan
      queda\[x \cdot {({\overline{y} + \overline{z}})}\], ahora solo
      aplicamos
      distributiva\[{({x \cdot \overline{y}})} + {({x \cdot \overline{z}})}\]
      y ahora aplicamos
      conmutabilidad\includegraphics[width=2.344cm,height=0.506cm]{./ObjectReplacements/Object 1700}y
      reagrupamos en funciones
      ``NAND''\includegraphics[width=1.58cm,height=0.467cm]{./ObjectReplacements/Object 1701}y
      ya tenemos la igualdad del axioma (5) de Sheffer en su parte
      derecha:\includegraphics[width=4.277cm,height=0.506cm]{./ObjectReplacements/Object 1702}.
    \item
      Los axiomas de Huntington de 1903 (H03) implican los de Sheffer
      re-interpretados de 1933 (S'33).
      Diremos\includegraphics[width=2.397cm,height=0.467cm]{./ObjectReplacements/Object 1703}.

      \begin{enumerate}
      \def\labelenumiv{\arabic{enumiv}.}
      \item
        Demostraremos
        que\includegraphics[width=2.397cm,height=0.467cm]{./ObjectReplacements/Object 1704}.

        \begin{enumerate}
        \def\labelenumv{\arabic{enumv}.}
        \item
          Definición\includegraphics[width=2.18cm,height=0.467cm]{./ObjectReplacements/Object 1705}.
          Es una aplicación.
        \item
          Axioma 3 de Sheffer en una nueva
          escritura\includegraphics[width=1.386cm,height=0.467cm]{./ObjectReplacements/Object 1706}.
        \item
          Conmutabilidad.
          \includegraphics[width=2.134cm,height=0.467cm]{./ObjectReplacements/Object 1707}.

          \begin{enumerate}
          \def\labelenumvi{\arabic{enumvi}.}
          \tightlist
          \item
            \includegraphics[width=7.308cm,height=0.506cm]{./ObjectReplacements/Object 1708}
          \item
            \includegraphics[width=10.497cm,height=0.506cm]{./ObjectReplacements/Object 1709}
          \item
            \includegraphics[width=6.971cm,height=0.506cm]{./ObjectReplacements/Object 1710}
          \item
            \includegraphics[width=6.396cm,height=0.506cm]{./ObjectReplacements/Object 1711}
          \end{enumerate}
        \item
          Existencia de una expresión constante.
          \includegraphics[width=2.482cm,height=0.467cm]{./ObjectReplacements/Object 1712}

          \begin{enumerate}
          \def\labelenumvi{\arabic{enumvi}.}
          \tightlist
          \item
            \includegraphics[width=3.438cm,height=0.506cm]{./ObjectReplacements/Object 1713}
          \item
            \includegraphics[width=3.881cm,height=0.506cm]{./ObjectReplacements/Object 1714}
          \item
            \includegraphics[width=3.881cm,height=0.506cm]{./ObjectReplacements/Object 1715}
          \item
            \includegraphics[width=1.842cm,height=0.506cm]{./ObjectReplacements/Object 1716}
          \end{enumerate}
        \item
          Definición\includegraphics[width=1.729cm,height=0.467cm]{./ObjectReplacements/Object 1717}.
          Es una única constante.
        \item
          Definición\includegraphics[width=2.944cm,height=0.506cm]{./ObjectReplacements/Object 1718}.
          Es una única constante.
        \item
          Definición\includegraphics[width=2.469cm,height=0.506cm]{./ObjectReplacements/Object 1719}.
          Es una aplicación.
        \item
          Definición\includegraphics[width=2.505cm,height=0.467cm]{./ObjectReplacements/Object 1720}.
          Es una aplicación.
        \item
          Elemento neutro del producto.
          \includegraphics[width=1.385cm,height=0.467cm]{./ObjectReplacements/Object 1721}.

          \begin{enumerate}
          \def\labelenumvi{\arabic{enumvi}.}
          \tightlist
          \item
            \includegraphics[width=5.609cm,height=0.506cm]{./ObjectReplacements/Object 1722}
          \item
            \includegraphics[width=4.886cm,height=0.506cm]{./ObjectReplacements/Object 1723}
          \end{enumerate}
        \item
          Elemento neutro de la suma.
          \includegraphics[width=1.552cm,height=0.467cm]{./ObjectReplacements/Object 1724}.

          \begin{enumerate}
          \def\labelenumvi{\arabic{enumvi}.}
          \tightlist
          \item
            \includegraphics[width=3.761cm,height=0.467cm]{./ObjectReplacements/Object 1725}
          \end{enumerate}
        \item
          Conmutabilidad del
          producto:\includegraphics[width=1.834cm,height=0.467cm]{./ObjectReplacements/Object 1726}.
          Solo hay que aplicar repetidas veces la conmutabilidad de la
          ``NAND''.
        \item
          Conmutabilidad del
          producto:\includegraphics[width=1.834cm,height=0.467cm]{./ObjectReplacements/Object 1727}.
          Solo hay que aplicar repetidas veces la conmutabilidad de la
          ``NOR''.
        \item
          Distributividad del producto respecto de la suma:
          \includegraphics[width=4.124cm,height=0.506cm]{./ObjectReplacements/Object 1728}

          \begin{enumerate}
          \def\labelenumvi{\arabic{enumvi}.}
          \tightlist
          \item
            \includegraphics[width=10.456cm,height=0.506cm]{./ObjectReplacements/Object 1729}
          \end{enumerate}
        \item
          Distributividad de la suma respecto del producto:
          \[{x + {({y \cdot z})}} = {{({x + y})} \cdot {({x + z})}}\]

          \begin{enumerate}
          \def\labelenumvi{\arabic{enumvi}.}
          \tightlist
          \item
            \[{{x + {({y \cdot z})}} = {({{x'}@{({y@z})}'})}}'{' = {({{({y@z})}@{x'}})}}'{' =}\]
          \item
            \includegraphics[width=7.19cm,height=0.506cm]{./ObjectReplacements/Object 1732}
          \item
            \includegraphics[width=5.681cm,height=0.506cm]{./ObjectReplacements/Object 1733}
          \end{enumerate}
        \item
          Existe elemento complementario.

          \begin{enumerate}
          \def\labelenumvi{\arabic{enumvi}.}
          \item
            \includegraphics[width=3.124cm,height=0.467cm]{./ObjectReplacements/Object 1734}

            \begin{enumerate}
            \def\labelenumvii{\arabic{enumvii}.}
            \tightlist
            \item
              \includegraphics[width=2.514cm,height=0.467cm]{./ObjectReplacements/Object 1735}
            \item
              \includegraphics[width=4.073cm,height=0.506cm]{./ObjectReplacements/Object 1736}
            \end{enumerate}
          \item
            \includegraphics[width=3.263cm,height=0.467cm]{./ObjectReplacements/Object 1737}

            \begin{enumerate}
            \def\labelenumvii{\arabic{enumvii}.}
            \tightlist
            \item
              \includegraphics[width=2.514cm,height=0.467cm]{./ObjectReplacements/Object 1738}
            \item
              \[{x + x}{' = {x'@x} = {x@x'} = 1}\]
            \end{enumerate}
          \end{enumerate}
        \item
          Definición:
          \includegraphics[width=3.194cm,height=0.467cm]{./ObjectReplacements/Object 1598}
        \end{enumerate}
      \end{enumerate}
    \item
      Lo más habitual es que la flecha de Sheffer se simbolice
      por\includegraphics[width=0.402cm,height=0.325cm]{./ObjectReplacements/Object 1599}cuando
      es la suma negada, y por
      \includegraphics[width=0.402cm,height=0.325cm]{./ObjectReplacements/Object 1600}cuando
      es el producto negado, tal como lo hemos hecho en el anterior
      texto. Ya que los razonamientos hechos con una de las 2
      interpretaciones son absolutamente simétricos y duales, en vez de
      hacer referencias a la suma o al producto lo hacemos a
      \includegraphics[width=1.473cm,height=0.506cm]{./ObjectReplacements/Object 1601}o
      a
      \includegraphics[width=1.339cm,height=0.467cm]{./ObjectReplacements/Object 1602}o
      idénticamente a
      \includegraphics[width=1.473cm,height=0.506cm]{./ObjectReplacements/Object 1603}o
      a\includegraphics[width=1.339cm,height=0.467cm]{./ObjectReplacements/Object 1604},
      por lo que en vez de diferenciar se usa habitualmente como símbolo
      \includegraphics[width=0.852cm,height=0.489cm]{./ObjectReplacements/Object 1605}o
      ``Sheffer's stroke''.
    \item
      Dualidad en el álgebra de Boole. {[}Metateorema{]}

      \begin{enumerate}
      \def\labelenumiv{\arabic{enumiv}.}
      \tightlist
      \item
        Cambiando simultáneamente todas las ocurrencias entre
        ``0,1,+,*'' por (respectivamente) las de ``1,0,*,+'' obtenemos a
        partir de una ex\-presión otra con el mismo valor de verdad que
        la primera. A medida que vaya\-mos introduciendo nuevos
        operadores veremos como extender esta propiedad de dualidad. Se
        demuestra por inducción sobre cualquier expresión bien formada,
        teniendo en cuenta que todos los axiomas son simétricos, duales.
      \end{enumerate}
    \end{enumerate}
  \end{enumerate}
\end{enumerate}

Prueba:

Primero hay que construir un lenguaje adecuado para las expresiones
booleanas (tal como está sirve para hacer un programa calculadora de
expresiones de Boo\-le, aunque no programable, esto es, no podemos
introducir símbolos nuevos con su definición, por ejemplo funciones
booleanas nuevas, ni tampoco se pueden asignar valores a variables).

Partimos de un lenguaje lógico-formal de 1\textsuperscript{er} orden
(lógica de predicados) con igualdad y pertenencia definida a través de
una teoría de conjuntos con 3 tipos dónde la única primitiva es la
pertenencia a un conjunto y la igualdad de elementos de tipo 0. No vamos
a definir esta parte, sino que la emplearemos de forma transparente, y
enumeraremos las transformaciones válidas. La igualdad para los
restantes tipos es la bi-inclusión clásica.

El lenguaje siguiente no da la probabilidad de hacer definiciones, con
todo es bastante expresivo:

\includegraphics[width=10.294cm,height=3.223cm]{./ObjectReplacements/Object 2066}

\includegraphics[width=14.067cm,height=2.695cm]{./ObjectReplacements/Object 216}

\includegraphics[width=14.545cm,height=2.21cm]{./ObjectReplacements/Object 2051}

\includegraphics[width=16.291cm,height=2.332cm]{./ObjectReplacements/Object 2050}

\includegraphics[width=16.274cm,height=2.461cm]{./ObjectReplacements/Object 2049}

\includegraphics[width=12.929cm,height=1.625cm]{./ObjectReplacements/Object 217}

\includegraphics[width=12.931cm,height=1.685cm]{./ObjectReplacements/Object 2048}

\includegraphics[width=12.931cm,height=1.804cm]{./ObjectReplacements/Object 220}\includegraphics[width=13.099cm,height=1.933cm]{./ObjectReplacements/Object 1184}

\includegraphics[width=15.946cm,height=10.809cm]{./ObjectReplacements/Object 2068}

\includegraphics[width=7.713cm,height=2.736cm]{./ObjectReplacements/Object 218}

\includegraphics[width=12.813cm,height=1.265cm]{./ObjectReplacements/Object 219}

\includegraphics[width=11.578cm,height=0.676cm]{./ObjectReplacements/Object 2073}

\includegraphics[width=12.084cm,height=0.676cm]{./ObjectReplacements/Object 2074}

\includegraphics[width=12.247cm,height=0.676cm]{./ObjectReplacements/Object 2075}

\includegraphics[width=14.757cm,height=3.286cm]{./ObjectReplacements/Object 2072}

\includegraphics[width=13.665cm,height=3.286cm]{./ObjectReplacements/Object 2076}

\includegraphics[width=14.002cm,height=3.286cm]{./ObjectReplacements/Object 2077}

\includegraphics[width=14.259cm,height=2.184cm]{./ObjectReplacements/Object 2078}

\includegraphics[width=8.631cm,height=1.813cm]{./ObjectReplacements/Object 2065}

\includegraphics[width=11.735cm,height=3.406cm]{./ObjectReplacements/Object 2079}

\includegraphics[width=13.571cm,height=2.916cm]{./ObjectReplacements/Object 2080}

\includegraphics[width=13.601cm,height=2.916cm]{./ObjectReplacements/Object 2067}

\includegraphics[width=13.615cm,height=2.916cm]{./ObjectReplacements/Object 2081}

\includegraphics[width=11.414cm,height=3.933cm]{./ObjectReplacements/Object 191}

\includegraphics[width=11.478cm,height=4.353cm]{./ObjectReplacements/Object 2071}

\includegraphics[width=12.181cm,height=4.353cm]{./ObjectReplacements/Object 2070}

\includegraphics[width=12.898cm,height=4.353cm]{./ObjectReplacements/Object 2069}

\includegraphics[width=5.978cm,height=2.471cm]{./ObjectReplacements/Object 17}

\begin{enumerate}
\def\labelenumi{\arabic{enumi}.}
\setcounter{enumi}{37}
\item
  \begin{enumerate}
  \def\labelenumii{\arabic{enumii}.}
  \item
    \begin{enumerate}
    \def\labelenumiii{\arabic{enumiii}.}
    \item
      Operadores ``nor'' y ``nand'' que representaremos respectivamente
      como\includegraphics[width=0.402cm,height=0.325cm]{./ObjectReplacements/Object 32}y\includegraphics[width=0.402cm,height=0.325cm]{./ObjectReplacements/Object 1606}.

      \begin{enumerate}
      \def\labelenumiv{\arabic{enumiv}.}
      \tightlist
      \item
        Definición de
        ``nor'':\includegraphics[width=3.616cm,height=0.467cm]{./ObjectReplacements/Object 33}.
      \item
        Definición de
        ``nand'':\includegraphics[width=3.616cm,height=0.467cm]{./ObjectReplacements/Object 34}.
      \end{enumerate}
    \item
      No asociatividad en general de ``nor'' y de ``nand''. Dar algún
      ejemplo
      en\includegraphics[width=0.637cm,height=0.61cm]{./ObjectReplacements/Object 199}.
    \item
      Cualquier expresión de las que hasta ahora se ha podido utilizar
      es expresable con solo funciones ``nand'' y con solo funciones
      ``nor''.

      \begin{enumerate}
      \def\labelenumiv{\arabic{enumiv}.}
      \item
        ``NOR'':

        \begin{enumerate}
        \def\labelenumv{\arabic{enumv}.}
        \tightlist
        \item
          \includegraphics[width=3.791cm,height=0.506cm]{./ObjectReplacements/Object 35}
        \item
          \includegraphics[width=3.928cm,height=0.506cm]{./ObjectReplacements/Object 36}
        \item
          \includegraphics[width=1.535cm,height=0.467cm]{./ObjectReplacements/Object 37}
        \end{enumerate}
      \item
        ``NAND'':

        \begin{enumerate}
        \def\labelenumv{\arabic{enumv}.}
        \tightlist
        \item
          \includegraphics[width=3.637cm,height=0.506cm]{./ObjectReplacements/Object 39}
        \item
          \includegraphics[width=3.791cm,height=0.506cm]{./ObjectReplacements/Object 40}
        \item
          \includegraphics[width=1.535cm,height=0.467cm]{./ObjectReplacements/Object 38}
        \end{enumerate}
      \end{enumerate}
    \item
      Ahora podemos establecer el álgebra de Boole en solo función de
      operadores ``nor'' o sólo de operadores ``nand''. Los postulados
      de Huntington se establecieron de esta manera, aunque con una
      propiedad que acortaba la longitud total del sistema de axiomas.
      Actualmente se han desarrollado en formas cada vez más cortas
      mediante el postulado de Robinson y los de Wolfram más
      recientemente. (La legibilidad de estos sistemas queda
      definitivamente aniquilada, ya que el sistema de pruebas es por lo
      general un software probador de teoremas automático).
    \item
      Extensión de la dualidad: solo hay que añadir que hay que
      intercambiar todos los operadores ``nor'' por ``nand'' y
      viceversa, además de las sustituciones ya enunciadas en 18.1.
    \item
      Nuevos operadores ``exor'' y ``exnor''.

      \begin{enumerate}
      \def\labelenumiv{\arabic{enumiv}.}
      \tightlist
      \item
        Definición del operador ``exor'':
        \includegraphics[width=6.726cm,height=0.506cm]{./ObjectReplacements/Object 43}
      \item
        Definición del operador ``exnor'':
        \includegraphics[width=6.727cm,height=0.506cm]{./ObjectReplacements/Object 44}
      \end{enumerate}
    \item
      Nueva extensión del teorema de dualidad: solo hay que intercambiar
      ``exor'' por ``ex\-nor '' y viceversa, además de todos los
      intercambios que anteriormente se han des\-crito en 22.
    \item
      Propiedad de elemento inverso (el inverso de cada elemento existe
      y es él mismo):

      \begin{enumerate}
      \def\labelenumiv{\arabic{enumiv}.}
      \item
        Para la operación ``exor'' :
        \includegraphics[width=3.087cm,height=0.563cm]{./ObjectReplacements/Object 57}

        \begin{enumerate}
        \def\labelenumv{\arabic{enumv}.}
        \tightlist
        \item
          Prueba:
        \item
          \includegraphics[width=5.122cm,height=0.506cm]{./ObjectReplacements/Object 159}
        \end{enumerate}
      \item
        Para la operación ``exnor'':
        \includegraphics[width=3.076cm,height=0.563cm]{./ObjectReplacements/Object 58}

        \begin{enumerate}
        \def\labelenumv{\arabic{enumv}.}
        \tightlist
        \item
          Prueba:
        \item
          \includegraphics[width=5.091cm,height=0.506cm]{./ObjectReplacements/Object 160}
        \end{enumerate}
      \end{enumerate}
    \item
      Valor para un elemento operado con su complementario:

      \begin{enumerate}
      \def\labelenumiv{\arabic{enumiv}.}
      \item
        \includegraphics[width=3.074cm,height=0.563cm]{./ObjectReplacements/Object 59}

        \begin{enumerate}
        \def\labelenumv{\arabic{enumv}.}
        \tightlist
        \item
          Prueba:
        \item
          \includegraphics[width=7.191cm,height=0.506cm]{./ObjectReplacements/Object 161}
        \end{enumerate}
      \item
        \includegraphics[width=3.087cm,height=0.563cm]{./ObjectReplacements/Object 60}

        \begin{enumerate}
        \def\labelenumv{\arabic{enumv}.}
        \tightlist
        \item
          Prueba:
        \item
          \includegraphics[width=7.205cm,height=0.506cm]{./ObjectReplacements/Object 162}
        \end{enumerate}
      \end{enumerate}
    \item
      Más valores de estas operaciones:

      \begin{enumerate}
      \def\labelenumiv{\arabic{enumiv}.}
      \item
        \includegraphics[width=3.074cm,height=0.563cm]{./ObjectReplacements/Object 61}

        \begin{enumerate}
        \def\labelenumv{\arabic{enumv}.}
        \tightlist
        \item
          Prueba:
        \item
          \includegraphics[width=6.837cm,height=0.506cm]{./ObjectReplacements/Object 163}
        \end{enumerate}
      \item
        \includegraphics[width=3.087cm,height=0.563cm]{./ObjectReplacements/Object 56}

        \begin{enumerate}
        \def\labelenumv{\arabic{enumv}.}
        \tightlist
        \item
          Prueba:
        \item
          \includegraphics[width=7.004cm,height=0.506cm]{./ObjectReplacements/Object 164}
        \end{enumerate}
      \end{enumerate}
    \item
      Elementos neutros:

      \begin{enumerate}
      \def\labelenumiv{\arabic{enumiv}.}
      \item
        \includegraphics[width=3.087cm,height=0.563cm]{./ObjectReplacements/Object 165}

        \begin{enumerate}
        \def\labelenumv{\arabic{enumv}.}
        \tightlist
        \item
          Prueba:
        \item
          \includegraphics[width=5.75cm,height=0.506cm]{./ObjectReplacements/Object 166}
        \end{enumerate}
      \item
        \includegraphics[width=3.076cm,height=0.563cm]{./ObjectReplacements/Object 168}

        \begin{enumerate}
        \def\labelenumv{\arabic{enumv}.}
        \tightlist
        \item
          Prueba:
        \item
          \includegraphics[width=5.872cm,height=0.506cm]{./ObjectReplacements/Object 167}
        \end{enumerate}
      \end{enumerate}
    \item
      Una propiedad de simetría:

      \begin{enumerate}
      \def\labelenumiv{\arabic{enumiv}.}
      \item
        \includegraphics[width=4.085cm,height=0.563cm]{./ObjectReplacements/Object 70}

        \begin{enumerate}
        \def\labelenumv{\arabic{enumv}.}
        \tightlist
        \item
          Prueba:
        \item
          \includegraphics[width=9.765cm,height=0.506cm]{./ObjectReplacements/Object 207}
        \item
          \includegraphics[width=9.74cm,height=0.529cm]{./ObjectReplacements/Object 169}
        \end{enumerate}
      \item
        \includegraphics[width=4.087cm,height=0.563cm]{./ObjectReplacements/Object 71}

        \begin{enumerate}
        \def\labelenumv{\arabic{enumv}.}
        \tightlist
        \item
          Prueba:
        \item
          \includegraphics[width=9.92cm,height=0.506cm]{./ObjectReplacements/Object 544}
        \item
          \includegraphics[width=9.28cm,height=0.529cm]{./ObjectReplacements/Object 170}
        \end{enumerate}
      \end{enumerate}
    \item
      Los operadores negados ``nexor'' y ``nexnor'' coinciden
      respectivamente con ``exnor'' y ``exor'' (y así no se producen
      nuevos operadores):

      \begin{enumerate}
      \def\labelenumiv{\arabic{enumiv}.}
      \tightlist
      \item
        \includegraphics[width=7.602cm,height=0.563cm]{./ObjectReplacements/Object 68}
      \item
        \[\forall a,{b \in}{}{}{}{\overline{a\odot b} \equiv \overline{a\odot b} = \overline{a}}\odot{b = a}\odot{\overline{b} = a}\oplus b\]
      \end{enumerate}
    \item
      Asociatividad de los nuevos operadores ``exor'' y ``exnor'':

      \begin{enumerate}
      \def\labelenumiv{\arabic{enumiv}.}
      \item
        \includegraphics[width=6.258cm,height=0.577cm]{./ObjectReplacements/Object 171}

        \begin{enumerate}
        \def\labelenumv{\arabic{enumv}.}
        \tightlist
        \item
          Prueba:
        \item
          \includegraphics[width=2.263cm,height=0.506cm]{./ObjectReplacements/Object 370}
        \item
          \includegraphics[width=5.537cm,height=0.506cm]{./ObjectReplacements/Object 371}
        \item
          \includegraphics[width=8.864cm,height=0.563cm]{./ObjectReplacements/Object 372}
        \item
          \includegraphics[width=8.864cm,height=0.563cm]{./ObjectReplacements/Object 373}
        \item
          \includegraphics[width=8.475cm,height=0.563cm]{./ObjectReplacements/Object 374}
        \item
          \includegraphics[width=8.782cm,height=0.506cm]{./ObjectReplacements/Object 375}
        \item
          \includegraphics[width=8.393cm,height=0.506cm]{./ObjectReplacements/Object 376}
        \item
          \includegraphics[width=8.992cm,height=0.506cm]{./ObjectReplacements/Object 377}
        \item
          \includegraphics[width=6.412cm,height=0.506cm]{./ObjectReplacements/Object 378}
        \item
          \includegraphics[width=6.412cm,height=0.563cm]{./ObjectReplacements/Object 379}
        \item
          \includegraphics[width=3.971cm,height=0.506cm]{./ObjectReplacements/Object 172}
        \end{enumerate}
      \item
        \includegraphics[width=6.262cm,height=0.577cm]{./ObjectReplacements/Object 174}

        \begin{enumerate}
        \def\labelenumv{\arabic{enumv}.}
        \tightlist
        \item
          Prueba:
        \item
          \includegraphics[width=2.265cm,height=0.506cm]{./ObjectReplacements/Object 380}
        \item
          \includegraphics[width=5.69cm,height=0.506cm]{./ObjectReplacements/Object 381}
        \item
          \includegraphics[width=9.326cm,height=0.563cm]{./ObjectReplacements/Object 382}
        \item
          \includegraphics[width=9.326cm,height=0.563cm]{./ObjectReplacements/Object 383}
        \item
          \includegraphics[width=9.4cm,height=0.563cm]{./ObjectReplacements/Object 384}
        \item
          \includegraphics[width=9.089cm,height=0.506cm]{./ObjectReplacements/Object 385}
        \item
          \includegraphics[width=8.393cm,height=0.506cm]{./ObjectReplacements/Object 386}
        \item
          \includegraphics[width=8.992cm,height=0.506cm]{./ObjectReplacements/Object 387}
        \item
          \includegraphics[width=6.565cm,height=0.506cm]{./ObjectReplacements/Object 388}
        \item
          \includegraphics[width=6.412cm,height=0.563cm]{./ObjectReplacements/Object 389}
        \item
          \includegraphics[width=3.972cm,height=0.506cm]{./ObjectReplacements/Object 173}
        \end{enumerate}
      \end{enumerate}
    \item
      Distributividad de
      ``\includegraphics[width=0.572cm,height=0.467cm]{./ObjectReplacements/Object 175}''
      respecto del producto lógico
      ``\includegraphics[width=0.445cm,height=0.467cm]{./ObjectReplacements/Object 176}''
      y de
      ``\includegraphics[width=0.573cm,height=0.467cm]{./ObjectReplacements/Object 177}''
      res\-pecto de la suma lógica
      ``\includegraphics[width=0.492cm,height=0.467cm]{./ObjectReplacements/Object 178}'':

      \begin{enumerate}
      \def\labelenumiv{\arabic{enumiv}.}
      \item
        \includegraphics[width=8.186cm,height=0.577cm]{./ObjectReplacements/Object 179}

        \begin{enumerate}
        \def\labelenumv{\arabic{enumv}.}
        \item
          Prueba:
        \item
          \includegraphics[width=4.23cm,height=0.506cm]{./ObjectReplacements/Object 221}
        \item
          \includegraphics[width=5.722cm,height=0.563cm]{./ObjectReplacements/Object 390}
        \item
          \includegraphics[width=5.292cm,height=0.506cm]{./ObjectReplacements/Object 391}
        \item
          \includegraphics[width=5.292cm,height=0.563cm]{./ObjectReplacements/Object 392}
        \item
          \includegraphics[width=9.11cm,height=0.563cm]{./ObjectReplacements/Object 393}
        \item
          \includegraphics[width=11.217cm,height=0.506cm]{./ObjectReplacements/Object 394}
        \item
          \includegraphics[width=11.896cm,height=0.506cm]{./ObjectReplacements/Object 395}
        \item
          \includegraphics[width=6.611cm,height=0.506cm]{./ObjectReplacements/Object 396}
        \item
          \includegraphics[width=8.842cm,height=0.506cm]{./ObjectReplacements/Object 397}
        \item
          \includegraphics[width=8.638cm,height=0.506cm]{./ObjectReplacements/Object 398}
        \item
          \includegraphics[width=2.417cm,height=0.467cm]{./ObjectReplacements/Object 399}
        \item
          \includegraphics[width=5.86cm,height=0.563cm]{./ObjectReplacements/Object 222}
        \item
          \includegraphics[width=5.445cm,height=0.506cm]{./ObjectReplacements/Object 400}
        \item
          \includegraphics[width=5.445cm,height=0.563cm]{./ObjectReplacements/Object 401}
        \item
          \includegraphics[width=9.264cm,height=0.563cm]{./ObjectReplacements/Object 402}
        \item
          \includegraphics[width=8.114cm,height=0.506cm]{./ObjectReplacements/Object 403}
        \item
          \includegraphics[width=7.006cm,height=0.506cm]{./ObjectReplacements/Object 404}
        \item
          \includegraphics[width=5.838cm,height=0.467cm]{./ObjectReplacements/Object 405}
        \item
          \includegraphics[width=12.024cm,height=0.467cm]{./ObjectReplacements/Object 406}
        \item
          \includegraphics[width=10.756cm,height=0.467cm]{./ObjectReplacements/Object 407}
        \item
          \includegraphics[width=9.128cm,height=0.506cm]{./ObjectReplacements/Object 408}
        \item
          \includegraphics[width=4.212cm,height=0.467cm]{./ObjectReplacements/Object 409}
        \item
          \includegraphics[width=3.992cm,height=0.506cm]{./ObjectReplacements/Object 410}
        \item
          \includegraphics[width=1.819cm,height=0.467cm]{./ObjectReplacements/Object 411}

          \includegraphics[width=10.414cm,height=0.563cm]{./ObjectReplacements/Object 223}

          \includegraphics[width=10.091cm,height=0.467cm]{./ObjectReplacements/Object 412}
        \item
          \includegraphics[width=4.23cm,height=0.506cm]{./ObjectReplacements/Object 413}
        \end{enumerate}

        Es seguro que la prueba anterior puede ser acortada
        drásticamente, así que si al\-guno encuentra una forma (quizás
        más directa) la pondremos en su lugar.
      \item
        \includegraphics[width=8.652cm,height=0.577cm]{./ObjectReplacements/Object 180}Se
        prueba como en el caso anterior, sólo que cambiando los
        operadores duales, y las dos constantes
        \includegraphics[width=1.161cm,height=0.492cm]{./ObjectReplacements/Object 192}
        entre sí y obtenemos el resultado que hemos enunciado.

        \begin{enumerate}
        \def\labelenumv{\arabic{enumv}.}
        \item
          Prueba:
        \item
          \includegraphics[width=4.694cm,height=0.506cm]{./ObjectReplacements/Object 414}
        \item
          \includegraphics[width=6.339cm,height=0.563cm]{./ObjectReplacements/Object 415}
        \item
          \includegraphics[width=5.907cm,height=0.506cm]{./ObjectReplacements/Object 416}
        \item
          \includegraphics[width=5.909cm,height=0.563cm]{./ObjectReplacements/Object 417}
        \item
          \includegraphics[width=10.342cm,height=0.563cm]{./ObjectReplacements/Object 418}
        \item
          \includegraphics[width=12.448cm,height=0.506cm]{./ObjectReplacements/Object 419}
        \item
          \includegraphics[width=11.583cm,height=0.506cm]{./ObjectReplacements/Object 594}
        \item
          \includegraphics[width=4.147cm,height=0.506cm]{./ObjectReplacements/Object 420}
        \item
          \includegraphics[width=6.897cm,height=0.506cm]{./ObjectReplacements/Object 421}
        \item
          \includegraphics[width=10.075cm,height=0.506cm]{./ObjectReplacements/Object 422}
        \item
          \includegraphics[width=10.022cm,height=0.506cm]{./ObjectReplacements/Object 423}
        \item
          \includegraphics[width=2.064cm,height=0.467cm]{./ObjectReplacements/Object 424}
        \item
          \includegraphics[width=6.168cm,height=0.563cm]{./ObjectReplacements/Object 425}
        \item
          \includegraphics[width=5.754cm,height=0.506cm]{./ObjectReplacements/Object 426}
        \item
          \includegraphics[width=5.756cm,height=0.563cm]{./ObjectReplacements/Object 427}
        \item
          \includegraphics[width=10.188cm,height=0.563cm]{./ObjectReplacements/Object 428}
        \item
          \includegraphics[width=8.29cm,height=0.575cm]{./ObjectReplacements/Object 429}
        \item
          \includegraphics[width=7.761cm,height=0.506cm]{./ObjectReplacements/Object 430}
        \item
          \includegraphics[width=10.888cm,height=0.506cm]{./ObjectReplacements/Object 431}\includegraphics[width=5.955cm,height=0.506cm]{./ObjectReplacements/Object 432}
        \item
          \includegraphics[width=10.888cm,height=0.506cm]{./ObjectReplacements/Object 593}\includegraphics[width=4.24cm,height=0.506cm]{./ObjectReplacements/Object 433}
        \item
          \includegraphics[width=12.511cm,height=0.506cm]{./ObjectReplacements/Object 434}
        \item
          \includegraphics[width=5.533cm,height=0.506cm]{./ObjectReplacements/Object 435}
        \item
          \includegraphics[width=4.42cm,height=0.506cm]{./ObjectReplacements/Object 436}
        \item
          \includegraphics[width=1.676cm,height=0.467cm]{./ObjectReplacements/Object 437}

          \includegraphics[width=10.876cm,height=0.563cm]{./ObjectReplacements/Object 438}

          \includegraphics[width=10.091cm,height=0.467cm]{./ObjectReplacements/Object 439}
        \item
          \includegraphics[width=4.694cm,height=0.506cm]{./ObjectReplacements/Object 181}
        \end{enumerate}
      \end{enumerate}
    \item
      Estructuras de anillo conmutativo con elemento unidad (es claro
      desde todas las pro\-piedades anteriormente demostradas):

      \begin{enumerate}
      \def\labelenumiv{\arabic{enumiv}.}
      \tightlist
      \item
        La más normal sería:
        \includegraphics[width=4.314cm,height=0.577cm]{./ObjectReplacements/Object 45}
      \item
        Su forma dual es
        :\includegraphics[width=4.47cm,height=0.577cm]{./ObjectReplacements/Object 46}
      \end{enumerate}
    \item
      Estructuras respectivas a 33 de bimódulo
      de\includegraphics[width=1.445cm,height=0.577cm]{./ObjectReplacements/Object 72}sobre
      el anillo
      \includegraphics[width=1.647cm,height=0.577cm]{./ObjectReplacements/Object 73}
      y el dual, de
      \includegraphics[width=1.445cm,height=0.577cm]{./ObjectReplacements/Object 74}sobre
      el anillo
      \includegraphics[width=1.803cm,height=0.577cm]{./ObjectReplacements/Object 75}.
    \end{enumerate}

    A continuación hablaremos sobre como son en general las álgebras de
    Boole, funda\-mentalmente las finitas, y veremos que efectivamente
    podemos llegar a teoremas que nos dicen de forma muy concreta cuales
    son estas álgebras de Boole. Vamos a ver for\-mas de generarlas y
    cuestiones parecidas.

    \begin{enumerate}
    \def\labelenumiii{\arabic{enumiii}.}
    \setcounter{enumiii}{17}
    \item
      En el caso que el cardinal
      de\includegraphics[width=0.51cm,height=0.563cm]{./ObjectReplacements/Object 47}sea
      finito,\includegraphics[width=1.431cm,height=0.577cm]{./ObjectReplacements/Object 48}.
      Para demostrarlo solo hay que darse cuenta
      que\includegraphics[width=2.706cm,height=0.667cm]{./ObjectReplacements/Object 182},
      que
      \includegraphics[width=6.339cm,height=0.506cm]{./ObjectReplacements/Object 183}y
      que\includegraphics[width=3.508cm,height=0.572cm]{./ObjectReplacements/Object 186}y
      así
      cuando\includegraphics[width=0.51cm,height=0.563cm]{./ObjectReplacements/Object 187}sea
      finito, su cardinal será un múlti\-plo de 2.
    \item
      Definición:

      \includegraphics[width=6.683cm,height=0.563cm]{./ObjectReplacements/Object 1142}
    \item
      \includegraphics[width=3.575cm,height=0.621cm]{./ObjectReplacements/Object 1144}

      \begin{enumerate}
      \def\labelenumiv{\arabic{enumiv}.}
      \item
        Reflexiva:
        \includegraphics[width=2.609cm,height=0.563cm]{./ObjectReplacements/Object 1164}

        \begin{enumerate}
        \def\labelenumv{\arabic{enumv}.}
        \tightlist
        \item
          \includegraphics[width=2.986cm,height=0.563cm]{./ObjectReplacements/Object 1145}
        \item

          \includegraphics[width=2.609cm,height=0.563cm]{./ObjectReplacements/Object 1158}
        \end{enumerate}
      \item
        Antisimétrica:
        \includegraphics[width=5.759cm,height=0.563cm]{./ObjectReplacements/Object 1147}

        \begin{enumerate}
        \def\labelenumv{\arabic{enumv}.}
        \tightlist
        \item
          \includegraphics[width=3.881cm,height=0.563cm]{./ObjectReplacements/Object 1146}
        \item
          \includegraphics[width=4.658cm,height=0.563cm]{./ObjectReplacements/Object 1148}
        \item
          \includegraphics[width=2.672cm,height=0.563cm]{./ObjectReplacements/Object 1149}
        \end{enumerate}
      \item
        Transitiva:
        \includegraphics[width=6.133cm,height=0.563cm]{./ObjectReplacements/Object 1150}

        \begin{enumerate}
        \def\labelenumv{\arabic{enumv}.}
        \tightlist
        \item

          \includegraphics[width=2.244cm,height=0.467cm]{./ObjectReplacements/Object 1151}
        \item
          \includegraphics[width=3.022cm,height=0.467cm]{./ObjectReplacements/Object 1152}
        \item
          \includegraphics[width=2.15cm,height=0.467cm]{./ObjectReplacements/Object 1153}
        \item
          \includegraphics[width=1.794cm,height=0.467cm]{./ObjectReplacements/Object 1154}
        \item
          \includegraphics[width=1.395cm,height=0.467cm]{./ObjectReplacements/Object 1155}
        \item
          \includegraphics[width=1.016cm,height=0.467cm]{./ObjectReplacements/Object 1156}
        \end{enumerate}
      \end{enumerate}
    \item
      Definición:\includegraphics[width=6.683cm,height=0.612cm]{./ObjectReplacements/Object 1143}
    \item
      \includegraphics[width=3.575cm,height=0.621cm]{./ObjectReplacements/Object 1177}

      \begin{enumerate}
      \def\labelenumiv{\arabic{enumiv}.}
      \item
        Reflexiva:
        \includegraphics[width=2.609cm,height=0.563cm]{./ObjectReplacements/Object 1180}

        \begin{enumerate}
        \def\labelenumv{\arabic{enumv}.}
        \tightlist
        \item
          \includegraphics[width=3.141cm,height=0.563cm]{./ObjectReplacements/Object 1182}
        \item

          \includegraphics[width=2.609cm,height=0.563cm]{./ObjectReplacements/Object 1183}
        \end{enumerate}
      \item
        Antisimétrica:
        \includegraphics[width=5.759cm,height=0.563cm]{./ObjectReplacements/Object 1186}

        \begin{enumerate}
        \def\labelenumv{\arabic{enumv}.}
        \tightlist
        \item
          \includegraphics[width=3.881cm,height=0.563cm]{./ObjectReplacements/Object 1187}
        \item
          \includegraphics[width=4.965cm,height=0.563cm]{./ObjectReplacements/Object 1188}
        \item
          \includegraphics[width=2.672cm,height=0.563cm]{./ObjectReplacements/Object 1189}
        \end{enumerate}
      \item
        Transitiva:
        \includegraphics[width=6.133cm,height=0.563cm]{./ObjectReplacements/Object 1190}

        \begin{enumerate}
        \def\labelenumv{\arabic{enumv}.}
        \tightlist
        \item

          \includegraphics[width=2.244cm,height=0.467cm]{./ObjectReplacements/Object 1191}
        \item
          \includegraphics[width=3.328cm,height=0.467cm]{./ObjectReplacements/Object 1192}
        \item
          \includegraphics[width=2.612cm,height=0.467cm]{./ObjectReplacements/Object 1193}
        \item
          \includegraphics[width=2.101cm,height=0.467cm]{./ObjectReplacements/Object 1194}
        \item
          \includegraphics[width=1.549cm,height=0.467cm]{./ObjectReplacements/Object 1195}
        \item
          \includegraphics[width=1.016cm,height=0.467cm]{./ObjectReplacements/Object 1170}
        \end{enumerate}
      \end{enumerate}
    \item
      Definición:\includegraphics[width=10.456cm,height=0.9cm]{./ObjectReplacements/Object 1136}
    \item
      \includegraphics[width=7.08cm,height=0.577cm]{./ObjectReplacements/Object 1138}.
    \item
      Definición:\includegraphics[width=12.307cm,height=0.834cm]{./ObjectReplacements/Object 1137}
    \item
      \includegraphics[width=7.641cm,height=0.577cm]{./ObjectReplacements/Object 1157}
    \item
      Definición\includegraphics[width=5.756cm,height=0.667cm]{./ObjectReplacements/Object 1159}
    \item
      \includegraphics[width=5.8cm,height=0.587cm]{./ObjectReplacements/Object 1173}
    \item
      Definición\includegraphics[width=5.794cm,height=0.667cm]{./ObjectReplacements/Object 1160}
    \item
      \includegraphics[width=5.856cm,height=0.587cm]{./ObjectReplacements/Object 1174}
    \item
      Definición\includegraphics[width=8.594cm,height=0.651cm]{./ObjectReplacements/Object 1199}
    \item
      Definición\[x,{y \in}\mspace{72mu}{x \geq y}\qquad\Rightarrow\qquad\left\lbrack {x,y} \right)_{}{: = {\{{{z \in} \mid {y \leq {z \land x} \geq {z \land z} \neq y}}\}}}\]
    \item
      Definición\includegraphics[width=9.744cm,height=0.651cm]{./ObjectReplacements/Object 1197}
    \item
      Definición\includegraphics[width=10.926cm,height=0.661cm]{./ObjectReplacements/Object 1198}
    \item
      \includegraphics[width=5.341cm,height=0.667cm]{./ObjectReplacements/Object 1168}
    \item
      \includegraphics[width=5.343cm,height=0.667cm]{./ObjectReplacements/Object 1169}
    \item
      \includegraphics[width=5.36cm,height=0.667cm]{./ObjectReplacements/Object 1104}
    \item
      \includegraphics[width=9.825cm,height=0.667cm]{./ObjectReplacements/Object 1171}
    \item
      \includegraphics[width=10.021cm,height=0.667cm]{./ObjectReplacements/Object 1172}
    \item
      Definición\includegraphics[width=5.463cm,height=0.661cm]{./ObjectReplacements/Object 1161}
    \item
      Definición\includegraphics[width=5.911cm,height=0.661cm]{./ObjectReplacements/Object 1162}
    \item
      \includegraphics[width=4.934cm,height=0.61cm]{./ObjectReplacements/Object 1163}
    \item
      Definición
      \includegraphics[width=5.914cm,height=0.681cm]{./ObjectReplacements/Object 1175}
    \item
      \includegraphics[width=1.496cm,height=0.624cm]{./ObjectReplacements/Object 1178}Pues
      es un álgebra de cardinal mayor o igual que
      \includegraphics[width=0.413cm,height=0.467cm]{./ObjectReplacements/Object 1307}
      y existe al
      menos\includegraphics[width=0.743cm,height=0.492cm]{./ObjectReplacements/Object 1141}.
    \item
      \includegraphics[width=3.574cm,height=0.624cm]{./ObjectReplacements/Object 1176}
    \item
      \includegraphics[width=11.675cm,height=0.905cm]{./ObjectReplacements/Object 1117}

      \begin{enumerate}
      \def\labelenumiv{\arabic{enumiv}.}
      \tightlist
      \item
        Prueba:
      \item
        \includegraphics[width=4.128cm,height=0.487cm]{./ObjectReplacements/Object 1179}
      \item
        \includegraphics[width=6.025cm,height=0.467cm]{./ObjectReplacements/Object 1185}
      \item
        \includegraphics[width=11.499cm,height=0.575cm]{./ObjectReplacements/Object 1181}
      \item
        \includegraphics[width=11.411cm,height=0.575cm]{./ObjectReplacements/Object 1106}
      \item
        \includegraphics[width=6.064cm,height=0.547cm]{./ObjectReplacements/Object 1107}
      \item
        \includegraphics[width=4.207cm,height=0.905cm]{./ObjectReplacements/Object 1108}
      \item
        \includegraphics[width=2.023cm,height=0.905cm]{./ObjectReplacements/Object 1109}
      \item
        \includegraphics[width=2.42cm,height=0.905cm]{./ObjectReplacements/Object 1110}
      \item
        \includegraphics[width=2.498cm,height=0.905cm]{./ObjectReplacements/Object 1111}
      \item
        \includegraphics[width=3.833cm,height=0.905cm]{./ObjectReplacements/Object 1118}
      \item
        \includegraphics[width=1.894cm,height=0.587cm]{./ObjectReplacements/Object 1112}
      \item
        \includegraphics[width=5.63cm,height=0.905cm]{./ObjectReplacements/Object 1113}
      \item
        \includegraphics[width=4.283cm,height=0.905cm]{./ObjectReplacements/Object 1114}
      \item
        \includegraphics[width=4.785cm,height=0.905cm]{./ObjectReplacements/Object 1115}
      \item
        \includegraphics[width=7.093cm,height=0.905cm]{./ObjectReplacements/Object 1119}
      \item
        \includegraphics[width=9.987cm,height=0.905cm]{./ObjectReplacements/Object 1116}
      \end{enumerate}
    \item
      \includegraphics[width=4.657cm,height=0.563cm]{./ObjectReplacements/Object 1120}
      Desde 57 es inmediato.
    \item
      Ahora vamos a construir una función inyectiva del álgebra de Boole
      de las partes de los átomos de
      \includegraphics[width=0.51cm,height=0.563cm]{./ObjectReplacements/Object 1308}
      (si este es finito) en el álgebra de Boole
      \includegraphics[width=0.51cm,height=0.563cm]{./ObjectReplacements/Object 1309}.
      Así cuando menos sa\-bremos que podemos interpretar esta álgebra
      de las partes de un conjunto de los áto\-mos de
      \includegraphics[width=0.51cm,height=0.563cm]{./ObjectReplacements/Object 1310}
      como un subálgebra de la que estamos estudiando.

      La función
      \includegraphics[width=0.471cm,height=0.319cm]{./ObjectReplacements/Object 1121}que
      vamos a definir va a quedar completamente definida en la fór\-mula
      que sigue. Tendremos que mostrar que está bien definida, que es
      inyectiva, que
      \includegraphics[width=4.032cm,height=0.506cm]{./ObjectReplacements/Object 1122},
      esto es que respeta la suma booleana en
      \includegraphics[width=0.51cm,height=0.563cm]{./ObjectReplacements/Object 1123}que
      viene como unión de conjuntos
      desde\includegraphics[width=2.538cm,height=0.639cm]{./ObjectReplacements/Object 1124},
      que\includegraphics[width=3.879cm,height=0.506cm]{./ObjectReplacements/Object 1125},
      esto es que respeta el producto booleano
      en\includegraphics[width=0.51cm,height=0.563cm]{./ObjectReplacements/Object 195}que
      viene como intersección de conjun\-tos
      desde\includegraphics[width=2.538cm,height=0.639cm]{./ObjectReplacements/Object 1127},
      y aunque ya no sería necesario, también veremos
      que\includegraphics[width=4.038cm,height=0.639cm]{./ObjectReplacements/Object 1128}.
      Así quedará clara la relación entre ambas álgebras.

      \begin{enumerate}
      \def\labelenumiv{\arabic{enumiv}.}
      \tightlist
      \item
        \includegraphics[width=11.557cm,height=6.468cm]{./ObjectReplacements/Object 1093}.
      \end{enumerate}
    \item
      \includegraphics[width=4.281cm,height=0.61cm]{./ObjectReplacements/Object 1166}Con
      la
      misma\includegraphics[width=0.471cm,height=0.319cm]{./ObjectReplacements/Object 1129}anterior.
    \item
      \includegraphics[width=4.307cm,height=0.61cm]{./ObjectReplacements/Object 1167}Con
      la misma\[\varphi\]anterior.
    \item
      \includegraphics[width=5.662cm,height=0.563cm]{./ObjectReplacements/Object 1165}Lo
      mejor sería demostrar que cualquier cadena completa saturada de 1
      a 0 tiene una longitud (número de elementos) siempre igual
      al\includegraphics[width=2.103cm,height=0.577cm]{./ObjectReplacements/Object 1131}.
      De ahí se sigue que si el cardinal es el dicho, tendríamos más de
      3 niveles, diferenciándose siempre los átomos y los hiperátomos.
    \item
      Para cada uno de los cardinales de
      \includegraphics[width=0.51cm,height=0.563cm]{./ObjectReplacements/Object 49},
      cuando son finitos, existe una estructura no solo de anillo
      conmutativo con unidad como en 34, sino también de cuerpo. La
      po\-demos encontrar explícitamente en el álgebra de las partes de
      un conjunto finito. Sólo nos queda ver que
      \includegraphics[width=0.471cm,height=0.319cm]{./ObjectReplacements/Object 1132}es
      sobreyectivo. Tenemos
      que\includegraphics[width=3.027cm,height=0.639cm]{./ObjectReplacements/Object 1133}.
      Así
      \includegraphics[width=0.471cm,height=0.319cm]{./ObjectReplacements/Object 1134}pasa
      a ser un isomorfismo de álgebras de Boole: esto es, en lo que a la
      estructu\-ra de álgebra de Boole se refiere, haciendo abstracción
      de las
      operaciones\includegraphics[width=0.492cm,height=0.467cm]{./ObjectReplacements/Object 1139}y
      \includegraphics[width=0.445cm,height=0.467cm]{./ObjectReplacements/Object 1135}concretas
      y los elementos
      concretos,\includegraphics[width=8.954cm,height=0.637cm]{./ObjectReplacements/Object 1140}.
    \item
      De 74 se deduce que si
      \includegraphics[width=0.51cm,height=0.563cm]{./ObjectReplacements/Object 1607}es
      un conjunto
      finito,\includegraphics[width=3.004cm,height=0.563cm]{./ObjectReplacements/Object 1608}.
    \item
      Estructuras respectivas a 35 de espacio vectorial
      de\includegraphics[width=1.445cm,height=0.577cm]{./ObjectReplacements/Object 76}sobre
      el cuerpo
      \includegraphics[width=1.647cm,height=0.577cm]{./ObjectReplacements/Object 77}y
      el dual de
      \includegraphics[width=1.445cm,height=0.577cm]{./ObjectReplacements/Object 78}sobre
      el cuerpo
      dual\includegraphics[width=1.803cm,height=0.577cm]{./ObjectReplacements/Object 79}.
      Éste último es el caso
      cuando\includegraphics[width=2.514cm,height=0.61cm]{./ObjectReplacements/Object 188}.
      Esto tendrá utilidad inmediata en los códigos de Hamming.
    \item
      Para calcular los inversos en los cuerpos finitos de
      cardinal\includegraphics[width=0.577cm,height=0.527cm]{./ObjectReplacements/Object 1609}correspondientes
      hay que re\-solver algunas ecuaciones sobre igualdades
      polinómicas. El producto del cuerpo fi\-nito asociado (en número
      de elementos) a nuestro álgebra de Boole no es en general igual al
      producto del ani\-llo booleano asociado. Tiene que ver con los
      cuerpos de
      Galois\includegraphics[width=1.48cm,height=0.557cm]{./ObjectReplacements/Object 50}.
      Estos cuerpos y los polinomios mencionados son de gran utilidad en
      teoría de codificación.
    \item
      Existen formulas sencillas para poner la suma
      ``\includegraphics[width=0.492cm,height=0.467cm]{./ObjectReplacements/Object 51}'',
      el producto
      ``\includegraphics[width=0.445cm,height=0.467cm]{./ObjectReplacements/Object 52}''
      en fun\-ción de las funciones
      ``\includegraphics[width=0.572cm,height=0.467cm]{./ObjectReplacements/Object 53}''
      y
      ``\includegraphics[width=0.445cm,height=0.467cm]{./ObjectReplacements/Object 54}'',
      y de
      ``\includegraphics[width=0.573cm,height=0.467cm]{./ObjectReplacements/Object 55}''
      y
      ``\includegraphics[width=0.492cm,height=0.467cm]{./ObjectReplacements/Object 62}'':

      \begin{enumerate}
      \def\labelenumiv{\arabic{enumiv}.}
      \item
        \includegraphics[width=3.665cm,height=0.506cm]{./ObjectReplacements/Object 63}

        \begin{enumerate}
        \def\labelenumv{\arabic{enumv}.}
        \tightlist
        \item
          Prueba:
        \item
          \includegraphics[width=2.928cm,height=0.506cm]{./ObjectReplacements/Object 251}
        \item
          \includegraphics[width=6.833cm,height=0.506cm]{./ObjectReplacements/Object 252}
        \item
          \includegraphics[width=10.791cm,height=0.563cm]{./ObjectReplacements/Object 253}
        \item
          \includegraphics[width=10.791cm,height=0.563cm]{./ObjectReplacements/Object 254}
        \item
          \includegraphics[width=11.098cm,height=0.506cm]{./ObjectReplacements/Object 255}
        \item
          \includegraphics[width=7.303cm,height=0.506cm]{./ObjectReplacements/Object 256}
        \item
          \includegraphics[width=6.461cm,height=0.506cm]{./ObjectReplacements/Object 257}
        \item
          \includegraphics[width=8.177cm,height=0.506cm]{./ObjectReplacements/Object 258}
        \item
          \includegraphics[width=4.537cm,height=0.482cm]{./ObjectReplacements/Object 259}
        \item
          \includegraphics[width=3.635cm,height=0.467cm]{./ObjectReplacements/Object 193}
        \end{enumerate}
      \item
        \includegraphics[width=3.667cm,height=0.506cm]{./ObjectReplacements/Object 64}

        \begin{enumerate}
        \def\labelenumv{\arabic{enumv}.}
        \tightlist
        \item
          Prueba:
        \item
          \includegraphics[width=3.083cm,height=0.506cm]{./ObjectReplacements/Object 260}
        \item
          \includegraphics[width=7.2cm,height=0.506cm]{./ObjectReplacements/Object 261}
        \item
          \includegraphics[width=11.772cm,height=0.563cm]{./ObjectReplacements/Object 262}
        \item
          \includegraphics[width=11.772cm,height=0.563cm]{./ObjectReplacements/Object 263}
        \item
          \includegraphics[width=11.308cm,height=0.506cm]{./ObjectReplacements/Object 264}
        \item
          \includegraphics[width=7.668cm,height=0.506cm]{./ObjectReplacements/Object 265}
        \item
          \includegraphics[width=6.826cm,height=0.506cm]{./ObjectReplacements/Object 266}
        \item
          \includegraphics[width=8.543cm,height=0.506cm]{./ObjectReplacements/Object 267}
        \item
          \includegraphics[width=4.593cm,height=0.482cm]{./ObjectReplacements/Object 268}
        \item
          \includegraphics[width=3.692cm,height=0.467cm]{./ObjectReplacements/Object 194}
        \end{enumerate}
      \end{enumerate}
    \item
      Si un
      anillo\includegraphics[width=2.379cm,height=0.577cm]{./ObjectReplacements/Object 190}es
      tal
      que\includegraphics[width=2.863cm,height=0.563cm]{./ObjectReplacements/Object 189},
      define de manera unívoca un álgebra de Boole (la estructura de la
      que hablamos se llama un anillo de Boole). Esta proposición, con
      ser matemáticamente importante, la vemos aquí como sólo una
      curiosidad. En el anillo no exigimos que sea conmutativo. La
      conmutatividad de la suma está asegurada para todo anillo, y la
      del producto está asegurada con la con\-dición de idempotencia
      impuesta a todos los elementos del anillo. La idempotencia de la
      suma también se deduce fácilmente de la idempotencia del producto.
      La suma lógica la establecemos
      \includegraphics[width=3.517cm,height=0.506cm]{./ObjectReplacements/Object 224}(como
      en 40.1, solo que aquí no su\-ponemos nada sobre álgebras de
      Boole), mientras que el producto lógico lo pone\-mos como idéntico
      al producto del anillo (idénticamente a lo anteriormente dicho).
      Nota: es importante darse cuenta que este anillo tendrá siempre
      divisores de cero, esto es habrá para cada elemento otro,
      distintos ambos de cero, que al multiplicarse dan cero, lo que
      impide que este anillo de Boole sea un dominio de integridad y por
      lo tanto también impide que sea cuerpo. El complementario,
      siguiendo la misma tó\-nica que en 40.1, se define
      como\includegraphics[width=1.598cm,height=0.467cm]{./ObjectReplacements/Object 225}.
      Por la idempotencia de la suma, la asociatividad de la suma y la
      existencia y unicidad de la suma tenemos que la doble negación es
      igual a la identidad. Sólo quedaría ver las distributividades.
      Cómo hasta ahora las demostraciones son cálculos que verifican la
      aserción. Así tenemos una forma de ir de cada álgebra de Boole a
      cada anillo de Boole, y un camino (exacta\-mente el inverso), que
      nos llevaría de cada anillo de Boole a cada álgebra de Boole. Las
      álgebras de Boole y los anillos de Boole son categorías
      equivalentes. A los ca\-minos los llamamos funtores. Daré solo un
      comienzo de estos teoremas:

      \begin{enumerate}
      \def\labelenumiv{\arabic{enumiv}.}
      \item
        {[}Axioma R0{]}
        \includegraphics[width=5.964cm,height=0.573cm]{./ObjectReplacements/Object 656}
      \item
        {[}Axioma R1{]}
        \includegraphics[width=6.274cm,height=0.573cm]{./ObjectReplacements/Object 655}
      \item
        {[}Axioma R2{]}
        \includegraphics[width=4.429cm,height=0.563cm]{./ObjectReplacements/Object 654}
      \item
        {[}Axioma R3{]}
        \includegraphics[width=4.784cm,height=0.61cm]{./ObjectReplacements/Object 653}
      \item
        {[}Axioma R4{]}
        \includegraphics[width=4.955cm,height=0.563cm]{./ObjectReplacements/Object 699}
      \item
        {[}Axioma R5{]}
        \includegraphics[width=5.339cm,height=0.573cm]{./ObjectReplacements/Object 700}
      \item
        {[}Axioma R6{]}
        \includegraphics[width=4.173cm,height=0.563cm]{./ObjectReplacements/Object 701}
      \item
        {[}Axioma R7{]}
        \includegraphics[width=4.494cm,height=0.563cm]{./ObjectReplacements/Object 765}
      \item
        {[}Axioma R8{]}
        \includegraphics[width=2.778cm,height=0.573cm]{./ObjectReplacements/Object 702}
      \item
        {[}Axioma R9{]}
        \includegraphics[width=6.392cm,height=0.573cm]{./ObjectReplacements/Object 703}
      \item
        {[}Axioma
        R10{]}\includegraphics[width=6.419cm,height=0.573cm]{./ObjectReplacements/Object 704}
      \item
        {[}Axioma
        BR{]}\includegraphics[width=2.863cm,height=0.563cm]{./ObjectReplacements/Object 759}
      \item
        \includegraphics[width=11.095cm,height=0.669cm]{./ObjectReplacements/Object 648}

        \begin{enumerate}
        \def\labelenumv{\arabic{enumv}.}
        \tightlist
        \item
          Prueba:
        \item
          \includegraphics[width=3.464cm,height=0.589cm]{./ObjectReplacements/Object 657}
        \item
          \includegraphics[width=4.113cm,height=0.721cm]{./ObjectReplacements/Object 619}
        \item
          \includegraphics[width=4.113cm,height=0.721cm]{./ObjectReplacements/Object 639}
        \item
          \includegraphics[width=3.039cm,height=0.653cm]{./ObjectReplacements/Object 640}
        \item
          \includegraphics[width=1.826cm,height=0.589cm]{./ObjectReplacements/Object 641}
        \end{enumerate}
      \item
        \includegraphics[width=7.895cm,height=0.639cm]{./ObjectReplacements/Object 642}

        \begin{enumerate}
        \def\labelenumv{\arabic{enumv}.}
        \tightlist
        \item
          Prueba:
        \item
          \includegraphics[width=7.151cm,height=0.647cm]{./ObjectReplacements/Object 643}
        \item
          \includegraphics[width=6.798cm,height=0.66cm]{./ObjectReplacements/Object 644}
        \item
          \includegraphics[width=6.193cm,height=0.653cm]{./ObjectReplacements/Object 645}
        \item
          \includegraphics[width=4.283cm,height=0.596cm]{./ObjectReplacements/Object 646}
        \item
          \includegraphics[width=2.048cm,height=0.589cm]{./ObjectReplacements/Object 647}
        \end{enumerate}
      \item
        \includegraphics[width=4.606cm,height=0.572cm]{./ObjectReplacements/Object 649}

        \begin{enumerate}
        \def\labelenumv{\arabic{enumv}.}
        \tightlist
        \item
          Prueba:
        \item
          \includegraphics[width=2.856cm,height=0.563cm]{./ObjectReplacements/Object 651}
        \item
          \includegraphics[width=5.777cm,height=0.639cm]{./ObjectReplacements/Object 652}
        \item
          \includegraphics[width=5.777cm,height=0.639cm]{./ObjectReplacements/Object 658}
        \item
          \includegraphics[width=3.471cm,height=0.563cm]{./ObjectReplacements/Object 659}
        \end{enumerate}
      \item
        \includegraphics[width=7.782cm,height=0.624cm]{./ObjectReplacements/Object 650}

        \begin{enumerate}
        \def\labelenumv{\arabic{enumv}.}
        \tightlist
        \item
          Prueba:
        \item
          \includegraphics[width=3.073cm,height=0.563cm]{./ObjectReplacements/Object 661}
        \item
          \includegraphics[width=4.784cm,height=0.61cm]{./ObjectReplacements/Object 662}
        \item
          \includegraphics[width=6.339cm,height=0.61cm]{./ObjectReplacements/Object 663}
        \item
          \includegraphics[width=5.579cm,height=0.61cm]{./ObjectReplacements/Object 668}
        \item
          \includegraphics[width=6.93cm,height=0.639cm]{./ObjectReplacements/Object 664}
        \item
          \includegraphics[width=4.073cm,height=0.639cm]{./ObjectReplacements/Object 665}
        \item
          \includegraphics[width=3.052cm,height=0.563cm]{./ObjectReplacements/Object 666}
        \item
          \includegraphics[width=2.461cm,height=0.563cm]{./ObjectReplacements/Object 667}
        \end{enumerate}
      \item
        \includegraphics[width=10.62cm,height=0.637cm]{./ObjectReplacements/Object 660}

        \begin{enumerate}
        \def\labelenumv{\arabic{enumv}.}
        \tightlist
        \item
          Prueba:
        \item
          \includegraphics[width=3.455cm,height=0.589cm]{./ObjectReplacements/Object 670}
        \item
          \includegraphics[width=2.872cm,height=0.589cm]{./ObjectReplacements/Object 671}
        \item
          \includegraphics[width=1.843cm,height=0.531cm]{./ObjectReplacements/Object 672}
        \item
          \includegraphics[width=1.252cm,height=0.531cm]{./ObjectReplacements/Object 673}
        \end{enumerate}
      \item
        \includegraphics[width=7.811cm,height=0.556cm]{./ObjectReplacements/Object 686}

        \begin{enumerate}
        \def\labelenumv{\arabic{enumv}.}
        \tightlist
        \item
          Prueba:
        \item
          \includegraphics[width=7.183cm,height=0.506cm]{./ObjectReplacements/Object 1761}
        \item
          \includegraphics[width=4.576cm,height=0.577cm]{./ObjectReplacements/Object 1760}
        \end{enumerate}
      \item
        \includegraphics[width=7.8cm,height=0.572cm]{./ObjectReplacements/Object 675}

        \begin{enumerate}
        \def\labelenumv{\arabic{enumv}.}
        \tightlist
        \item
          Prueba:
        \item
          \includegraphics[width=3.604cm,height=0.531cm]{./ObjectReplacements/Object 676}
        \item
          \includegraphics[width=5.265cm,height=0.506cm]{./ObjectReplacements/Object 677}
        \item
          \includegraphics[width=5.166cm,height=0.563cm]{./ObjectReplacements/Object 678}
        \item
          \includegraphics[width=2.208cm,height=0.467cm]{./ObjectReplacements/Object 679}
        \item
          \includegraphics[width=1.025cm,height=0.467cm]{./ObjectReplacements/Object 674}
        \end{enumerate}
      \item
        \includegraphics[width=7.638cm,height=0.572cm]{./ObjectReplacements/Object 681}

        \begin{enumerate}
        \def\labelenumv{\arabic{enumv}.}
        \tightlist
        \item
          \begin{enumerate}
          \def\labelenumvi{\arabic{enumvi}.}
          \tightlist
          \item
            \begin{enumerate}
            \def\labelenumvii{\arabic{enumvii}.}
            \tightlist
            \item
              Prueba:
            \item
              \includegraphics[width=2.228cm,height=0.467cm]{./ObjectReplacements/Object 682}
            \item
              \includegraphics[width=5.265cm,height=0.506cm]{./ObjectReplacements/Object 683}
            \item
              \includegraphics[width=5.166cm,height=0.563cm]{./ObjectReplacements/Object 684}
            \item
              \includegraphics[width=2.208cm,height=0.467cm]{./ObjectReplacements/Object 685}
            \item
              \includegraphics[width=1.025cm,height=0.467cm]{./ObjectReplacements/Object 680}
            \end{enumerate}
          \end{enumerate}
        \end{enumerate}
      \item
        \includegraphics[width=6.712cm,height=0.572cm]{./ObjectReplacements/Object 687}

        \begin{enumerate}
        \def\labelenumv{\arabic{enumv}.}
        \tightlist
        \item
          Prueba:
        \item
          \includegraphics[width=2.348cm,height=0.519cm]{./ObjectReplacements/Object 689}
        \item
          \includegraphics[width=3.655cm,height=0.575cm]{./ObjectReplacements/Object 690}
        \item
          \includegraphics[width=3.649cm,height=0.575cm]{./ObjectReplacements/Object 691}
        \item
          \includegraphics[width=2.217cm,height=0.467cm]{./ObjectReplacements/Object 692}
        \item
          \includegraphics[width=1.626cm,height=0.467cm]{./ObjectReplacements/Object 688}
        \end{enumerate}
      \item
        \includegraphics[width=6.712cm,height=0.572cm]{./ObjectReplacements/Object 694}

        \begin{enumerate}
        \def\labelenumv{\arabic{enumv}.}
        \tightlist
        \item
          Prueba:
        \item
          \includegraphics[width=2.351cm,height=0.519cm]{./ObjectReplacements/Object 695}
        \item
          \includegraphics[width=3.649cm,height=0.575cm]{./ObjectReplacements/Object 696}
        \item
          \includegraphics[width=3.655cm,height=0.575cm]{./ObjectReplacements/Object 697}
        \item
          \includegraphics[width=2.217cm,height=0.467cm]{./ObjectReplacements/Object 698}
        \item
          \includegraphics[width=1.626cm,height=0.467cm]{./ObjectReplacements/Object 693}
        \end{enumerate}
      \item
        \includegraphics[width=5.36cm,height=0.601cm]{./ObjectReplacements/Object 714}

        \begin{enumerate}
        \def\labelenumv{\arabic{enumv}.}
        \tightlist
        \item
          Prueba:
        \item
          \includegraphics[width=2.171cm,height=0.519cm]{./ObjectReplacements/Object 716}
        \item
          \includegraphics[width=2.175cm,height=0.519cm]{./ObjectReplacements/Object 717}
        \item
          \includegraphics[width=2.171cm,height=0.519cm]{./ObjectReplacements/Object 715}
        \item
          \includegraphics[width=2.159cm,height=0.575cm]{./ObjectReplacements/Object 718}
        \end{enumerate}
      \item
        \includegraphics[width=3.267cm,height=0.519cm]{./ObjectReplacements/Object 669}

        \begin{enumerate}
        \def\labelenumv{\arabic{enumv}.}
        \tightlist
        \item
          Prueba:
        \item
          \includegraphics[width=2.155cm,height=0.519cm]{./ObjectReplacements/Object 706}
        \item
          \includegraphics[width=2.722cm,height=0.519cm]{./ObjectReplacements/Object 707}
        \item
          \includegraphics[width=1.563cm,height=0.519cm]{./ObjectReplacements/Object 708}
        \end{enumerate}
      \item
        \includegraphics[width=4.565cm,height=0.572cm]{./ObjectReplacements/Object 720}

        \begin{enumerate}
        \def\labelenumv{\arabic{enumv}.}
        \tightlist
        \item
          Prueba:
        \item
          \includegraphics[width=1.99cm,height=0.519cm]{./ObjectReplacements/Object 721}
        \item
          \includegraphics[width=2.933cm,height=0.519cm]{./ObjectReplacements/Object 722}
        \item
          \includegraphics[width=2.321cm,height=0.519cm]{./ObjectReplacements/Object 723}
        \item
          \includegraphics[width=1.485cm,height=0.467cm]{./ObjectReplacements/Object 724}
        \item
          \includegraphics[width=0.785cm,height=0.467cm]{./ObjectReplacements/Object 725}
        \item
          \includegraphics[width=3.395cm,height=0.519cm]{./ObjectReplacements/Object 726}
        \item
          \includegraphics[width=1.612cm,height=0.519cm]{./ObjectReplacements/Object 727}
        \end{enumerate}
      \item
        \includegraphics[width=4.57cm,height=0.572cm]{./ObjectReplacements/Object 719}

        \begin{enumerate}
        \def\labelenumv{\arabic{enumv}.}
        \tightlist
        \item
          Prueba:
        \item
          \includegraphics[width=1.993cm,height=0.519cm]{./ObjectReplacements/Object 729}
        \item
          \includegraphics[width=2.933cm,height=0.519cm]{./ObjectReplacements/Object 730}
        \item
          \includegraphics[width=2.321cm,height=0.519cm]{./ObjectReplacements/Object 731}
        \item
          \includegraphics[width=1.485cm,height=0.467cm]{./ObjectReplacements/Object 732}
        \item
          \includegraphics[width=0.785cm,height=0.467cm]{./ObjectReplacements/Object 733}
        \item
          \includegraphics[width=3.39cm,height=0.519cm]{./ObjectReplacements/Object 734}
        \item
          \includegraphics[width=1.612cm,height=0.519cm]{./ObjectReplacements/Object 728}
        \end{enumerate}
      \item
        \includegraphics[width=6.881cm,height=0.601cm]{./ObjectReplacements/Object 741}

        \begin{enumerate}
        \def\labelenumv{\arabic{enumv}.}
        \tightlist
        \item
          Prueba:
        \item
          \includegraphics[width=3.258cm,height=0.575cm]{./ObjectReplacements/Object 742}
        \item
          \includegraphics[width=2.859cm,height=0.575cm]{./ObjectReplacements/Object 751}
        \item
          \includegraphics[width=1.529cm,height=0.467cm]{./ObjectReplacements/Object 744}
        \item
          \includegraphics[width=0.774cm,height=0.467cm]{./ObjectReplacements/Object 745}
        \end{enumerate}
      \item
        \includegraphics[width=6.876cm,height=0.601cm]{./ObjectReplacements/Object 746}

        \begin{enumerate}
        \def\labelenumv{\arabic{enumv}.}
        \tightlist
        \item
          Prueba:
        \item
          \includegraphics[width=3.258cm,height=0.575cm]{./ObjectReplacements/Object 747}
        \item
          \includegraphics[width=2.995cm,height=0.575cm]{./ObjectReplacements/Object 748}
        \item
          \includegraphics[width=1.496cm,height=0.467cm]{./ObjectReplacements/Object 749}
        \item
          \includegraphics[width=0.774cm,height=0.467cm]{./ObjectReplacements/Object 750}
        \end{enumerate}
      \item
        \includegraphics[width=6.676cm,height=0.573cm]{./ObjectReplacements/Object 743}

        \begin{enumerate}
        \def\labelenumv{\arabic{enumv}.}
        \tightlist
        \item
          Prueba:
        \item
          \includegraphics[width=4.399cm,height=0.575cm]{./ObjectReplacements/Object 752}
        \item
          \includegraphics[width=4.745cm,height=0.575cm]{./ObjectReplacements/Object 754}
        \item
          \includegraphics[width=3.558cm,height=0.575cm]{./ObjectReplacements/Object 753}
        \item
          \includegraphics[width=2.314cm,height=0.575cm]{./ObjectReplacements/Object 756}
        \item
          \includegraphics[width=1.692cm,height=0.519cm]{./ObjectReplacements/Object 757}
        \item
          \includegraphics[width=0.774cm,height=0.467cm]{./ObjectReplacements/Object 758}
        \end{enumerate}
      \item
        \includegraphics[width=5.703cm,height=0.573cm]{./ObjectReplacements/Object 705}

        \begin{enumerate}
        \def\labelenumv{\arabic{enumv}.}
        \tightlist
        \item
          Prueba:
        \item
          \includegraphics[width=2.559cm,height=0.575cm]{./ObjectReplacements/Object 709}
        \item
          \includegraphics[width=3.498cm,height=0.575cm]{./ObjectReplacements/Object 710}
        \item
          \includegraphics[width=2.886cm,height=0.575cm]{./ObjectReplacements/Object 711}
        \item
          \includegraphics[width=1.496cm,height=0.467cm]{./ObjectReplacements/Object 712}
        \item
          \includegraphics[width=0.774cm,height=0.467cm]{./ObjectReplacements/Object 713}
        \end{enumerate}
      \item
        \includegraphics[width=4.341cm,height=0.572cm]{./ObjectReplacements/Object 735}

        \begin{enumerate}
        \def\labelenumv{\arabic{enumv}.}
        \tightlist
        \item
          Prueba:
        \item
          \includegraphics[width=6.817cm,height=0.573cm]{./ObjectReplacements/Object 737}
        \item
          \includegraphics[width=6.495cm,height=0.573cm]{./ObjectReplacements/Object 738}
        \item
          \includegraphics[width=2.205cm,height=0.467cm]{./ObjectReplacements/Object 739}
        \end{enumerate}
      \item
        \includegraphics[width=4.568cm,height=0.572cm]{./ObjectReplacements/Object 760}
      \item
        \includegraphics[width=10.278cm,height=0.482cm]{./ObjectReplacements/Object 793}
      \item
        \includegraphics[width=5.754cm,height=0.573cm]{./ObjectReplacements/Object 736}

        \begin{enumerate}
        \def\labelenumv{\arabic{enumv}.}
        \tightlist
        \item
          Prueba:
        \item
          \includegraphics[width=1.695cm,height=0.519cm]{./ObjectReplacements/Object 740}
        \item
          \includegraphics[width=2.304cm,height=0.575cm]{./ObjectReplacements/Object 766}
        \item
          \includegraphics[width=1.351cm,height=0.519cm]{./ObjectReplacements/Object 767}
        \end{enumerate}
      \item
        \includegraphics[width=8.066cm,height=0.693cm]{./ObjectReplacements/Object 755}

        \begin{enumerate}
        \def\labelenumv{\arabic{enumv}.}
        \tightlist
        \item
          Prueba:
        \item
          \includegraphics[width=4.847cm,height=0.693cm]{./ObjectReplacements/Object 768}
        \item
          \includegraphics[width=2.939cm,height=0.519cm]{./ObjectReplacements/Object 769}
        \item
          \includegraphics[width=1.316cm,height=0.467cm]{./ObjectReplacements/Object 770}
        \item
          \includegraphics[width=0.97cm,height=0.467cm]{./ObjectReplacements/Object 771}
        \end{enumerate}
      \item
        \includegraphics[width=8.066cm,height=0.693cm]{./ObjectReplacements/Object 773}

        \begin{enumerate}
        \def\labelenumv{\arabic{enumv}.}
        \tightlist
        \item
          Prueba:
        \item
          \includegraphics[width=4.847cm,height=0.693cm]{./ObjectReplacements/Object 774}
        \item
          \includegraphics[width=1.535cm,height=0.519cm]{./ObjectReplacements/Object 775}
        \item
          \includegraphics[width=1.321cm,height=0.467cm]{./ObjectReplacements/Object 1762}
        \item
          \includegraphics[width=1.669cm,height=0.467cm]{./ObjectReplacements/Object 1763}
        \item
          \includegraphics[width=1.316cm,height=0.467cm]{./ObjectReplacements/Object 776}
        \item
          \includegraphics[width=0.97cm,height=0.467cm]{./ObjectReplacements/Object 772}
        \end{enumerate}
      \item
        \includegraphics[width=7.96cm,height=0.601cm]{./ObjectReplacements/Object 867}

        \begin{enumerate}
        \def\labelenumv{\arabic{enumv}.}
        \tightlist
        \item
          Prueba:
        \item
          \includegraphics[width=4.29cm,height=0.575cm]{./ObjectReplacements/Object 877}
        \item
          \includegraphics[width=4.949cm,height=0.704cm]{./ObjectReplacements/Object 878}
        \item
          \includegraphics[width=3.42cm,height=0.637cm]{./ObjectReplacements/Object 879}
        \item
          \includegraphics[width=2.559cm,height=0.575cm]{./ObjectReplacements/Object 880}
        \item
          \includegraphics[width=0.774cm,height=0.467cm]{./ObjectReplacements/Object 881}
        \end{enumerate}
      \end{enumerate}

      Esta fórmula que acabamos de exponer es la fórmula universal para
      el inverso en cualquier grupo, o in\-cluso, para cualquier
      operación con neutro asociativa, siempre que existan los inversos,
      tanto el total como los individuales.

      \begin{enumerate}
      \def\labelenumiv{\arabic{enumiv}.}
      \setcounter{enumiv}{37}
      \item
        \includegraphics[width=6.308cm,height=0.563cm]{./ObjectReplacements/Object 761}

        \begin{enumerate}
        \def\labelenumv{\arabic{enumv}.}
        \tightlist
        \item
          Prueba:
        \item
          \includegraphics[width=2.247cm,height=0.575cm]{./ObjectReplacements/Object 762}
        \item
          \includegraphics[width=2.928cm,height=0.519cm]{./ObjectReplacements/Object 763}
        \item
          \includegraphics[width=4.106cm,height=0.575cm]{./ObjectReplacements/Object 764}
        \item
          \includegraphics[width=2.551cm,height=0.519cm]{./ObjectReplacements/Object 777}
        \item
          \includegraphics[width=2.245cm,height=0.575cm]{./ObjectReplacements/Object 778}
        \item
          \includegraphics[width=3.441cm,height=0.519cm]{./ObjectReplacements/Object 788}
        \item
          \includegraphics[width=6.431cm,height=0.575cm]{./ObjectReplacements/Object 789}
        \item
          \includegraphics[width=2.288cm,height=0.467cm]{./ObjectReplacements/Object 779}
        \end{enumerate}
      \end{enumerate}

      El grupo aditivo de un anillo con unidad multiplicativa por ambos
      lados es siempre un grupo abeliano.

      \begin{enumerate}
      \def\labelenumiv{\arabic{enumiv}.}
      \setcounter{enumiv}{38}
      \item
        \includegraphics[width=5.166cm,height=0.563cm]{./ObjectReplacements/Object 780}

        \begin{enumerate}
        \def\labelenumv{\arabic{enumv}.}
        \tightlist
        \item
          Prueba:
        \item
          \includegraphics[width=2.175cm,height=0.519cm]{./ObjectReplacements/Object 781}
        \item
          \includegraphics[width=3.364cm,height=0.575cm]{./ObjectReplacements/Object 782}
        \item
          \includegraphics[width=2.223cm,height=0.519cm]{./ObjectReplacements/Object 783}
        \item
          \includegraphics[width=1.609cm,height=0.467cm]{./ObjectReplacements/Object 784}
        \end{enumerate}
      \end{enumerate}

      En el grupo aditivo de un anillo booleano
      \includegraphics[width=0.51cm,height=0.563cm]{./ObjectReplacements/Object 786}
      es siempre
      \includegraphics[width=1.582cm,height=0.519cm]{./ObjectReplacements/Object 785}.

      \begin{enumerate}
      \def\labelenumiv{\arabic{enumiv}.}
      \setcounter{enumiv}{39}
      \item
        \includegraphics[width=5.387cm,height=0.563cm]{./ObjectReplacements/Object 799}
      \item
        \includegraphics[width=4.937cm,height=0.563cm]{./ObjectReplacements/Object 787}

        \begin{enumerate}
        \def\labelenumv{\arabic{enumv}.}
        \tightlist
        \item
          Prueba:
        \item
          \includegraphics[width=1.154cm,height=0.467cm]{./ObjectReplacements/Object 802}
        \item
          \includegraphics[width=2.332cm,height=0.519cm]{./ObjectReplacements/Object 803}
        \item
          \includegraphics[width=2.215cm,height=0.527cm]{./ObjectReplacements/Object 804}
        \item
          \includegraphics[width=1.739cm,height=0.467cm]{./ObjectReplacements/Object 805}
        \item
          \includegraphics[width=0.774cm,height=0.467cm]{./ObjectReplacements/Object 806}
        \end{enumerate}
      \item
        \includegraphics[width=7.832cm,height=0.573cm]{./ObjectReplacements/Object 800}
      \item
        \includegraphics[width=5.846cm,height=0.563cm]{./ObjectReplacements/Object 790}

        \begin{enumerate}
        \def\labelenumv{\arabic{enumv}.}
        \tightlist
        \item
          Prueba:
        \item
          \includegraphics[width=8.243cm,height=0.57cm]{./ObjectReplacements/Object 791}
        \item
          \includegraphics[width=3.113cm,height=0.57cm]{./ObjectReplacements/Object 792}
        \item
          \includegraphics[width=2.441cm,height=0.482cm]{./ObjectReplacements/Object 794}
        \item
          \includegraphics[width=3.976cm,height=0.467cm]{./ObjectReplacements/Object 795}
        \item
          \includegraphics[width=2.332cm,height=0.467cm]{./ObjectReplacements/Object 796}
        \item
          \includegraphics[width=2.14cm,height=0.467cm]{./ObjectReplacements/Object 797}
        \item
          \includegraphics[width=1.82cm,height=0.467cm]{./ObjectReplacements/Object 798}
        \end{enumerate}
      \end{enumerate}

      La operación multiplicativa de un anillo de Boole
      \includegraphics[width=0.51cm,height=0.563cm]{./ObjectReplacements/Object 801}
      es siempre abeliana. Un anillo de Boole es una subcategoría de la
      categoría de los anillos conmutativos.

      \begin{enumerate}
      \def\labelenumiv{\arabic{enumiv}.}
      \setcounter{enumiv}{43}
      \item
        \includegraphics[width=4.958cm,height=0.563cm]{./ObjectReplacements/Object 808}

        \begin{enumerate}
        \def\labelenumv{\arabic{enumv}.}
        \tightlist
        \item
          Prueba:
        \item
          \includegraphics[width=1.154cm,height=0.467cm]{./ObjectReplacements/Object 809}
        \item
          \includegraphics[width=2.332cm,height=0.519cm]{./ObjectReplacements/Object 810}
        \item
          \includegraphics[width=2.215cm,height=0.527cm]{./ObjectReplacements/Object 811}
        \item
          \includegraphics[width=1.739cm,height=0.467cm]{./ObjectReplacements/Object 812}
        \item
          \includegraphics[width=0.774cm,height=0.467cm]{./ObjectReplacements/Object 813}
        \end{enumerate}
      \item
        \includegraphics[width=5.159cm,height=0.563cm]{./ObjectReplacements/Object 816}

        \begin{enumerate}
        \def\labelenumv{\arabic{enumv}.}
        \tightlist
        \item
          Prueba:
        \item
          \includegraphics[width=1.388cm,height=0.467cm]{./ObjectReplacements/Object 817}
        \item
          \[{} = {x\oplus\left( {x\oplus 1} \right)} = {}\]
        \item
          \includegraphics[width=2.565cm,height=0.519cm]{./ObjectReplacements/Object 819}
        \item
          \includegraphics[width=1.709cm,height=0.467cm]{./ObjectReplacements/Object 820}
        \item
          \includegraphics[width=0.764cm,height=0.467cm]{./ObjectReplacements/Object 815}
        \end{enumerate}
      \item
        \includegraphics[width=5.085cm,height=0.563cm]{./ObjectReplacements/Object 822}

        \begin{enumerate}
        \def\labelenumv{\arabic{enumv}.}
        \tightlist
        \item
          Prueba:
        \item
          \includegraphics[width=1.309cm,height=0.467cm]{./ObjectReplacements/Object 823}
        \item
          \includegraphics[width=3.212cm,height=0.575cm]{./ObjectReplacements/Object 825}
        \item
          \includegraphics[width=1.709cm,height=0.467cm]{./ObjectReplacements/Object 828}
        \item
          \includegraphics[width=0.764cm,height=0.467cm]{./ObjectReplacements/Object 821}
        \end{enumerate}
      \item
        \includegraphics[width=5.105cm,height=0.563cm]{./ObjectReplacements/Object 827}

        \begin{enumerate}
        \def\labelenumv{\arabic{enumv}.}
        \tightlist
        \item
          Prueba:
        \item
          \includegraphics[width=1.307cm,height=0.467cm]{./ObjectReplacements/Object 829}
        \item
          \includegraphics[width=3.21cm,height=0.575cm]{./ObjectReplacements/Object 831}
        \item
          \includegraphics[width=1.73cm,height=0.467cm]{./ObjectReplacements/Object 832}
        \item
          \includegraphics[width=0.785cm,height=0.467cm]{./ObjectReplacements/Object 826}
        \end{enumerate}
      \item
        \includegraphics[width=5.096cm,height=0.563cm]{./ObjectReplacements/Object 830}

        \begin{enumerate}
        \def\labelenumv{\arabic{enumv}.}
        \tightlist
        \item
          Prueba:
        \item
          \includegraphics[width=1.307cm,height=0.467cm]{./ObjectReplacements/Object 833}
        \item
          \includegraphics[width=3.434cm,height=0.575cm]{./ObjectReplacements/Object 834}
        \item
          \includegraphics[width=1.73cm,height=0.467cm]{./ObjectReplacements/Object 835}
        \item
          \includegraphics[width=0.785cm,height=0.467cm]{./ObjectReplacements/Object 824}
        \end{enumerate}
      \item
        \includegraphics[width=6.16cm,height=0.563cm]{./ObjectReplacements/Object 836}

        \begin{enumerate}
        \def\labelenumv{\arabic{enumv}.}
        \tightlist
        \item
          Prueba:
        \item
          \includegraphics[width=1.341cm,height=0.467cm]{./ObjectReplacements/Object 837}
        \item
          \includegraphics[width=3.267cm,height=0.519cm]{./ObjectReplacements/Object 838}
        \item
          \includegraphics[width=3.267cm,height=0.519cm]{./ObjectReplacements/Object 839}
        \item
          \includegraphics[width=1.341cm,height=0.467cm]{./ObjectReplacements/Object 840}
        \end{enumerate}
      \item
        \includegraphics[width=8.25cm,height=0.573cm]{./ObjectReplacements/Object 842}

        \begin{enumerate}
        \def\labelenumv{\arabic{enumv}.}
        \tightlist
        \item
          Prueba:
        \item
          \includegraphics[width=2.909cm,height=0.506cm]{./ObjectReplacements/Object 846}
        \item
          \includegraphics[width=3.812cm,height=0.506cm]{./ObjectReplacements/Object 847}
        \item
          \includegraphics[width=7.982cm,height=0.506cm]{./ObjectReplacements/Object 848}
        \item
          \includegraphics[width=7.451cm,height=0.506cm]{./ObjectReplacements/Object 849}
        \item
          \includegraphics[width=2.872cm,height=0.506cm]{./ObjectReplacements/Object 850}
        \item
          \includegraphics[width=4.646cm,height=0.506cm]{./ObjectReplacements/Object 851}
        \item
          \includegraphics[width=7.347cm,height=0.506cm]{./ObjectReplacements/Object 852}
        \item
          \includegraphics[width=4.65cm,height=0.482cm]{./ObjectReplacements/Object 853}
        \item
          \includegraphics[width=7.451cm,height=0.506cm]{./ObjectReplacements/Object 854}
        \item
          \includegraphics[width=4.727cm,height=0.563cm]{./ObjectReplacements/Object 855}
        \end{enumerate}
      \item
        \includegraphics[width=8.661cm,height=0.577cm]{./ObjectReplacements/Object 843}

        \begin{enumerate}
        \def\labelenumv{\arabic{enumv}.}
        \tightlist
        \item
          Prueba:
        \item
          \includegraphics[width=2.755cm,height=0.506cm]{./ObjectReplacements/Object 845}
        \item
          \includegraphics[width=3.658cm,height=0.506cm]{./ObjectReplacements/Object 856}
        \item
          \includegraphics[width=4.293cm,height=0.506cm]{./ObjectReplacements/Object 857}
        \item
          \includegraphics[width=3.351cm,height=0.506cm]{./ObjectReplacements/Object 858}
        \item
          \includegraphics[width=4.293cm,height=0.506cm]{./ObjectReplacements/Object 859}
        \item
          \includegraphics[width=7.731cm,height=0.482cm]{./ObjectReplacements/Object 860}
        \end{enumerate}
      \item
        \includegraphics[width=8.83cm,height=0.577cm]{./ObjectReplacements/Object 844}

        \begin{enumerate}
        \def\labelenumv{\arabic{enumv}.}
        \tightlist
        \item
          Prueba:
        \item
          \includegraphics[width=2.755cm,height=0.506cm]{./ObjectReplacements/Object 861}
        \item
          \includegraphics[width=2.328cm,height=0.506cm]{./ObjectReplacements/Object 863}
        \item
          \includegraphics[width=4.63cm,height=0.506cm]{./ObjectReplacements/Object 864}
        \item
          \includegraphics[width=3.706cm,height=0.506cm]{./ObjectReplacements/Object 865}
        \item
          \includegraphics[width=3.505cm,height=0.506cm]{./ObjectReplacements/Object 866}
        \item
          \includegraphics[width=6.366cm,height=0.506cm]{./ObjectReplacements/Object 862}
        \item
          \includegraphics[width=12.575cm,height=0.506cm]{./ObjectReplacements/Object 868}
        \item
          \includegraphics[width=11.155cm,height=0.506cm]{./ObjectReplacements/Object 869}
        \item
          \includegraphics[width=13.995cm,height=0.506cm]{./ObjectReplacements/Object 870}
        \item
          \includegraphics[width=11.367cm,height=0.467cm]{./ObjectReplacements/Object 871}
        \item
          \includegraphics[width=9.998cm,height=0.467cm]{./ObjectReplacements/Object 872}
        \item
          \includegraphics[width=8.11cm,height=0.467cm]{./ObjectReplacements/Object 873}
        \item
          \includegraphics[width=5.419cm,height=0.467cm]{./ObjectReplacements/Object 876}
        \item
          \includegraphics[width=3.463cm,height=0.467cm]{./ObjectReplacements/Object 874}
        \item
          \includegraphics[width=3.6cm,height=0.519cm]{./ObjectReplacements/Object 875}
        \item
          \includegraphics[width=7.731cm,height=0.482cm]{./ObjectReplacements/Object 841}
        \end{enumerate}
      \end{enumerate}
    \end{enumerate}
  \item
    Estudio de las funciones booleanas, sobre álgebras de Boole finitas.

    \begin{enumerate}
    \def\labelenumiii{\arabic{enumiii}.}
    \tightlist
    \item
      Estudiaremos sólo las
      funciones\includegraphics[width=1.729cm,height=0.563cm]{./ObjectReplacements/Object 65}dónde\includegraphics[width=2.514cm,height=0.61cm]{./ObjectReplacements/Object 66}.
      Esto
      es\includegraphics[width=1.984cm,height=0.61cm]{./ObjectReplacements/Object 67}.
    \item
      Formas normales. Toda
      función\includegraphics[width=1.984cm,height=0.61cm]{./ObjectReplacements/Object 80}se
      puede poner en la
      forma\includegraphics[width=6.232cm,height=1.177cm]{./ObjectReplacements/Object 81}dónde
      \includegraphics[width=3.175cm,height=0.506cm]{./ObjectReplacements/Object 1313}.
    \item
      En general, el número de funciones de un
      conjunto\includegraphics[width=0.506cm,height=0.467cm]{./ObjectReplacements/Object 29}de
      cardinal\includegraphics[width=1.272cm,height=0.531cm]{./ObjectReplacements/Object 30}en
      un
      conjunto\includegraphics[width=0.547cm,height=0.467cm]{./ObjectReplacements/Object 31}de
      cardinal\includegraphics[width=1.295cm,height=0.531cm]{./ObjectReplacements/Object 89}será\includegraphics[width=0.87cm,height=0.598cm]{./ObjectReplacements/Object 90}.
      Así una función
      de\includegraphics[width=0.422cm,height=0.467cm]{./ObjectReplacements/Object 95}variables
      en
      \includegraphics[width=0.506cm,height=0.467cm]{./ObjectReplacements/Object 91}con
      valores en
      \includegraphics[width=0.547cm,height=0.467cm]{./ObjectReplacements/Object 92}será\includegraphics[width=1.131cm,height=0.633cm]{./ObjectReplacements/Object 93}.
      Para el caso en
      que\includegraphics[width=1.993cm,height=0.61cm]{./ObjectReplacements/Object 94}que
      toma como
      argumento\includegraphics[width=0.422cm,height=0.467cm]{./ObjectReplacements/Object 96}
      variables, obtenemos que el número de funciones será
      de\includegraphics[width=0.944cm,height=0.676cm]{./ObjectReplacements/Object 97}.
      En una pequeña tabla vemos como crece esta cantidad:
    \end{enumerate}
  \end{enumerate}
\end{enumerate}

{\def\LTcaptype{none} % do not increment counter
\begin{longtable}[]{@{}ll@{}}
\toprule\noalign{}
\endhead
\bottomrule\noalign{}
\endlastfoot
\includegraphics[width=0.564cm,height=0.467cm]{./ObjectReplacements/Object 99}
variables & Número de funciones distintas:
\includegraphics[width=0.967cm,height=0.656cm]{./ObjectReplacements/Object 98} \\
0 & 2 \\
1 & 4 \\
2 & 16 \\
3 & 256 \\
4 & 65536 \\
5 & 4294967296 \\
6 & 18446744073709551616 \\
7 & 340282366920938463463374607431768211456 \\
\end{longtable}
}

\begin{enumerate}
\def\labelenumi{\arabic{enumi}.}
\setcounter{enumi}{38}
\item
  \begin{enumerate}
  \def\labelenumii{\arabic{enumii}.}
  \item
    \begin{enumerate}
    \def\labelenumiii{\arabic{enumiii}.}
    \item
      El punto anterior es fácil de probar:

      \begin{enumerate}
      \def\labelenumiv{\arabic{enumiv}.}
      \item
        Para el caso
        de\includegraphics[width=1.122cm,height=0.467cm]{./ObjectReplacements/Object 100}y
        de\includegraphics[width=1.109cm,height=0.467cm]{./ObjectReplacements/Object 101}es
        fácil probar (por enumeración) la vali\-dez de la fórmula. Más
        tarde mostraremos tablas de todas las funciones hasta
        \includegraphics[width=1.111cm,height=0.467cm]{./ObjectReplacements/Object 102}inclusive.
      \item
        Para el caso general, la Hipótesis de Inducción será
        :\includegraphics[width=12.017cm,height=0.686cm]{./ObjectReplacements/Object 103}
      \item
        Veremos si para el
        caso\includegraphics[width=1.12cm,height=0.467cm]{./ObjectReplacements/Object 104}se
        sigue cumpliendo la fórmula anterior. Pero esto es claro: al
        añadir una variable en el argumento tendremos todas las
        funcio\-nes del
        caso\includegraphics[width=1.683cm,height=0.467cm]{./ObjectReplacements/Object 106}:\includegraphics[width=1.205cm,height=0.557cm]{./ObjectReplacements/Object 107}para
        el
        valor\includegraphics[width=0.422cm,height=0.467cm]{./ObjectReplacements/Object 108}de
        la nueva variable y
        otros\includegraphics[width=1.205cm,height=0.557cm]{./ObjectReplacements/Object 109}para
        el
        valor\includegraphics[width=0.411cm,height=0.467cm]{./ObjectReplacements/Object 110}de
        la nueva variable, y no quedan otros casos. Las funciones
        totales
        para\includegraphics[width=1.12cm,height=0.467cm]{./ObjectReplacements/Object 111}serán
        :

        \begin{enumerate}
        \def\labelenumv{\arabic{enumv}.}
        \tightlist
        \item
          \includegraphics[width=12.033cm,height=0.7cm]{./ObjectReplacements/Object 440}
        \item
          \includegraphics[width=13.714cm,height=0.7cm]{./ObjectReplacements/Object 441}
        \item
          \includegraphics[width=9.216cm,height=0.686cm]{./ObjectReplacements/Object 442}
        \end{enumerate}
      \end{enumerate}
    \end{enumerate}
  \end{enumerate}
\end{enumerate}

Y así queda establecida la fórmula.

\begin{enumerate}
\def\labelenumi{\arabic{enumi}.}
\setcounter{enumi}{39}
\item
  \begin{enumerate}
  \def\labelenumii{\arabic{enumii}.}
  \item
    \begin{enumerate}
    \def\labelenumiii{\arabic{enumiii}.}
    \item
      Como se ve en el punto anterior, el crecimiento es desmesurado al
      compararlo al crecimiento lineal de los argumentos. En un futuro,
      cuan\-do intentemos hacer reducciones de expresiones booleanas,
      este crecimiento nos im\-pedirá construir métodos eficaces para
      resolver las minimizaciones.
    \item
      Aunque hemos visto que podemos poner las expresiones booleanas en
      los conjuntos de
      operadores\includegraphics[width=6.078cm,height=0.492cm]{./ObjectReplacements/Object 112},
      por comprensibilidad y para una lectura normal se utilizan
      frecuentemente los dos primeros conjuntos unidos, pu\-diendo
      variarse bien el orden de los operadores
      binarios:\includegraphics[width=3.551cm,height=0.506cm]{./ObjectReplacements/Object 113}.
      El primer conjunto
      \includegraphics[width=1.647cm,height=0.506cm]{./ObjectReplacements/Object 114}será
      el que estudiemos por defecto, el segundo se tratará con una
      simetría de dualidad (hay que tener algunos cuidados). El conjunto
      elegido de operaciones se llamará desarrollo en sumas de productos
      de términos simples (una variable, o su negada, o una constante).
      También se llamará desarrollo por minitérminos o SOP (inglés) o
      SdP. El segundo será el desarrollo en producto de sumas de
      términos simples (una variable, o su negada, o una constante).
      También se llamará desarrollo por maxitérminos o POS (inglés) o
      PdS.
    \item
      En el desarrollo por minitérminos expresamos sólo los términos de
      la expresión en que la función tiene como
      valor\includegraphics[width=0.411cm,height=0.467cm]{./ObjectReplacements/Object 115}.
      Veamos:
    \item
      \includegraphics[width=12.857cm,height=1.207cm]{./ObjectReplacements/Object 116}

      \includegraphics[width=6.817cm,height=0.603cm]{./ObjectReplacements/Object 226}
    \item
      Para más sencillez:

      \includegraphics[width=13.78cm,height=1.125cm]{./ObjectReplacements/Object 117}\includegraphics[width=10.149cm,height=0.54cm]{./ObjectReplacements/Object 227}
    \item
      Cuándo la sigma (el minitérmino) es en todos los casos (para todas
      las iotas) completo (es un producto
      de\includegraphics[width=3.708cm,height=0.467cm]{./ObjectReplacements/Object 807}),
      el valor de la función en esa iota concreta indica si el
      minitérmino aparece o no.
    \item
      En el desarrollo por maxitérminos de una función
      de\includegraphics[width=0.448cm,height=0.467cm]{./ObjectReplacements/Object 883}variables,
      los
      \includegraphics[width=2.612cm,height=0.467cm]{./ObjectReplacements/Object 882}son
      sumatorios
      de\includegraphics[width=3.708cm,height=0.467cm]{./ObjectReplacements/Object 814}.
      Así termina consis\-tiendo la función en un producto de
      maxitérminos. Los maxitérminos indican un
      \includegraphics[width=0.422cm,height=0.467cm]{./ObjectReplacements/Object 884}de
      la función.
    \item
      Además de expresar las funciones por cadenas de símbolos que
      constituyen un térmi\-no, existe una posibilidad de expresar estas
      funciones por tablas lineales o por cua\-dros (tablas
      bidimensionales). Para cada combinación de valores booleanos a la
      en\-trada de una función obtenemos un valor booleano de salida.
    \item
      Para que las funciones booleanas representen algo de interés para
      la ingeniería, lo primero que debemos tener es una forma de
      representar la información que queremos procesar. ¿Cómo
      representamos un número?. ¿Cómo una letra?. Haremos un alto en la
      exposición de funciones booleanas, para detallar más esta
      pregunta, poder respon\-derla y así ver para qué estamos viendo
      las álgebras de Boole.
    \end{enumerate}
  \end{enumerate}
\item
  Representación de la información.

  \begin{enumerate}
  \def\labelenumii{\arabic{enumii}.}
  \tightlist
  \item
    Un alfabeto es un conjunto de valores diferentes que podemos aplicar
    para represen\-tar información. En el sentido que nosotros lo
    utilizamos el alfabeto de la escritura en español no solo se compone
    de las letras del abecedario, digamos en minúsculas, sino además, de
    otro conjunto similar pero en mayúsculas. A esto hay que añadir
    todos los signos de puntuación en párrafos. Además hay que añadir
    todas las vocales que son susceptibles de tener tilde o diéresis,
    además, el guion para separar las palabras en dos y por último un
    elemento muy frecuente que suele pasar desapercibido: el espa\-cio
    en blanco. Por último los guarismos de los números del 0 al 9. Si
    consideramos los números naturales estos guarismos son un alfabeto
    de los números naturales. Como vemos los alfabetos tienen la común
    propiedad de ser finitos. Para los núme\-ros naturales (como para
    cualquier otra cosa que representar) basta con
    \includegraphics[width=1.984cm,height=0.681cm]{./ObjectReplacements/Object 201}
    .
  \item
    Para nosotros un lenguaje sobre un alfabeto será un conjunto, donde
    los elementos serán ristras de letras de ese alfabeto. Una ristra de
    letras (una ristra finita) de ese al\-fabeto no tiene porqué
    pertenecer al lenguaje. Por ejemplo, podemos establecer un lenguaje
    de los números naturales sobre un alfabeto cualquiera. Si es sobre
    el
    alfabe\-to\includegraphics[width=0.672cm,height=0.681cm]{./ObjectReplacements/Object 228},
    de cardinal 2, diremos que una palabra de nuestro lenguaje es bien
    un 0 o bien un 1 seguido de una ristra finita cualquiera de 0s y 1s
    (formalmente se suele re\-presentar
    como\includegraphics[width=3.055cm,height=0.593cm]{./ObjectReplacements/Object 229},
    dónde
    \textquotesingle{}\includegraphics[width=0.445cm,height=0.467cm]{./ObjectReplacements/Object 230}\textquotesingle{}
    quiere decir seguido o conca\-tenado, y
    \textquotesingle\textquotesingle{}\includegraphics[width=0.975cm,height=0.542cm]{./ObjectReplacements/Object 231}\textquotesingle\textquotesingle{}
    quiere decir una ristra de letras de dentro del paréntesis, con la
    única condición que sea finita, y que puede ser vacía).
  \item
    Una palabra del español podemos representarla sobre el alfabeto
    castellano en minús\-culas unido al alfabeto castellano en
    mayúsculas, unido a la tilde y la diéresis cómo el conjunto de
    palabras del diccionario de la Real Academia de la Lengua,
    represen\-tando una palabra con una diéresis como cigüeña como
    ``cig\#ueña'', y si lleva una til\-de como en Julián como
    ``Juli\textasciitilde an''. Un nombre más habitual que el de
    lenguaje en electrónica suele ser el de código.
  \item
    Antes de seguir el desarrollo, veamos algunos casos particulares de
    códigos: código binario natural de cualquier longitud, código
    binario natural de longitud fija, códigos 5 entre 2 (o biquinarios),
    códigos BCD (BCD natural, BCD exceso 3, BCD auto-complementario
    Aitken), códigos con distancias sucesivas 1 (para longitud fija)
    (códigos continuos), códigos continuos con distancia 1 entre el
    primer elemento y el último elemento o circulares. Códigos
    especulares. Si son circulares y especulares se llaman códigos Gray
    (de una determinada longitud fija). Códigos con redundancia de
    infor\-mación, entre los que destacan los códigos Hamming.
  \end{enumerate}
\end{enumerate}

Podéis ver que en los códigos antes explicitados y sus propiedades hacen
referencia necesaria a algo más que al conjunto de las palabras de un
lenguaje, que se puede re\-sumir en un orden en las palabras. Esto viene
dado normalmente por el significado de las palabras.

\begin{enumerate}
\def\labelenumi{\arabic{enumi}.}
\setcounter{enumi}{41}
\item
  \begin{enumerate}
  \def\labelenumii{\arabic{enumii}.}
  \item
    Hemos hablado de alfabetos, lenguajes sobre un alfabeto y después
    sobre el orden sobre esas palabras. Si llamamos a las reglas que
    cumplen las palabras de un lenguaje respecto a un alfabeto la
    gramática de ese lenguaje, el último elemento, el del orden, o más
    en general los significados, son la semántica.
  \item
    En 3 hablamos de un lenguaje de los números naturales sobre un
    alfabeto\includegraphics[width=0.672cm,height=0.681cm]{./ObjectReplacements/Object 232}.
    En general: ¿qué número significa 1001? ¿y el 1011101?. A qué valor
    natural apunta cada cadena es algo exterior al léxico y su
    gramática. Vamos a construir una semánti\-ca: llamamos en una cadena
    de letras del alfabeto
    \includegraphics[width=2.018cm,height=0.601cm]{./ObjectReplacements/Object 233}a
    \includegraphics[width=0.496cm,height=0.601cm]{./ObjectReplacements/Object 234}el
    dígito me\-nos significativo
    (el\includegraphics[width=0.734cm,height=0.467cm]{./ObjectReplacements/Object 351})
    y
    al\includegraphics[width=0.938cm,height=0.601cm]{./ObjectReplacements/Object 235}el
    dígito más significativo
    (el\includegraphics[width=0.965cm,height=0.467cm]{./ObjectReplacements/Object 885}).
    Pri\-mero asignamos unos valores naturales a los del
    alfabeto\includegraphics[width=4.036cm,height=1.348cm]{./ObjectReplacements/Object 236},
    y según el
    subíndice\includegraphics[width=0.448cm,height=0.467cm]{./ObjectReplacements/Object 238}de
    la
    letra\includegraphics[width=0.515cm,height=0.601cm]{./ObjectReplacements/Object 239}obtenemos
    el
    valor\includegraphics[width=1.935cm,height=0.7cm]{./ObjectReplacements/Object 237}para
    esa letra en ese lugar concreto. Así el valor de una palabra de
    nuestro código
    será\includegraphics[width=2.639cm,height=1.139cm]{./ObjectReplacements/Object 240}.
    Más en general si
    \includegraphics[width=0.711cm,height=0.617cm]{./ObjectReplacements/Object 269}es
    un alfabeto de
    cardinal\includegraphics[width=0.448cm,height=0.467cm]{./ObjectReplacements/Object 270},
    esto es, desde los dígitos
    \includegraphics[width=3.96cm,height=0.663cm]{./ObjectReplacements/Object 271}y
    establecemos la
    función\includegraphics[width=3.542cm,height=0.617cm]{./ObjectReplacements/Object 272},
    el valor de una palabra del lenguaje
    (código)\includegraphics[width=2.522cm,height=0.7cm]{./ObjectReplacements/Object 273}será\includegraphics[width=4.551cm,height=1.162cm]{./ObjectReplacements/Object 274},
    y
    \includegraphics[width=0.868cm,height=0.506cm]{./ObjectReplacements/Object 275}es
    la longitud de la
    representación\includegraphics[width=0.392cm,height=0.467cm]{./ObjectReplacements/Object 276}.
    Este tipo de representación es la más usada en el ámbito de los
    números desde hace seiscientos o setecientos años más o menos en
    Occidente, sólo que para base 10 (diez dígitos distintos). Nosotros
    usaremos también bastante la base 2, código que llamamos
    habitualmente binario natural.
  \item
    Aquí es conveniente tener claro cómo pasamos de una base a otra. En
    general lo ha\-cemos por pasos:

    \begin{enumerate}
    \def\labelenumiii{\arabic{enumiii}.}
    \item
      Paso de una
      base\includegraphics[width=0.487cm,height=0.467cm]{./ObjectReplacements/Object 277}a
      base\includegraphics[width=0.628cm,height=0.467cm]{./ObjectReplacements/Object 278}.
      Esto lo hacemos por aplicación directa de la fórmula expresada
      anteriormente en 2.6.
    \item
      Paso de
      base\includegraphics[width=0.628cm,height=0.467cm]{./ObjectReplacements/Object 280}a
      una
      base\includegraphics[width=0.487cm,height=0.467cm]{./ObjectReplacements/Object 279}.
      Esto se hace por el proceso inverso al ante\-rior, dividimos
      sucesivamente por
      \includegraphics[width=0.487cm,height=0.467cm]{./ObjectReplacements/Object 311}los
      cocientes y nos vamos quedando con los restos, dónde los primeros
      son los dígitos de más bajo peso, y los últimos los de más alto.
      El proceso termina naturalmente cuando el cociente a dividir es
      más pequeño que la
      base\includegraphics[width=0.487cm,height=0.467cm]{./ObjectReplacements/Object 312},
      cociente que pasa a ser el dígito
      \includegraphics[width=0.926cm,height=0.467cm]{./ObjectReplacements/Object 313}de
      la conver\-sión.
    \item
      Paso de una
      base\includegraphics[width=0.614cm,height=0.531cm]{./ObjectReplacements/Object 281}a
      una base
      \includegraphics[width=0.616cm,height=0.531cm]{./ObjectReplacements/Object 282}:

      \begin{enumerate}
      \def\labelenumiv{\arabic{enumiv}.}
      \tightlist
      \item
        Caso
        que\includegraphics[width=1.96cm,height=0.591cm]{./ObjectReplacements/Object 283}.
        Cada dígito
        de\includegraphics[width=0.614cm,height=0.531cm]{./ObjectReplacements/Object 284}lo
        desarrollamos como su co\-rrespondiente
        en\includegraphics[width=0.466cm,height=0.467cm]{./ObjectReplacements/Object 285}dígitos
        de
        base\includegraphics[width=0.616cm,height=0.531cm]{./ObjectReplacements/Object 286}(sin
        ahorrar los 0 a la izquierda). Así hemos terminado.
      \item
        Caso
        que\includegraphics[width=1.939cm,height=0.591cm]{./ObjectReplacements/Object 287}.
        Comenzando por el
        \includegraphics[width=0.714cm,height=0.467cm]{./ObjectReplacements/Object 288}(el
        más a la derecha) agrupamos los dígitos
        de\includegraphics[width=1.155cm,height=0.467cm]{./ObjectReplacements/Object 290}dígitos
        de\includegraphics[width=0.614cm,height=0.531cm]{./ObjectReplacements/Object 289}.
        Cada grupo
        de\includegraphics[width=0.43cm,height=0.467cm]{./ObjectReplacements/Object 291}dí\-gitos
        de\includegraphics[width=0.614cm,height=0.531cm]{./ObjectReplacements/Object 292}es
        exactamente un dígito
        de\includegraphics[width=0.616cm,height=0.531cm]{./ObjectReplacements/Object 293},
        hacemos así la sustitución predicha y hemos terminado. El único
        posible problema es el grupo
        de\includegraphics[width=0.43cm,height=0.467cm]{./ObjectReplacements/Object 296}dígitos
        más alto, en el caso que tenga
        entre\includegraphics[width=0.411cm,height=0.467cm]{./ObjectReplacements/Object 294}y\includegraphics[width=0.993cm,height=0.467cm]{./ObjectReplacements/Object 295}dígitos.
        Ese grupo se rellena
        con\includegraphics[width=0.573cm,height=0.531cm]{./ObjectReplacements/Object 297}(el
        que representa
        a\includegraphics[width=0.422cm,height=0.467cm]{./ObjectReplacements/Object 298})
        por la izquierda hasta completar la
        longitud\includegraphics[width=0.43cm,height=0.467cm]{./ObjectReplacements/Object 299}.
      \item
        Caso
        que\includegraphics[width=4.688cm,height=0.591cm]{./ObjectReplacements/Object 300}.
        Lo dividimos en dos pasos: pasamos
        primero\includegraphics[width=0.614cm,height=0.531cm]{./ObjectReplacements/Object 301}a\includegraphics[width=0.624cm,height=0.591cm]{./ObjectReplacements/Object 302}y
        entonces
        de\includegraphics[width=1.842cm,height=0.591cm]{./ObjectReplacements/Object 303}a\includegraphics[width=0.616cm,height=0.531cm]{./ObjectReplacements/Object 304}.
      \item
        Caso en que no hay una relación
        entre\includegraphics[width=0.614cm,height=0.531cm]{./ObjectReplacements/Object 310}y\includegraphics[width=0.616cm,height=0.531cm]{./ObjectReplacements/Object 309}de
        las anteriores. Enton\-ces pasamos de
        base\includegraphics[width=0.614cm,height=0.531cm]{./ObjectReplacements/Object 308}a
        base\includegraphics[width=0.628cm,height=0.467cm]{./ObjectReplacements/Object 307},
        y entonces de
        base\includegraphics[width=0.628cm,height=0.467cm]{./ObjectReplacements/Object 306}a\includegraphics[width=0.616cm,height=0.531cm]{./ObjectReplacements/Object 305}.
        Este método es general, pero el anterior es más eficiente para
        los casos espe\-cificados.
      \end{enumerate}
    \end{enumerate}
  \item
    Para describir letras, esto es, caracteres alfabéticos de la lengua
    natural, utilizamos generalmente el código ASCII (de 7 bits el
    estándar original o de 8 dígitos binarios o bits, el extendido de
    Microsoft). Existen otros códigos, como el EBCDIC de IBM y las
    implementaciones UTF8, UTF16 y UTF32 de Unicode. Permiten
    sobradamente trasladar los lenguajes naturales escritos a un
    alfabeto binario.
  \item
    El conjunto de los naturales es insuficiente para muchas
    aplicaciones. La primera ampliación es el
    conjunto\includegraphics[width=0.559cm,height=0.385cm]{./ObjectReplacements/Object 314}de
    los enteros. Las formas de representar números enteros es variada
    pero la vamos a resumir en tres formas:

    \begin{enumerate}
    \def\labelenumiii{\arabic{enumiii}.}
    \item
      Magnitud y Signo
      (\includegraphics[width=1.291cm,height=0.467cm]{./ObjectReplacements/Object 629}).
      Representamos el valor absoluto del número como un número en
      binario natural y le añadimos
      un\includegraphics[width=0.422cm,height=0.467cm]{./ObjectReplacements/Object 318}o
      un\includegraphics[width=0.411cm,height=0.467cm]{./ObjectReplacements/Object 319}en
      el\includegraphics[width=0.926cm,height=0.467cm]{./ObjectReplacements/Object 317},
      dónde
      el\includegraphics[width=0.422cm,height=0.467cm]{./ObjectReplacements/Object 320}representa
      el signo
      ``\includegraphics[width=0.554cm,height=0.467cm]{./ObjectReplacements/Object 316}''
      y
      el\includegraphics[width=0.411cm,height=0.467cm]{./ObjectReplacements/Object 321}representa
      el signo
      ``\includegraphics[width=0.492cm,height=0.467cm]{./ObjectReplacements/Object 315}''.
      Este lenguaje sería en
      general\includegraphics[width=6.932cm,height=0.84cm]{./ObjectReplacements/Object 633},
      habida cuenta
      que\includegraphics[width=1.55cm,height=0.617cm]{./ObjectReplacements/Object 323}.
      Esta es una traducción directa de lo que hacemos cuando escribimos
      un número entero con signo. En general tenemos 2
      representa\-ciones para el número 0.
    \item
      En Complemento a la base
      \includegraphics[width=0.487cm,height=0.467cm]{./ObjectReplacements/Object 1314}(\includegraphics[width=0.963cm,height=0.467cm]{./ObjectReplacements/Object 630}).
      En este método de representación tene\-mos una forma de
      representar los números positivos y otra los negativos. Los
      nú\-meros positivos (incluido el
      \includegraphics[width=0.607cm,height=0.601cm]{./ObjectReplacements/Object 324}de
      valor\includegraphics[width=0.422cm,height=0.467cm]{./ObjectReplacements/Object 325})
      se representan de forma idéntica a como se hace
      en\includegraphics[width=1.291cm,height=0.467cm]{./ObjectReplacements/Object 204}.
      La novedad está en los números negativos, cuyo
      \includegraphics[width=0.926cm,height=0.467cm]{./ObjectReplacements/Object 326}es\includegraphics[width=0.411cm,height=0.467cm]{./ObjectReplacements/Object 327},
      seguido de un dígito cualquiera distinto del más alto
      \includegraphics[width=1.438cm,height=0.601cm]{./ObjectReplacements/Object 328}y
      seguido de una ristra finita cualquiera de dígitos.

      \begin{enumerate}
      \def\labelenumiv{\arabic{enumiv}.}
      \item
        Primero introduciremos una operación, la complementación de
        dígitos que designaremos como un operador
        prefijo\includegraphics[width=1.207cm,height=0.601cm]{./ObjectReplacements/Object 329},
        de forma que dada una ris\-tra cualquiera sobre un
        alfabeto\includegraphics[width=0.711cm,height=0.617cm]{./ObjectReplacements/Object 330},
        formalmente el
        lenguaje\includegraphics[width=1.549cm,height=0.838cm]{./ObjectReplacements/Object 331},
        que definimos como:

        \includegraphics[width=8.211cm,height=0.665cm]{./ObjectReplacements/Object 322}

        \includegraphics[width=8.243cm,height=0.72cm]{./ObjectReplacements/Object 631}

        \includegraphics[width=6.608cm,height=0.639cm]{./ObjectReplacements/Object 632}\includegraphics[width=12.95cm,height=0.723cm]{./ObjectReplacements/Object 634}
      \end{enumerate}

      Otra forma de llamar a esta operación es complemento a la base
      menos uno. Permite hacer operaciones sin ningún tipo de acarreo,
      dígito a dígito sin de\-pendencias de la posición. Algunas
      propiedades de esta operación son:

      \includegraphics[width=6.156cm,height=0.723cm]{./ObjectReplacements/Object 333}

      \includegraphics[width=10.312cm,height=0.734cm]{./ObjectReplacements/Object 334}

      Esta operación es interesante porque nos va a permitir definir la
      complementa\-ción a la base
      \includegraphics[width=0.496cm,height=0.321cm]{./ObjectReplacements/Object 335}
      de forma sencilla:

      \includegraphics[width=8.885cm,height=0.762cm]{./ObjectReplacements/Object 336}

      Las propiedades de esta complementación son:

      \includegraphics[width=4.731cm,height=0.617cm]{./ObjectReplacements/Object 338}

      \includegraphics[width=9.28cm,height=0.734cm]{./ObjectReplacements/Object 337}

      \begin{enumerate}
      \def\labelenumiv{\arabic{enumiv}.}
      \setcounter{enumiv}{1}
      \item
        Ahora ya podemos saber como interpretar los números negativos en
        represen\-tación
        \includegraphics[width=0.963cm,height=0.467cm]{./ObjectReplacements/Object 196}.
        El
        \includegraphics[width=0.926cm,height=0.467cm]{./ObjectReplacements/Object 339}será\includegraphics[width=0.411cm,height=0.467cm]{./ObjectReplacements/Object 340},
        a continuación no tendremos un
        dígito\includegraphics[width=1.064cm,height=0.601cm]{./ObjectReplacements/Object 341},
        sino cualquier otro, y una ristra finita cualquiera de dígitos
        sobre el
        alfabeto\includegraphics[width=0.711cm,height=0.617cm]{./ObjectReplacements/Object 342}.
        Para cualquier palabra sobre el alfabeto dicho definimos el
        valor en Complemento a la
        base\includegraphics[width=0.487cm,height=0.467cm]{./ObjectReplacements/Object 1315}con
        1 y con 0.

        \includegraphics[width=7.461cm,height=0.838cm]{./ObjectReplacements/Object 332}

        \includegraphics[width=6.137cm,height=1.309cm]{./ObjectReplacements/Object 635}

        \includegraphics[width=0.466cm,height=0.467cm]{./ObjectReplacements/Object 636}

        \includegraphics[width=5.046cm,height=1.309cm]{./ObjectReplacements/Object 637}
      \item
        Ahora nos disponemos a representar el lenguaje de las
        representaciones ente\-ras
        en\includegraphics[width=0.979cm,height=0.467cm]{./ObjectReplacements/Object 886},
        y podemos ver bien que valores tienen:
      \end{enumerate}

      \includegraphics[width=9.342cm,height=0.85cm]{./ObjectReplacements/Object 344}

      \includegraphics[width=8.821cm,height=0.85cm]{./ObjectReplacements/Object 345}

      \includegraphics[width=5.796cm,height=0.661cm]{./ObjectReplacements/Object 346}

      \includegraphics[width=10.051cm,height=1.566cm]{./ObjectReplacements/Object 347}

      Y así queda definido el lenguaje de las representaciones en
      complemento a una base y su semántica de valores
      en\includegraphics[width=0.557cm,height=0.385cm]{./ObjectReplacements/Object 348}.
    \item
      Exceso
      a\includegraphics[width=0.422cm,height=0.467cm]{./ObjectReplacements/Object 349}.
      Este tipo de representación no necesita signo siendo, el valor de
      la representación su valor en binario natural
      menos\includegraphics[width=0.422cm,height=0.467cm]{./ObjectReplacements/Object 350}.
      La representación
      \includegraphics[width=0.422cm,height=0.467cm]{./ObjectReplacements/Object 352}tiene
      como
      valor\includegraphics[width=0.774cm,height=0.467cm]{./ObjectReplacements/Object 353}.
      En general el valor de la
      representación\includegraphics[width=0.392cm,height=0.467cm]{./ObjectReplacements/Object 355}en
      Exceso a
      \includegraphics[width=0.422cm,height=0.467cm]{./ObjectReplacements/Object 354}será\includegraphics[width=1.549cm,height=0.506cm]{./ObjectReplacements/Object 356},
      siendo la
      función\includegraphics[width=1.706cm,height=0.467cm]{./ObjectReplacements/Object 887}la
      correspondiente al bina\-rio natural.
    \item
      La siguiente ampliación
      es\includegraphics[width=0.52cm,height=0.385cm]{./ObjectReplacements/Object 357}.
      Existen de entrada dos formas comunes de repre\-sentación: en
      punto fijo y en punto flotante.
    \item
      En punto fijo tendremos siempre que saber dónde se encuentra el
      punto o coma deci\-mal, conociendo cual es la longitud de la parte
      fraccionaria
      (\includegraphics[width=1.219cm,height=0.697cm]{./ObjectReplacements/Object 358})
      y cual la de la parte entera
      (\includegraphics[width=1.325cm,height=0.61cm]{./ObjectReplacements/Object 359}).
      Hasta el momento solo hemos utilizado
      \includegraphics[width=3.256cm,height=0.829cm]{./ObjectReplacements/Object 360}dónde\includegraphics[width=1.563cm,height=0.697cm]{./ObjectReplacements/Object 361}.
    \item
      Representaciones de binario natural en punto fijo. Vamos a ver, en
      un primer mo\-mento, estas representaciones solo si son positivas,
      esto es, representaciones de
      \includegraphics[width=0.741cm,height=0.568cm]{./ObjectReplacements/Object 362}en
      binario natural, que será la base para el resto de
      representaciones más com\-plejas. En principio la parte entera
      vendrá dada por una expresión del
      tipo\includegraphics[width=3.604cm,height=0.76cm]{./ObjectReplacements/Object 363},\includegraphics[width=3.313cm,height=0.61cm]{./ObjectReplacements/Object 365},\includegraphics[width=3.244cm,height=0.887cm]{./ObjectReplacements/Object 364},\includegraphics[width=0.824cm,height=0.467cm]{./ObjectReplacements/Object 366},\includegraphics[width=7.832cm,height=0.887cm]{./ObjectReplacements/Object 368}y\includegraphics[width=3.256cm,height=0.829cm]{./ObjectReplacements/Object 367}.
      Para evaluar el valor de la representación la fórmula es la
      clásica\includegraphics[width=3.874cm,height=1.307cm]{./ObjectReplacements/Object 369}.
      En general aunque ponemos explícita\-mente el punto, no es
      necesario una vez se conocen las longitudes de las partes entera y
      fraccionaria.
    \item
      El siguiente paso es ¿cómo pasamos un número en punto fijo de una
      base a otra?. Para esto lo más sencillo es separar el número en
      punto fijo en dos partes, la entera y la fraccionaria. La parte
      entera, un número natural, seguirá siendo entera en cualquier
      base, por lo que aplicamos las reglas de conversión que ya
      conocemos para n-ario na\-tural, ya vistas.
    \end{enumerate}

    ¿Y la parte fraccionaria?. También sigue siendo un número
    \includegraphics[width=1.499cm,height=0.538cm]{./ObjectReplacements/Object 105}para
    cualquier base, y en cualquier base sigue representándose como una
    cadena de\-trás del punto fijo. A partir de ahora la
    representación\includegraphics[width=0.392cm,height=0.467cm]{./ObjectReplacements/Object 443},
    representará siempre la parte fraccionaria, esto
    es\includegraphics[width=2.351cm,height=0.506cm]{./ObjectReplacements/Object 444},
    y su representación será siempre tal que confundiremos
    \includegraphics[width=0.392cm,height=0.467cm]{./ObjectReplacements/Object 445}con
    \includegraphics[width=0.803cm,height=0.467cm]{./ObjectReplacements/Object 446}y\includegraphics[width=3.748cm,height=0.663cm]{./ObjectReplacements/Object 447}.
    Veremos los distintos ca\-sos:

    \begin{enumerate}
    \def\labelenumiii{\arabic{enumiii}.}
    \setcounter{enumiii}{7}
    \item
      Casos en que existe una relación
      entre\includegraphics[width=0.714cm,height=0.635cm]{./ObjectReplacements/Object 448}y\includegraphics[width=0.716cm,height=0.635cm]{./ObjectReplacements/Object 450}tal
      que\includegraphics[width=7.322cm,height=0.568cm]{./ObjectReplacements/Object 449}.
      Los cambios son idénticos a los realizados para la parte entera
      excepto que los grupos de dígitos se cogen desde el punto decimal
      hacia la derecha, esto es, en sentido inverso al que tomábamos
      para los naturales.
    \item
      Caso general. Hemos de utilizar una base intermedia, la habitual
      base
      10\includegraphics[width=0.905cm,height=0.635cm]{./ObjectReplacements/Object 451}.
      El método es el mismo explicado en su momento para números
      naturales en
      \includegraphics[width=2.798cm,height=0.467cm]{./ObjectReplacements/Object 343}.
      Solo que da pues ver como pasamos de
      base\includegraphics[width=0.905cm,height=0.635cm]{./ObjectReplacements/Object 452}a\includegraphics[width=0.714cm,height=0.635cm]{./ObjectReplacements/Object 453}y
      viceversa.

      \begin{enumerate}
      \def\labelenumiv{\arabic{enumiv}.}
      \item
        El caso inmediato
        es\includegraphics[width=1.794cm,height=0.635cm]{./ObjectReplacements/Object 454}.
        Utilizamos la siguiente
        fórmula:\includegraphics[width=4.621cm,height=1.319cm]{./ObjectReplacements/Object 457}.
      \item
        El caso inmediato
        es\includegraphics[width=1.794cm,height=0.635cm]{./ObjectReplacements/Object 455}.
        Se trata del proceso inverso al anterior, esto es, como las
        potencias son negativas realizamos, en el caso anterior
        divi\-siones sucesivas
        por\includegraphics[width=0.453cm,height=0.492cm]{./ObjectReplacements/Object 456},
        el cardinal de la base. Pro lo tanto el proceso que nos atañe
        será el de multiplicar sucesivamente por la
        base\includegraphics[width=0.453cm,height=0.492cm]{./ObjectReplacements/Object 458},
        el número
        fraccionario\includegraphics[width=1.011cm,height=0.492cm]{./ObjectReplacements/Object 459}.
        En cada multiplicación por la base la parte entera será un
        dígito de la nueva base, y los obtenemos desde el primer lugar a
        la derecha del punto decimal hacia la derecha. En cada
        multiplicación nos quedamos con el resto, esto es, la parte
        fraccionaria del resultado (la parte entera ya ha sido
        incorporada al resultado final).
      \item
        Podemos observar fácilmente que en los dos casos anteriores el
        procedimien\-to de colocar decimales en la nueva base no tiene
        por qué ser un proceso que termine, esto es, finito. Es
        claro\includegraphics[width=0.811cm,height=0.467cm]{./ObjectReplacements/Object 460}en
        base
        \includegraphics[width=0.686cm,height=0.617cm]{./ObjectReplacements/Object 461}tiene
        una representa\-ción
        finita\includegraphics[width=2.193cm,height=0.628cm]{./ObjectReplacements/Object 462},
        pero en
        base\includegraphics[width=0.684cm,height=0.617cm]{./ObjectReplacements/Object 463},
        en binario, y en base
        \includegraphics[width=0.707cm,height=0.492cm]{./ObjectReplacements/Object 464}y
        en otras bases tenemos
        que\includegraphics[width=1.131cm,height=0.506cm]{./ObjectReplacements/Object 466}\includegraphics[width=0.554cm,height=0.467cm]{./ObjectReplacements/Object 467}\includegraphics[width=1.162cm,height=0.667cm]{./ObjectReplacements/Object 468}\includegraphics[width=0.554cm,height=0.467cm]{./ObjectReplacements/Object 469}\includegraphics[width=1.182cm,height=0.669cm]{./ObjectReplacements/Object 470}\includegraphics[width=0.554cm,height=0.467cm]{./ObjectReplacements/Object 471}\includegraphics[width=0.923cm,height=0.619cm]{./ObjectReplacements/Object 472}\includegraphics[width=0.554cm,height=0.467cm]{./ObjectReplacements/Object 473}\includegraphics[width=0.967cm,height=0.667cm]{./ObjectReplacements/Object 474}\includegraphics[width=0.554cm,height=0.467cm]{./ObjectReplacements/Object 475}\includegraphics[width=1.182cm,height=0.667cm]{./ObjectReplacements/Object 476}\includegraphics[width=0.554cm,height=0.467cm]{./ObjectReplacements/Object 477}\includegraphics[width=0.967cm,height=0.619cm]{./ObjectReplacements/Object 478}\includegraphics[width=0.554cm,height=0.467cm]{./ObjectReplacements/Object 479}\includegraphics[width=0.967cm,height=0.667cm]{./ObjectReplacements/Object 480}\includegraphics[width=0.554cm,height=0.467cm]{./ObjectReplacements/Object 481}\includegraphics[width=1.184cm,height=0.667cm]{./ObjectReplacements/Object 482}\includegraphics[width=0.554cm,height=0.467cm]{./ObjectReplacements/Object 483}\includegraphics[width=0.972cm,height=0.619cm]{./ObjectReplacements/Object 484}\includegraphics[width=1.362cm,height=0.467cm]{./ObjectReplacements/Object 465}\includegraphics[width=1.184cm,height=0.619cm]{./ObjectReplacements/Object 607}\includegraphics[width=1.362cm,height=0.467cm]{./ObjectReplacements/Object 608}\includegraphics[width=1.161cm,height=0.619cm]{./ObjectReplacements/Object 609}\includegraphics[width=1.362cm,height=0.467cm]{./ObjectReplacements/Object 610}\includegraphics[width=1.161cm,height=0.663cm]{./ObjectReplacements/Object 602}\includegraphics[width=1.362cm,height=0.467cm]{./ObjectReplacements/Object 603}\includegraphics[width=1.168cm,height=0.619cm]{./ObjectReplacements/Object 600}\includegraphics[width=1.009cm,height=0.467cm]{./ObjectReplacements/Object 601}.

        Una vez hemos utilizado la representación en binario, las de
        base 4 y la octal (base 8) son inmediatas por reagrupamientos.
        Igualmente la representación en base 3 es inmediata ya
        que\includegraphics[width=2.369cm,height=0.631cm]{./ObjectReplacements/Object 485},
        que corresponde
        a\includegraphics[width=5.482cm,height=0.704cm]{./ObjectReplacements/Object 486},
        y confundiremos
        habitualmente\includegraphics[width=0.644cm,height=0.635cm]{./ObjectReplacements/Object 488}con
        \includegraphics[width=0.453cm,height=0.492cm]{./ObjectReplacements/Object 487}y
        queda claro que para cualquier otra posición tenemos
        que\includegraphics[width=5.733cm,height=0.683cm]{./ObjectReplacements/Object 489}.
        Igualmente tenemos por reagrupamien\-to la representación para
        base 9.
      \item
        Es importante darnos cuenta que en el caso de pasos entre
        \includegraphics[width=0.714cm,height=0.635cm]{./ObjectReplacements/Object 496}y\includegraphics[width=0.716cm,height=0.635cm]{./ObjectReplacements/Object 497}tal
        que
        \includegraphics[width=7.322cm,height=0.568cm]{./ObjectReplacements/Object 498},
        si la representación fuente es finita la destino también los
        será, y si la fuente es infinita con un periodo, el destino
        también lo será. En general si pasamos un número de
        base\includegraphics[width=0.714cm,height=0.635cm]{./ObjectReplacements/Object 604}a
        base\includegraphics[width=1.007cm,height=0.635cm]{./ObjectReplacements/Object 605}la
        representación seguirá guardando su carácter finito no
        periódi\-co si así ocurre en
        \includegraphics[width=0.714cm,height=0.635cm]{./ObjectReplacements/Object 606}(como
        vemos al pasar de base 3 a base 6).
      \item
        Habitualmente utilizaremos el guarismo confundido con el
        dígito\includegraphics[width=1.046cm,height=0.496cm]{./ObjectReplacements/Object 490}\includegraphics[width=1.487cm,height=0.467cm]{./ObjectReplacements/Object 491}\includegraphics[width=1.55cm,height=0.492cm]{./ObjectReplacements/Object 492}.
        Para dígitos superiores, si los hubiera, hay varias
        es\-trategias. En Electrónica Digital se suele utilizar el
        alfabeto latino, sin dife\-renciar mayúsculas de minúsculas, lo
        más habitual para trabajar en base 16
        (hexadecimal),\includegraphics[width=9.68cm,height=0.7cm]{./ObjectReplacements/Object 493}y
        en general podemos utilizar más letras, y más letras aún de
        otros abecedarios (y aún distinguir entre mayúsculas y
        minúsculas, acentuarlas de diferentes formas e inventarnos
        nuevas letras-gráficos, pero en general no es solución general
        del problema). Una solución general es utilizar cadenas que
        conten\-gan el índice del dígito correspondiente en base 10
        (comprensible para todos) en una cadena que exprese la base en
        la que nos encontramos
        con\includegraphics[width=6.535cm,height=0.639cm]{./ObjectReplacements/Object 494},
        que es la solución que hemos adop\-tado en el
        programa\includegraphics[width=1.87cm,height=0.467cm]{./ObjectReplacements/Object 495}que
        podéis disponer para hacer conversio\-nes etc en cualquier base.
        En este programa (en realidad un conjunto de clases del lenguaje
        de programación C++), además de representar y operar con
        dígi\-tos de cualquier base, podemos representar y operar con
        números naturales en cualquier base
        (representación\includegraphics[width=2.798cm,height=0.467cm]{./ObjectReplacements/Object 599})
        y con números enteros de cualquier base (representación interna
        en\includegraphics[width=0.963cm,height=0.467cm]{./ObjectReplacements/Object 597}).
        Los naturales los repre\-sentamos
        como\includegraphics[width=10.236cm,height=0.67cm]{./ObjectReplacements/Object 595}y
        los enteros por defecto entran y salen como
        en\includegraphics[width=0.963cm,height=0.467cm]{./ObjectReplacements/Object 596}.
        Hay algunas facili\-dades para sacar los números enteros en
        pantalla
        en\includegraphics[width=1.291cm,height=0.467cm]{./ObjectReplacements/Object 598}.
      \item
        Por lo demás es fácil conseguir la representación de la parte
        fraccionaria
        en\includegraphics[width=0.963cm,height=0.467cm]{./ObjectReplacements/Object 611}si
        ésta es negativa. Se realizan las operaciones de la misma forma
        que la hacíamos con la parte entera, solo que allí teníamos en
        cuenta la longi\-tud de la representación, pero aquí, para la
        parte fraccionaria el complemento lo conseguimos básicamente con
        el complemento de la parte fraccionaria a la
        unidad\includegraphics[width=2.185cm,height=0.492cm]{./ObjectReplacements/Object 612}.
      \end{enumerate}
    \end{enumerate}
  \item
    El último tipo de representación que vamos a ver es la
    representación en punto flo\-tante
    (\includegraphics[width=2.524cm,height=0.467cm]{./ObjectReplacements/Object 888},
    de forma que en C, C++ al tipo de números con decimales que
    habitualmente llamamos reales, toman el nombre de su forma de
    representación :
    \includegraphics[width=1.021cm,height=0.467cm]{./ObjectReplacements/Object 889}).
    Este tipo está estandarizado y tenemos un documento que lo detalla
    bastante bien. Este tipo tiene ventajas (relativas) con respecto al
    punto fijo porque amplía bastante el rango de valores a representar
    con el mismo número de dígitos, y porque se controla el error de
    representación en forma proporcional al valor absoluto del número,
    esto es, trabajamos con errores relativos, mientras en punto fijo se
    trabaja con errores absolutos que son fijos en todo el rango de
    representaciones. Como parte negativa hay que decir que los
    circuitos para trabajar con punto flotante son más complejos y
    voluminosos que los que teníamos para punto fijo, que eran los
    mismos circuitos que teníamos para trabajar con enteros.
  \item
    Siguiendo con los códigos que utilizamos en Electrónica Digital, los
    más comunes son los códigos o lenguajes que tienen una longitud
    fija. Por ver un poco la noción de longitud, ésta determina
    básicamente la cantidad de valores de diferentes, de palabras
    diferentes que puede tener el lenguaje. Si trabajamos con el
    lenguaje\includegraphics[width=0.457cm,height=0.492cm]{./ObjectReplacements/Object 617}sobre
    el
    alfabeto\includegraphics[width=0.482cm,height=0.492cm]{./ObjectReplacements/Object 616},
    definiremos:

    \includegraphics[width=3.607cm,height=0.603cm]{./ObjectReplacements/Object 615}

    \includegraphics[width=4.263cm,height=0.603cm]{./ObjectReplacements/Object 618}

    \includegraphics[width=7.336cm,height=0.61cm]{./ObjectReplacements/Object 638}

    \includegraphics[width=7.176cm,height=0.61cm]{./ObjectReplacements/Object 614}

    \includegraphics[width=9.938cm,height=0.64cm]{./ObjectReplacements/Object 620}

    A este último lo suelo llamar código saturado de
    longitud\includegraphics[width=0.453cm,height=0.492cm]{./ObjectReplacements/Object 621}sobre
    el
    alfabeto\includegraphics[width=0.482cm,height=0.492cm]{./ObjectReplacements/Object 622},
    y será la referencia para los distintos códigos de longitud fija.
    Para un código de este tipo el número de palabras máximo permitido
    será\includegraphics[width=1.235cm,height=0.633cm]{./ObjectReplacements/Object 623}.
    Si trabajamos en binario natural de longitud
    fija\includegraphics[width=0.453cm,height=0.492cm]{./ObjectReplacements/Object 624},
    representaremos desde
    el\includegraphics[width=0.453cm,height=0.492cm]{./ObjectReplacements/Object 626}al\includegraphics[width=1.251cm,height=0.594cm]{./ObjectReplacements/Object 625}.
  \item
    Por ejemplo, el código
    \includegraphics[width=3.092cm,height=0.492cm]{./ObjectReplacements/Object 890}está
    dentro
    de\includegraphics[width=1.505cm,height=0.774cm]{./ObjectReplacements/Object 891},
    pero no son idénticos. Enumeraré los valores en una tabla de dos
    columnas:
  \end{enumerate}
\end{enumerate}

{\def\LTcaptype{none} % do not increment counter
\begin{longtable}[]{@{}ll@{}}
\toprule\noalign{}
\includegraphics[width=3.092cm,height=0.492cm]{./ObjectReplacements/Object 892}
&
\includegraphics[width=1.505cm,height=0.774cm]{./ObjectReplacements/Object 893} \\
\midrule\noalign{}
\endhead
\bottomrule\noalign{}
\endlastfoot
0000 & 0000 \\
0001 & 0001 \\
0010 & 0010 \\
0011 & 0011 \\
0100 & 0100 \\
0101 & 0101 \\
0110 & 0110 \\
0111 & 0111 \\
1000 & 1000 \\
1001 & 1001 \\
-\/-\/-\/- & 1010 \\
-\/-\/-\/- & 1011 \\
-\/-\/-\/- & 1100 \\
-\/-\/-\/- & 1101 \\
-\/-\/-\/- & 1110 \\
-\/-\/-\/- & 1111 \\
\end{longtable}
}

\begin{enumerate}
\def\labelenumi{\arabic{enumi}.}
\setcounter{enumi}{42}
\item
  \begin{enumerate}
  \def\labelenumii{\arabic{enumii}.}
  \item
    Los
    códigos\includegraphics[width=0.961cm,height=0.492cm]{./ObjectReplacements/Object 894}son
    códigos de longitud fija, por lo general 4, pero lo fundamen\-tal es
    que remedan en binario el
    alfabeto\includegraphics[width=2.503cm,height=0.635cm]{./ObjectReplacements/Object 895}.
    Dependiendo de la finali\-dad hay varios: el más sencillo el que
    acabamos de dar.
    El\includegraphics[width=3.739cm,height=0.492cm]{./ObjectReplacements/Object 896},
    es de longitud 4, pero con el 0 en 0011 y el 9 en 1100. Comprobaréis
    fácilmente que es un código dónde el
    \includegraphics[width=1.946cm,height=0.676cm]{./ObjectReplacements/Object 897}coincide
    con la negación lógica bit a bit. Esto fa\-cilita el hacer
    operaciones directamente
    en\includegraphics[width=0.905cm,height=0.635cm]{./ObjectReplacements/Object 898}.
    El
    código\includegraphics[width=2.85cm,height=0.492cm]{./ObjectReplacements/Object 899},
    es un código auto-complementario como el anterior pero además
    mantiene un sistema de pesos 2-4-2-1, en orden con su valores
    decimales sería:

    \includegraphics[width=11.098cm,height=1.048cm]{./ObjectReplacements/Object 900}
  \item
    Otro tipo de códigos son los códigos continuos y los cíclicos.
    Decimos que un código es continuo si entre una palabra y la
    siguiente en el orden propio (valor semántico o significado) solo
    varía un dígito. Además es cíclico si entre la primera palabra y la
    última varía a su vez un solo bit. Para casos con más de dos valores
    en el alfabeto ha\-bría que afinar la definición, pero para el caso
    que nos ocupa, que es el código Gray de
    \includegraphics[width=1.397cm,height=0.467cm]{./ObjectReplacements/Object 901}basta
    con lo dicho y los ejemplos que a continuación se van a dar. Estos
    códigos son además códigos reflejados, esto es, se obtienen por un
    sistema especular. Para un solo bit sería:

    0

    1

    Para dos bits colocamos el anterior tal cual y lo copiamos en forma
    especular hacia abajo:

    0

    1

    1

    0

    Ahora añadimos 0 en el bit izquierdo a los primeros y 1 a los
    copiados:

    00

    01

    11

    10

    Ya hemos obtenido el código Gray de 2 bits. Pongo el proceso en una
    tabla para Gray de 3 bits. Comenzamos por el Gray de 2 bits en la
    primera columna, lo reflejamos en la segunda, y añadimos lo 0s y 1s
    para distinguir los bits más altos de los más bajos, quedando el
    Gray de 3 bits en la última columna:

    \includegraphics[width=2.469cm,height=3.995cm]{./ObjectReplacements/Object 902}

    Para 4 bits sería:

    \includegraphics[width=3.101cm,height=8.027cm]{./ObjectReplacements/Object 903}

    Y así sucesivamente.
  \item
    Los códigos biquinarios (el 2 entre 5 en concreto, entre los muchos
    biquinarios que de hecho se han utilizado), es un código que
    mantiene el número de 1s en cada pala\-bra del código que tiene
    longitud constante 5. Además proviene de un código ponde\-rado, pero
    tal como lo ponemos aquí ya no lo es:

    \includegraphics[width=10.142cm,height=0.506cm]{./ObjectReplacements/Object 904}

    \includegraphics[width=10.098cm,height=0.506cm]{./ObjectReplacements/Object 905}

    Proviene del siguiente, que es ponderado, de 7 bits de longitud, y
    ponderación
    \includegraphics[width=3.854cm,height=0.467cm]{./ObjectReplacements/Object 908}:

    \includegraphics[width=12.137cm,height=0.506cm]{./ObjectReplacements/Object 907}

    \includegraphics[width=12.365cm,height=0.506cm]{./ObjectReplacements/Object 906}
  \item
    El siguiente y último código es el Johnson (Johnson-Möbius) de 5
    bits de longitud, que es un código continuo y progresivo, pero que
    puede ser de longitud fija
    de\includegraphics[width=0.453cm,height=0.492cm]{./ObjectReplacements/Object 913},
    con una capacidad de
    \includegraphics[width=0.843cm,height=0.492cm]{./ObjectReplacements/Object 914}valores
    distintos, así el de 5 bits es apropiado para re\-presentar
    un\includegraphics[width=2.738cm,height=0.492cm]{./ObjectReplacements/Object 915},
    y es apropiado para tratamiento muy rápido de la in\-formación
    mediante registros de desplazamientos y otros dispositivos:

    \includegraphics[width=10.072cm,height=0.506cm]{./ObjectReplacements/Object 909}

    \includegraphics[width=9.97cm,height=0.506cm]{./ObjectReplacements/Object 911}
  \end{enumerate}
\item
  Tratamiento del error en códigos de longitud fija.

  \begin{enumerate}
  \def\labelenumii{\arabic{enumii}.}
  \item
    Utilizaremos en principio varios métodos para esta finalidad:
    códigos Reed-Solomon, códigos lineales de grupo (dentro de estos se
    encuentran los códigos Hamming), códi\-gos Golay, distintos bits de
    paridad, códigos de
    repetición\includegraphics[width=1.48cm,height=0.519cm]{./ObjectReplacements/Object 1001}u
    otros, códigos CRC, y de suma igual. En general: en unos detectamos
    un error y devolvemos una peti\-ción de reenvío o similar, en otros
    cubrimos los casos más importantes y probables y la detección y
    corrección ha de ser autónoma por el receptor. Es importante saber
    que es lo que suele pasar cuando hay una comunicación de tramas de
    bits:

    \begin{enumerate}
    \def\labelenumiii{\arabic{enumiii}.}
    \item
      Tenemos un Emisor y un Receptor de mensaje, que comparten un
      protocolo común de actuación además del alfabeto básico y el
      código a utilizar.
    \item
      Entre el Emisor y el Receptor tenemos un Medio: una Línea de
      transmisión, sea esta cableada o no.
    \item
      El Medio/Línea es ruidoso: existe la probabilidad que se cambien
      uno símbolos de alfabeto por otros, obteniendo así una palabra del
      código válida o no. Si la palabra re\-cibida en el Emisor es
      recibida con cambios que no la convierten en una palabra prohibida
      (no del código común), el error pasará desapercibido. Por otra
      parte si per\-cibimos como Receptor un error podremos bien
      corregirlo, bien no.
    \item
      Suposiciones varias que hacen el tratamiento de errores manejable
      (suposiciones ra\-zonables).

      \begin{enumerate}
      \def\labelenumiv{\arabic{enumiv}.}
      \tightlist
      \item
        La longitud de la palabra enviada permanece inalterable. No se
        pierden bits por el trayecto: se transmutarán en otros pero no
        seguirá siendo
        un\includegraphics[width=1.939cm,height=0.467cm]{./ObjectReplacements/Object 916}.
      \item
        La probabilidad de error será pequeña, esto
        es\includegraphics[width=3.918cm,height=0.54cm]{./ObjectReplacements/Object 917},
        aunque de hecho no importaría esta otra
        situación\includegraphics[width=3.907cm,height=0.54cm]{./ObjectReplacements/Object 918},
        caso en que cambiamos 1s por 0s y viceversa, y cambiamos la
        probabilidad
        por\includegraphics[width=4.48cm,height=0.54cm]{./ObjectReplacements/Object 919}.
        El problema grave aparece
        cuando\includegraphics[width=5.36cm,height=0.556cm]{./ObjectReplacements/Object 920}con
        \includegraphics[width=2.17cm,height=0.482cm]{./ObjectReplacements/Object 921}(por
        poner un límite real\-mente permisivo. Así permitimos
        que\includegraphics[width=5.639cm,height=0.54cm]{./ObjectReplacements/Object 922},
        y por la reducción vista antes tenemos
        que\includegraphics[width=3.859cm,height=0.54cm]{./ObjectReplacements/Object 923}.
        Con un error de
        \includegraphics[width=0.951cm,height=0.467cm]{./ObjectReplacements/Object 924}el
        sistema es plenamente aleatorio. No hay absolutamente nada que
        hacer.
      \item
        El error en un bit ha de ser independiente de los que hay
        alrededor. Esto es a ve\-ces claramente no realista. La idea es
        que en una palabra
        recibida,\includegraphics[width=2.76cm,height=0.531cm]{./ObjectReplacements/Object 925}dados
        \includegraphics[width=8.594cm,height=0.54cm]{./ObjectReplacements/Object 926}.
      \item
        Por la suposición 2 y 3, dada la palabra que se
        envía\includegraphics[width=2.76cm,height=0.531cm]{./ObjectReplacements/Object 927},
        tene\-mos entonces
        que\includegraphics[width=2.574cm,height=0.482cm]{./ObjectReplacements/Object 928},
        tal que
        para\includegraphics[width=3.298cm,height=0.467cm]{./ObjectReplacements/Object 929}entonces\includegraphics[width=2.709cm,height=0.54cm]{./ObjectReplacements/Object 930},
        y
        así\includegraphics[width=6.115cm,height=0.591cm]{./ObjectReplacements/Object 931}.
        En
        general\includegraphics[width=4.637cm,height=0.591cm]{./ObjectReplacements/Object 932}.
      \item
        La probabilidad de error sobre un valor 0 o un valor 1 ha de ser
        muy parecida, de forma que la asimetría sea inapreciable.
      \end{enumerate}
    \item
      Hay varios conceptos importantes en cuanto a los errores: el de
      distancia Hamming y el de paridad.
    \item
      En un código de longitud fija, digamos
      un\includegraphics[width=1.972cm,height=0.492cm]{./ObjectReplacements/Object 933},
      la distancia Hamming en\-tre dos palabras de
      longitud\includegraphics[width=0.453cm,height=0.492cm]{./ObjectReplacements/Object 934}sobre
      un
      alfabeto\includegraphics[width=0.482cm,height=0.492cm]{./ObjectReplacements/Object 935},
      digamos
      \includegraphics[width=2.589cm,height=0.591cm]{./ObjectReplacements/Object 936}y
      \includegraphics[width=2.591cm,height=0.591cm]{./ObjectReplacements/Object 937},
      definimos:

      \includegraphics[width=6.473cm,height=0.826cm]{./ObjectReplacements/Object 938}

      En palabras es el número de bits que son diferentes entre amabas
      palabras (posición por posición).
    \item
      La función definida anteriormente, formalmente,
      \[d_{\mathtt{\mathrm{H}}}^{\mathtt{\mathrm{n}}}:L^{\mathtt{\mathrm{n}}}{\left( \mathtt{\mathrm{}_{2}} \right) \times L^{\mathtt{\mathrm{n}}}}{\left( \mathtt{\mathrm{}_{2}} \right)\rightarrow{\lbrack{0,n}\rbrack} \subset {\mathbb{N}} \subset {\mathbb{R}}}::\left( \mathtt{\mathrm{a}\mathrm{,}\mathrm{b}} \right)@\mathit{card}\left\{ {{\iota \in {\lbrack{0,{n - 1}}\rbrack}} \mid {l_{\iota}^{a} \neq l_{\iota}^{b}}} \right\}\]es
      una distancia bien definida matemáticamente, la distancia entre
      dos palabras iguales es siempre 0, y si son distintas es
      necesariamente distinta de 0. Siempre es un número mayor o igual
      que 0 como corresponde a un cardinal de un conjunto. Además es
      si\-métrica, esto es, la distancia entre dos palabras
      \includegraphics[width=1.478cm,height=0.492cm]{./ObjectReplacements/Object 940}
      es idéntica a la distancia
      \includegraphics[width=1.478cm,height=0.492cm]{./ObjectReplacements/Object 941}.
      Por último se cumple la desigualdad triangular, para tres palabras
      cuales\-quiera\includegraphics[width=3.179cm,height=0.693cm]{./ObjectReplacements/Object 942}la
      distancia Hamming
      cumple\includegraphics[width=5.442cm,height=0.616cm]{./ObjectReplacements/Object 943}.
      Estas tres propiedades pueden verse que normales entre las
      distancias euclidianas normales. De hecho esta distancia nos
      habilita para ver una geometría
      en\includegraphics[width=2.198cm,height=0.693cm]{./ObjectReplacements/Object 944},
      dónde podemos poner como puntos
      los\includegraphics[width=2.441cm,height=0.467cm]{./ObjectReplacements/Object 945}o
      bien
      los\includegraphics[width=2.512cm,height=0.467cm]{./ObjectReplacements/Object 946},
      en general, a estos, vistos desde el pris\-ma geométrico, se les
      llama\includegraphics[width=1.792cm,height=0.467cm]{./ObjectReplacements/Object 947}.
    \item
      Para un código cualquiera de longitud fija definimos ahora el
      concepto
      de\includegraphics[width=6.096cm,height=0.492cm]{./ObjectReplacements/Object 948},
      que es,
      dado\includegraphics[width=1.473cm,height=0.61cm]{./ObjectReplacements/Object 949},\includegraphics[width=6.137cm,height=1.147cm]{./ObjectReplacements/Object 952}.
      Para el lenguaje
      completo\includegraphics[width=2.027cm,height=0.676cm]{./ObjectReplacements/Object 950},
      esta distancia mínima es 1. La idea que sigue es muy intuitiva: si
      la distancia entre dos palabras es 0, entonces las dos palabras
      son en realidad la misma, son el mismo pun\-to-palabra del
      espacio-código de longitud fija. Si las palabras son de longitud
      24, 24 es la distancia más alejada entre dos palabras y así.

      Para determinar la distancia mínima de un
      \includegraphics[width=1.972cm,height=0.492cm]{./ObjectReplacements/Object 951}hay
      algunos trucos que nos ayudarán:

      \begin{enumerate}
      \def\labelenumiv{\arabic{enumiv}.}
      \item
        Con la primera distancia 1 que encontremos podemos dejar de
        comprobar. La distancia mínima de ese código es 1.
      \item
        \includegraphics[width=2.499cm,height=0.676cm]{./ObjectReplacements/Object 953}.
        Si\includegraphics[width=1.665cm,height=0.527cm]{./ObjectReplacements/Object 954},
        entonces
        \includegraphics[width=2.385cm,height=0.619cm]{./ObjectReplacements/Object 955}.
      \item
        Si\includegraphics[width=1.944cm,height=0.527cm]{./ObjectReplacements/Object 956},
        entonces\includegraphics[width=2.385cm,height=0.619cm]{./ObjectReplacements/Object 957}.
        Este resultado, tipo cota de la dis\-tancia mínima en función
        del cardinal del código, se puede extender para distancia mínima
        2, ... Por
        ejemplo\includegraphics[width=2.395cm,height=0.619cm]{./ObjectReplacements/Object 958}entonces\includegraphics[width=1.522cm,height=0.527cm]{./ObjectReplacements/Object 959}.
      \item
        Una forma aumentar la distancia mínima en uno cuando la
        distancia mínima es 1, y a veces en otras ocasiones, es añadir
        un bit de paridad. Un bit de paridad es un bit que nos dice si
        el número de 1s que hay en una palabra es par o impar. Un bit de
        paridad par, es un bit que vale 1 si y solo si el número de 1s
        que hay en el resto de la palabra es impar (con él incluido el
        número de 1s de la palabra es par). Un bit de paridad impar es
        un bit que vale 1 si y solo si el número de 1s que hay en el
        resto de la palabra es par (con él incluido el número de 1s de
        la palabra es impar). Si
        \includegraphics[width=3.427cm,height=0.676cm]{./ObjectReplacements/Object 960}y
        el bit de paridad lo ponemos (añadimos) en la posición
        \includegraphics[width=0.453cm,height=0.492cm]{./ObjectReplacements/Object 961},
        la fórmula
        \includegraphics[width=4.819cm,height=0.944cm]{./ObjectReplacements/Object 962}
        nos da el valor del bit de paridad par. El de paridad impar es
        exactamente el inverso del formulado. La ra\-zón por la que la
        fórmula funciona es sencilla, si recordamos
        que\includegraphics[width=1.628cm,height=0.467cm]{./ObjectReplacements/Object 963}y
        que\includegraphics[width=1.628cm,height=0.467cm]{./ObjectReplacements/Object 964}.
        Así, para un código de un solo bit, el de paridad sería
        simple\-mente la repetición del bit (el bit de paridad par sería
        simplemente del mismo va\-lor que el ya existente). De esta
        forma los dos bits son iguales y el número de 1s será o 0 o 2,
        esto es, siempre par. Si fuese el código original de 2 bits de
        longitud, el bit de paridad debería valer 1 si solo uno de los
        dos (no los dos) valiese 1. Esto vuelve a ser la suma exclusiva.
        Y así sucesivamente.
      \item
        Podemos realizar un bit de paridad par sobre los 0s de una
        palabra. Esto tiene mayor interés cuando la longitud de la
        palabra es impar. Solo hay que negar cada bit y utilizar la
        misma fórmula
        anterior:\includegraphics[width=4.821cm,height=0.944cm]{./ObjectReplacements/Object 965}.
        Si la lon\-gitud de la palabra es par podéis comprobar que es lo
        mismo poner que no poner todas las negaciones.
      \item
        Si la distancia mínima de un código es 1, entonces, en general
        no podremos lle\-gar a detectar fallos. Si un bit de una palabra
        nos llega equivocado tenemos que la palabra resultante puede
        estar perfectamente en el código, y así pasará inadvertido el
        error.
      \item
        La distancia mínima requerida para conseguir detectar un error
        de un solo bit es 2. Así un error de un bit, cuando menos nos
        dejará la palabra errónea a distancia 1 de cualquier palabra
        válida del código. Luego la palabra con el error no estará en el
        código. En general para poder detectar errores de hasta
        \includegraphics[width=1.155cm,height=0.467cm]{./ObjectReplacements/Object 966}necesitaremos
        que el código sea de distancia
        mínima\includegraphics[width=0.99cm,height=0.467cm]{./ObjectReplacements/Object 967}.
        El razonamiento es igual al expre\-sado para errores de 1 bit.
        Si la distancia mínima del código es
        \includegraphics[width=0.99cm,height=0.467cm]{./ObjectReplacements/Object 968}entonces,
        en el peor caso,
        \includegraphics[width=1.76cm,height=0.467cm]{./ObjectReplacements/Object 969}pueden
        dejarnos la palabra a distancia 1 (como míni\-mo) de cualquier
        palabra válida del código. Luego podemos detectar que la
        pala\-bra recibida no está en el código.
      \item
        Podemos dar un paso más e intentar no solo detectar sino
        corregir el bit equivo\-cado. Para esto necesitamos que la
        palabra con el bit equivocado quede de forma tal que sea claro
        desde cual palabra del código correcto se ha producido el error.
        No podremos tener dos palabras distintas del código a distancia
        mínima de la pa\-labra equivocada, sino solamente una. La
        corrección consistiría en cambiar la pa\-labra equivocada por la
        única más cercana del código. El principio de corrección es el
        máxima verosimilitud.
      \item
        Para poder
        corregir\includegraphics[width=1.76cm,height=0.467cm]{./ObjectReplacements/Object 970}necesitamos
        una distancia
        mínima\includegraphics[width=1.351cm,height=0.467cm]{./ObjectReplacements/Object 971}.
        Para corregir un solo bit necesitamos una distancia mínima de
        tres. Esto es claro. Si la distancia fuera dos, al producirse el
        error, quedaría a distancia uno de más de una palabra del
        código. Sin embargo si la distancia mínima es tres, al
        producirse un error quedará a distancia 2 como mínimo de todas
        las demás palabras del códi\-go. En general desde una
        palabra\includegraphics[width=0.665cm,height=0.635cm]{./ObjectReplacements/Object 1020}se
        ha producido con error la
        palabra\includegraphics[width=0.534cm,height=0.492cm]{./ObjectReplacements/Object 1021}tal
        que\includegraphics[width=2.559cm,height=0.854cm]{./ObjectReplacements/Object 1022}.
        Podemos decir que la palabra errónea está dentro del radio de la
        esfera que rodea
        a\includegraphics[width=0.665cm,height=0.635cm]{./ObjectReplacements/Object 1023}.
        La pregunta que surge es ¿puede existir una palabra
        \includegraphics[width=0.665cm,height=0.635cm]{./ObjectReplacements/Object 1024}con
        \[\mathtt{\mathrm{i} \neq \mathrm{j}}\]tal que
        \includegraphics[width=2.559cm,height=0.854cm]{./ObjectReplacements/Object 1027}?.
        Sabemos
        que\includegraphics[width=5.803cm,height=0.854cm]{./ObjectReplacements/Object 1028}.
        Esto quiere decir, que:

        \includegraphics[width=9.851cm,height=0.854cm]{./ObjectReplacements/Object 1029}

        \includegraphics[width=13.591cm,height=1.037cm]{./ObjectReplacements/Object 1030}

        \includegraphics[width=10.557cm,height=0.938cm]{./ObjectReplacements/Object 1031}

        \includegraphics[width=6.433cm,height=0.938cm]{./ObjectReplacements/Object 1032}

        \includegraphics[width=5.528cm,height=0.938cm]{./ObjectReplacements/Object 1033}

        \includegraphics[width=8.163cm,height=0.938cm]{./ObjectReplacements/Object 1034}

        \includegraphics[width=6.428cm,height=0.938cm]{./ObjectReplacements/Object 1035}

        Luego la corrección será
        \includegraphics[width=1.349cm,height=0.635cm]{./ObjectReplacements/Object 1036}necesariamente.
        Siempre habrá una sola pala\-bra de nuestro código que esté a
        distancia menor o igual
        que\includegraphics[width=0.466cm,height=0.467cm]{./ObjectReplacements/Object 1037},
        estando las de\-más palabras del código a distancia mayor o
        igual
        que\includegraphics[width=0.967cm,height=0.467cm]{./ObjectReplacements/Object 1038}.
      \item
        El método de construcción de detectores/correctores de error de
        Hamming es en principio el más fácil de usar, es el método
        lineal de codificación por grupo. Se trata de utilizar una
        matriz (de 0s y 1s) para convertir
        el\includegraphics[width=1.97cm,height=0.492cm]{./ObjectReplacements/Object 972}original,
        en
        un\includegraphics[width=2.769cm,height=0.547cm]{./ObjectReplacements/Object 973}.
        La condición principal es que la transformación sea inyecti\-va
        y\includegraphics[width=5.075cm,height=0.746cm]{./ObjectReplacements/Object 988}.
        Para esto lo fundamental es que la matriz de codificación (de
        dimensión\includegraphics[width=1.859cm,height=0.547cm]{./ObjectReplacements/Object 974})
        mantenga el original (con
        \includegraphics[width=0.43cm,height=0.467cm]{./ObjectReplacements/Object 975}colum\-nas
        dónde el único 1 está en la posición
        \includegraphics[width=0.944cm,height=0.519cm]{./ObjectReplacements/Object 977}con\includegraphics[width=1.625cm,height=0.482cm]{./ObjectReplacements/Object 976}),
        esto es, siendo un bloque matricial completo la matriz
        identidad\includegraphics[width=0.97cm,height=0.563cm]{./ObjectReplacements/Object 989},
        uno de los dos bloques de construcción de la
        matriz\includegraphics[width=0.453cm,height=0.492cm]{./ObjectReplacements/Object 990}de
        codificación Hamming. El otro bloque lo de\-notaremos
        por\includegraphics[width=0.974cm,height=0.563cm]{./ObjectReplacements/Object 991}.
        Esto además lo conseguiremos si la matriz de decodifica\-ción
        \includegraphics[width=0.674cm,height=0.603cm]{./ObjectReplacements/Object 992}de
        dimensión
        \includegraphics[width=1.85cm,height=0.531cm]{./ObjectReplacements/Object 993}es
        tal que
        \includegraphics[width=3.865cm,height=0.908cm]{./ObjectReplacements/Object 994}.
        Si el resultado de la primera decodificación no es cero es que
        hay error. Entonces, si llamamos
        \includegraphics[width=0.534cm,height=0.48cm]{./ObjectReplacements/Object 995}a
        la palabra recibida y
        \includegraphics[width=0.466cm,height=0.48cm]{./ObjectReplacements/Object 996}a
        la palabra enviada (siempre correcta), ha de existir un vector
        columna\includegraphics[width=1.73cm,height=0.515cm]{./ObjectReplacements/Object 997},
        con un único uno
        en\includegraphics[width=0.349cm,height=0.235cm]{./ObjectReplacements/Object 998},
        tal
        que\includegraphics[width=4.487cm,height=0.612cm]{./ObjectReplacements/Object 999}.
        Por otra parte, condición necesaria y sufi\-ciente para que la
        operación sea inyectiva es que no contenga columnas igua\-les ni
        columnas vector 0.

        Ejemplo:

        Supongamos que nos llega un código de 2 bits cuyos bits
        llamaremos\includegraphics[width=0.954cm,height=0.531cm]{./ObjectReplacements/Object 978},
        a su vez los bits codificados por la matriz generadora
        serán\includegraphics[width=1.593cm,height=0.531cm]{./ObjectReplacements/Object 979}.
        En general para poder corregir un bit necesitamos distancia 3:
        ¿cuántos bits habrá que añadir?. En general si queremos corregir
        códigos, hemos de añadir un número
        \includegraphics[width=0.33cm,height=0.467cm]{./ObjectReplacements/Object 1039}de
        bits a
        los\includegraphics[width=0.422cm,height=0.467cm]{./ObjectReplacements/Object 1040}originales
        tal que la distancia sea mayor o igual
        que\includegraphics[width=0.418cm,height=0.467cm]{./ObjectReplacements/Object 980},
        y probando, tenemos
        que\includegraphics[width=1.094cm,height=0.467cm]{./ObjectReplacements/Object 983}ya
        cumple. Así que la matriz de codificación de grupo será de
        dimensión\includegraphics[width=1.012cm,height=0.467cm]{./ObjectReplacements/Object 984}.

        \includegraphics[width=3.33cm,height=0.586cm]{./ObjectReplacements/Object 1010}
      \end{enumerate}
    \end{enumerate}
  \end{enumerate}
\end{enumerate}

\[\begin{array}{l}
{= {\begin{pmatrix}
1 & 0 \\
0 & 1 \\
a & b \\
c & d
\end{pmatrix} \cdot \begin{pmatrix}
1 & 0 & 1 \\
0 & 1 & 1
\end{pmatrix}} =} \\
{= \begin{pmatrix}
1 & 0 & 1 \\
0 & 1 & 1 \\
a & b & {a\oplus b} \\
c & d & {c\oplus d}
\end{pmatrix}}
\end{array}\]

\begin{enumerate}
\def\labelenumi{\arabic{enumi}.}
\setcounter{enumi}{44}
\item
  \begin{enumerate}
  \def\labelenumii{\arabic{enumii}.}
  \item
    \begin{enumerate}
    \def\labelenumiii{\arabic{enumiii}.}
    \item
      \begin{enumerate}
      \def\labelenumiv{\arabic{enumiv}.}
      \item
        \begin{enumerate}
        \def\labelenumv{\arabic{enumv}.}
        \tightlist
        \item
          \begin{enumerate}
          \def\labelenumvi{\arabic{enumvi}.}
          \tightlist
          \item
            \begin{enumerate}
            \def\labelenumvii{\arabic{enumvii}.}
            \tightlist
            \item
              \begin{enumerate}
              \def\labelenumviii{\arabic{enumviii}.}
              \tightlist
              \item
                \begin{enumerate}
                \def\labelenumix{\arabic{enumix}.}
                \tightlist
                \item
                  \begin{enumerate}
                  \def\labelenumx{\arabic{enumx}.}
                  \tightlist
                  \item
                  \end{enumerate}
                \end{enumerate}
              \end{enumerate}
            \end{enumerate}
          \end{enumerate}
        \end{enumerate}

        La matriz izquierda de la primera fila es la matriz generadora
        de paridades, la de\-recha son el código original de 4 bits en
        forma de una palabra por cada columna. Por último la matriz
        inferior son los vectores codificados con tres bits de paridad
        insertos. Solo queda ver la transformación inversa
        correspondiente. Para esto vea\-mos que si observamos matriz
        izquierda de la primera
        fila\includegraphics[width=1.974cm,height=1.265cm]{./ObjectReplacements/Object 986}.
        Tomare\-mos como matriz decodificadora (no necesariamente
        inyectiva) a

        \begin{enumerate}
        \def\labelenumv{\arabic{enumv}.}
        \item
          \begin{enumerate}
          \def\labelenumvi{\arabic{enumvi}.}
          \item
            \includegraphics[width=3.223cm,height=0.67cm]{./ObjectReplacements/Object 987}

            \includegraphics[width=3.514cm,height=1.074cm]{./ObjectReplacements/Object 1000}
          \end{enumerate}
        \end{enumerate}

        Ahora bien, sabemos que:

        \includegraphics[width=2.094cm,height=0.593cm]{./ObjectReplacements/Object 1002}

        De dónde obtenemos un sistema de ecuaciones, con estos
        resultados (distintos del resultado trivial):

        \includegraphics[width=2.872cm,height=1.074cm]{./ObjectReplacements/Object 1004}

        Que hace
        que\[\mathtt{\mathrm{H}} \cdot \mathtt{\mathrm{o}}\]quede como:

        \[\begin{pmatrix}
        1 & 0 & 1 \\
        0 & 1 & 1
        \end{pmatrix}\rightarrow\begin{pmatrix}
        1 & 0 & 1 \\
        0 & 1 & 1 \\
        a & b & {a\oplus b} \\
        c & d & {c\oplus d}
        \end{pmatrix}\]

        El sistema de ecuaciones se genera como se ve a continuación:

        \includegraphics[width=5.547cm,height=2.18cm]{./ObjectReplacements/Object 1018}

        \includegraphics[width=8.916cm,height=1.159cm]{./ObjectReplacements/Object 982}

        \includegraphics[width=8.551cm,height=1.155cm]{./ObjectReplacements/Object 985}

        \[\begin{bmatrix}
        {c = \overline{a}} & {{a \cdot b} = {\overline{a} \cdot d}} & {1\oplus{a \cdot \overline{b}}\oplus{c \cdot \overline{d}}} \\
        {{a \cdot b} = {\overline{a} \cdot d}} & {d = \overline{b}} & {1\oplus{b \cdot \overline{a}}\oplus{d \cdot \overline{c}}}
        \end{bmatrix}\]

        \includegraphics[width=6.523cm,height=1.067cm]{./ObjectReplacements/Object 1017}

        \includegraphics[width=5.92cm,height=1.124cm]{./ObjectReplacements/Object 1019}

        \includegraphics[width=6.07cm,height=1.067cm]{./ObjectReplacements/Object 1041}

        \includegraphics[width=7.664cm,height=1.556cm]{./ObjectReplacements/Object 1042}

        \includegraphics[width=2.805cm,height=1.074cm]{./ObjectReplacements/Object 1043}

        \includegraphics[width=3.45cm,height=1.074cm]{./ObjectReplacements/Object 1044}

        \includegraphics[width=2.027cm,height=2.18cm]{./ObjectReplacements/Object 1045}

        \includegraphics[width=5.419cm,height=2.18cm]{./ObjectReplacements/Object 981}
      \end{enumerate}
    \end{enumerate}
  \end{enumerate}
\item
  Funciones de Boole
  de\includegraphics[width=2.367cm,height=0.467cm]{./ObjectReplacements/Object 910}en
  \includegraphics[width=2.198cm,height=0.467cm]{./ObjectReplacements/Object 912}:
  tablas lineales, 2-dimensiona\-les, arreglos de tablas 2-dimensionales
  y otras formas de representación. Simplificación en 2 capas de
  puertas.
\end{enumerate}

\end{document}
